% Circulation sanguine — SVTEEHB 3ème — Cours signé OnBuch+
% Programme officiel MINESEC. Compiler avec : tectonic N09-circulation-sanguine.tex
\newcommand{\DOCMATIERE}{SVTEEHB}
\newcommand{\DOCNIVEAU}{3ème}
\newcommand{\DOCMODULE}{SVTEEHB – classe de 3ème}
\newcommand{\DOCLECON}{N09}
\newcommand{\DOCTITRE}{Circulation sanguine}
% OnBuch+ — préambule « cours signés » ENRICHI (compilé avec tectonic / XeLaTeX)
% Système visuel « Pop » d'OnBuch+ : fond crème (jamais blanc), bordures encre
% épaisses, ombres « dures » décalées, titres Archivo Black en capitales.
% Version MATHÉMATIQUES : graphes (pgfplots/tikz), boîtes pédagogiques
% (définition, propriété, méthode, attention, le savais-tu, exercice résolu,
% exercice, TP, bilan), page de garde et signature OnBuch+ ancrée en bas.
%
% Macros attendues dans main.tex AVANT \input{preamble} :
%   \DOCMATIERE  \DOCNIVEAU  \DOCMODULE  \DOCLECON (numéro)  \DOCTITRE
\documentclass[11pt]{article}

\usepackage{fontspec}
\usepackage[french]{babel}
\usepackage[a4paper,margin=2cm,top=3.4cm,bottom=2.8cm,headheight=1.4cm,footskip=1.3cm]{geometry}
\usepackage{amsmath,amssymb,amsfonts}
\usepackage{mathtools} % \xmapsto, \coloneqq... souvent utilisés par les modèles
\usepackage{cancel} % simplifications barrées (bilans de forces)
% \abs{z} (module, valeur absolue) et \norm{u} (norme), utilisés par les modèles
\providecommand{\abs}{}\let\abs\relax\DeclarePairedDelimiter{\abs}{\lvert}{\rvert}
\providecommand{\norm}{}\let\norm\relax\DeclarePairedDelimiter{\norm}{\lVert}{\rVert}
\usepackage[table]{xcolor} % [table] : fournit aussi \rowcolors (alternance de lignes)
\usepackage{pagecolor}
\usepackage{tikz}
\usetikzlibrary{arrows.meta,positioning,calc,shapes.geometric,shapes.misc,shapes.arrows,decorations.pathmorphing,decorations.pathreplacing,decorations.markings,patterns,shadows,mindmap,trees,fit,backgrounds,matrix,intersections,angles,quotes}
\usepackage{pgfplots}
\usepackage[version=4]{mhchem} % équations nucléaires \ce{^{235}_{92}U}
\usepackage[european,siunitx]{circuitikz} % schémas électriques
\pgfplotsset{compat=1.18}
\usepgfplotslibrary{fillbetween,groupplots} % aires entre courbes (calcul intégral)
\usepackage{enumitem}
\usepackage{fancyhdr}
% [most] charge toutes les bibliothèques tcolorbox (listings, minted, raster,
% documentation...) et gonfle inutilement la mémoire TeX sur un document long
% (~300 boîtes) : on ne charge que ce qui est réellement utilisé.
\usepackage[skins,breakable]{tcolorbox}
\usepackage{graphicx}
\usepackage{colortbl}
\usepackage{booktabs}
\usepackage{tabularx}
\usepackage{multirow}
\usepackage{diagbox} % case d'en-tête diagonale des tableaux à double entrée
% Colonnes extensibles centrée / gauche / droite (C, L, R), souvent utilisées par les modèles
\newcolumntype{C}{>{\centering\arraybackslash}X}
\newcolumntype{L}{>{\raggedright\arraybackslash}X}
\newcolumntype{R}{>{\raggedleft\arraybackslash}X}
\usepackage{array}
\usepackage{multicol}
\usepackage{siunitx}
% parse-numbers=false : accepte aussi \SI{2,5 \times 10^{-3}}{...}
\sisetup{locale=FR,output-decimal-marker={,},group-minimum-digits=4,parse-numbers=false}
\DeclareSIUnit\ohm{\ensuremath{\symup{\Omega}}} % U+2126 absent de la police
\DeclareSIUnit\division{div} % réglages d'oscilloscope (ms/div, V/div)
\DeclareSIUnit\tr{tr} % tours (vitesse de rotation, ex. \SI{1500}{\tr\per\minute})
\DeclareSIUnit\molar{M} % concentration molaire (ex. \SI{3}{\molar}, \SI{1}{\micro\molar})
\DeclareSIUnit\an{an} % année (le modèle confond parfois avec \year, un compteur TeX) — ex. \SI{5}{\kilo\meter\per\an}
\DeclareSIUnit\FCFA{FCFA} % franc CFA (ex. \SI{500}{\FCFA\per\kilo\gram})
\DeclareSIUnit\degreeBrix{\textdegree Brix} % teneur en sucre d'un sirop/jus (réfractométrie)
\DeclareSIUnit\month{mois} % le modèle utilise parfois \month, un compteur TeX, comme unité de durée
\DeclareSIUnit\microliter{\micro\liter} % le modèle écrit parfois \microliter en un seul token
\DeclareSIUnit\million{millions} % dénombrement (ex. \SI{15}{\million\per\mL} spermatozoïdes)
\usepackage{qrcode}
% \degree : commande fréquemment utilisée par les modèles pour un angle ou une
% température en dehors de \SI{}{} (non fournie par les paquets chargés).
\providecommand{\degree}{^{\circ}}
% \qed (fin de démonstration), utilisé par les modèles sans amsthm
\providecommand{\qed}{\hfill$\square$}
% \barre : double barre des valeurs interdites dans un tableau de variations (tkz-tab)
\providecommand{\barre}{\ensuremath{\|}}
% environnement proof (amsthm non chargé) utilisé par les modèles
\ifdefined\proof\else\newenvironment{proof}[1][Démonstration]{\par\noindent\textit{#1.}\ }{\qed\par}\fi
\usepackage{lastpage}
\usepackage{hyperref}
\hypersetup{hidelinks}

% ============================================================
% Palette « Pop » OnBuch+ — mêmes hex que la classe P (plus_ui.dart)
% ============================================================
\definecolor{poporange}{HTML}{F59321}
\definecolor{popdark}{HTML}{E07A0C}
\definecolor{popcream}{HTML}{FFF6E8}
\definecolor{popink}{HTML}{1C1714}
\definecolor{poppurple}{HTML}{7A5AE0}
\definecolor{popgreen}{HTML}{2E9E6A}
\definecolor{poppink}{HTML}{E0457B}
\definecolor{popblue}{HTML}{2F7DE1}
\definecolor{popgold}{HTML}{C98A12}
\definecolor{popmuted}{HTML}{8A7F76}
\definecolor{popline}{HTML}{EADFD2}
\definecolor{popgray}{HTML}{8A7F76}
\colorlet{popgrey}{popgray}
% Alias tolérants (le modèle nomme parfois les couleurs différemment)
\colorlet{popwhite}{white}
\colorlet{popblack}{popink}
\colorlet{popred}{poppink}
\colorlet{popyellow}{popgold}
% Teintes claires (fonds de tableaux/encadrés) — le modèle nomme parfois
% les couleurs avec un suffixe L (ex. popblueL) sans les avoir définies.
\colorlet{poporangeL}{poporange!20}
\colorlet{popdarkL}{popdark!20}
\colorlet{poppurpleL}{poppurple!20}
\colorlet{popgreenL}{popgreen!20}
\colorlet{poppinkL}{poppink!20}
\colorlet{popblueL}{popblue!20}
\colorlet{popgoldL}{popgold!20}
\colorlet{popredL}{popred!20}
\colorlet{popgrayL}{popgray!20}
% Teintes claires pour fonds de tableaux / zones
\colorlet{popblueL}{popblue!10!white}
\colorlet{poporangeL}{poporange!14!white}
\colorlet{popgreenL}{popgreen!12!white}
\colorlet{poppurpleL}{poppurple!12!white}
\colorlet{poppinkL}{poppink!12!white}
\colorlet{popgoldL}{popgold!14!white}

\pagecolor{popcream}

% ============================================================
% Typographie
% ============================================================
\defaultfontfeatures{Ligatures=TeX}
\setmainfont{PlusJakartaSans-Medium.ttf}[Path=../../../fonts/,BoldFont=PlusJakartaSans-Bold.ttf]
% Maths en sans-serif (Fira Math) avec chiffres et lettres droites de Plus
% Jakarta, pour que les formules se fondent dans le texte.
\usepackage[math-style=ISO,bold-style=ISO]{unicode-math}
\setmathfont{FiraMath-Regular.otf}[Path=../../../fonts/]
\setmathfont{PlusJakartaSans-Medium.ttf}[Path=../../../fonts/,range={up/num,up/{Latin,latin},\mathup}]
\usepackage{icomma} % virgule décimale sans espace en mode math
% Caractères mathématiques Unicode absents des polices de texte (Plus Jakarta,
% Archivo Black) : rendus par leur équivalent mathématique, en texte comme en
% formule — sinon ils s'affichent en carré vide (ex. « Arithmétique dans ℤ »).
\usepackage{newunicodechar}
\newunicodechar{ℝ}{\ensuremath{\mathbb{R}}}
\newunicodechar{ℤ}{\ensuremath{\mathbb{Z}}}
\newunicodechar{ℂ}{\ensuremath{\mathbb{C}}}
\newunicodechar{ℕ}{\ensuremath{\mathbb{N}}}
\newunicodechar{ℚ}{\ensuremath{\mathbb{Q}}}
\newunicodechar{𝔻}{\ensuremath{\mathbb{D}}}
\newunicodechar{α}{\ensuremath{\alpha}}
\newunicodechar{β}{\ensuremath{\beta}}
\newunicodechar{γ}{\ensuremath{\gamma}}
\newunicodechar{δ}{\ensuremath{\delta}}
\newunicodechar{ε}{\ensuremath{\varepsilon}}
\newunicodechar{θ}{\ensuremath{\theta}}
\newunicodechar{λ}{\ensuremath{\lambda}}
\newunicodechar{μ}{\ensuremath{\mu}}
\newunicodechar{σ}{\ensuremath{\sigma}}
\newunicodechar{τ}{\ensuremath{\tau}}
\newunicodechar{φ}{\ensuremath{\varphi}}
\newunicodechar{ψ}{\ensuremath{\psi}}
\newunicodechar{ω}{\ensuremath{\omega}}
\newunicodechar{Σ}{\ensuremath{\Sigma}}
% majuscules produites par \MakeUppercase dans les titres (α-aminés -> Α)
\newunicodechar{Α}{\ensuremath{\alpha}}
\newunicodechar{Β}{\ensuremath{\beta}}
\newunicodechar{Γ}{\ensuremath{\Gamma}}
\newunicodechar{Δ}{\ensuremath{\Delta}}
\newunicodechar{Ε}{\ensuremath{\varepsilon}}
\newunicodechar{Θ}{\ensuremath{\Theta}}
\newunicodechar{Λ}{\ensuremath{\Lambda}}
\newunicodechar{Μ}{\ensuremath{\mu}}
\newunicodechar{Τ}{\ensuremath{\tau}}
\newunicodechar{Φ}{\ensuremath{\Phi}}
\newunicodechar{Ψ}{\ensuremath{\Psi}}
\newunicodechar{Ω}{\ensuremath{\Omega}}
\newunicodechar{↦}{\ensuremath{\mapsto}}
\newunicodechar{∈}{\ensuremath{\in}}
\newunicodechar{∉}{\ensuremath{\notin}}
\newunicodechar{∧}{\ensuremath{\wedge}}
\newunicodechar{⇔}{\ensuremath{\Leftrightarrow}}
\newunicodechar{⟺}{\ensuremath{\Longleftrightarrow}}
\newunicodechar{⇒}{\ensuremath{\Rightarrow}}
\newunicodechar{⟹}{\ensuremath{\Longrightarrow}}
\newunicodechar{∘}{\ensuremath{\circ}}
\newunicodechar{∪}{\ensuremath{\cup}}
\newunicodechar{∩}{\ensuremath{\cap}}
\newunicodechar{≡}{\ensuremath{\equiv}}
\newunicodechar{⊂}{\ensuremath{\subset}}
\newunicodechar{⊥}{\ensuremath{\perp}}
\newunicodechar{∃}{\ensuremath{\exists}}
\newunicodechar{∀}{\ensuremath{\forall}}
\newunicodechar{∛}{\ensuremath{\sqrt[3]{\,}}}
\newunicodechar{∜}{\ensuremath{\sqrt[4]{\,}}}
\newunicodechar{⃗}{\ensuremath{{}^{\to}}}
\newunicodechar{ⁿ}{\ifmmode^{n}\else\textsuperscript{\ensuremath{n}}\fi}
\newunicodechar{ˣ}{\ifmmode^{x}\else\textsuperscript{\ensuremath{x}}\fi}
\newunicodechar{ᵇ}{\ifmmode^{b}\else\textsuperscript{\ensuremath{b}}\fi}
\newunicodechar{ʳ}{\ifmmode^{r}\else\textsuperscript{\ensuremath{r}}\fi}
\newunicodechar{ᵘ}{\ifmmode^{u}\else\textsuperscript{\ensuremath{u}}\fi}
\newunicodechar{ʸ}{\ifmmode^{y}\else\textsuperscript{\ensuremath{y}}\fi}
\newunicodechar{ᵃ}{\ifmmode^{a}\else\textsuperscript{\ensuremath{a}}\fi}
\newunicodechar{ᵉ}{\ifmmode^{e}\else\textsuperscript{\ensuremath{e}}\fi}
\newunicodechar{ᵏ}{\ifmmode^{k}\else\textsuperscript{\ensuremath{k}}\fi}
\newunicodechar{ᵖ}{\ifmmode^{p}\else\textsuperscript{\ensuremath{p}}\fi}
\newunicodechar{ᴺ}{\ifmmode^{N}\else\textsuperscript{\ensuremath{N}}\fi}
\newunicodechar{ᶜ}{\ifmmode^{c}\else\textsuperscript{\ensuremath{c}}\fi}
\newunicodechar{ʰ}{\ifmmode^{h}\else\textsuperscript{\ensuremath{h}}\fi}
\newunicodechar{ᵠ}{\ifmmode^{q}\else\textsuperscript{\ensuremath{q}}\fi}
\newunicodechar{ₙ}{\ifmmode_{n}\else\textsubscript{\ensuremath{n}}\fi}
\newunicodechar{ₐ}{\ifmmode_{a}\else\textsubscript{\ensuremath{a}}\fi}
\newunicodechar{ᵢ}{\ifmmode_{i}\else\textsubscript{\ensuremath{i}}\fi}
\newunicodechar{ᵣ}{\ifmmode_{r}\else\textsubscript{\ensuremath{r}}\fi}
\newunicodechar{ₖ}{\ifmmode_{k}\else\textsubscript{\ensuremath{k}}\fi}
\newunicodechar{ᵦ}{\ifmmode_{\beta}\else\textsubscript{\ensuremath{\beta}}\fi}
\newunicodechar{ₓ}{\ifmmode_{x}\else\textsubscript{\ensuremath{x}}\fi}
\newfontfamily\popdisplay{ArchivoBlack-Regular.ttf}[Path=../../../fonts/]
\newfontfamily\poplogo{SpaceGrotesk-Bold.ttf}[Path=../../../fonts/]
\newfontfamily\popsemi{PlusJakartaSans-SemiBold.ttf}[Path=../../../fonts/]
\newfontfamily\popxbold{PlusJakartaSans-ExtraBold.ttf}[Path=../../../fonts/]
\newfontfamily\popxboldls{PlusJakartaSans-ExtraBold.ttf}[Path=../../../fonts/,LetterSpace=3.0]

\setlist[enumerate]{leftmargin=1.7em,itemsep=5pt,topsep=4pt}
\setlist[itemize]{leftmargin=1.7em,itemsep=4pt,topsep=4pt,label={\color{poporange}\textbullet}}
\setlength{\parindent}{0pt}
\setlength{\parskip}{5pt}
\renewcommand{\arraystretch}{1.25}

% pgfplots : style Pop par défaut
\pgfplotsset{
  every axis/.append style={axis line style={popink,line width=1pt},tick style={popink},
    grid style={popline,line width=0.5pt},label style={font=\small},tick label style={font=\footnotesize}},
  popcurve/.style={poporange,line width=1.8pt,no markers,smooth},
}
% Même style disponible dans un \draw TikZ simple (hors pgfplots)
\tikzset{popcurve/.style={poporange,line width=1.8pt}}
\tikzset{popgrid/.style={popline,very thin}}
\colorlet{popcurve}{poporange} % le nom du style est parfois utilisé comme couleur

% ============================================================
% Pill / squiggle
% ============================================================
\newcommand{\poppill}[3]{%
  \tikz[baseline=-0.55ex]\node[draw=popink,line width=1.2pt,rounded corners=2.6pt,%
    fill=#2,inner xsep=6.5pt,inner ysep=3pt]{%
    \popxboldls\fontsize{7}{7}\selectfont\color{#3}\MakeUppercase{#1}};%
}
\newcommand{\popsquiggle}[1]{%
  \tikz[baseline=-0.4ex]{
    \draw[#1,line width=2.6pt,line cap=round]
      plot [smooth] coordinates {(0,0.09) (0.34,0.01) (0.68,0.21) (1.02,0.01) (1.36,0.10)};
  }%
}
% Mot-clé surligné (marqueur orange sous le texte)
\newcommand{\cle}[1]{{\popxbold\color{popdark}#1}}

% ============================================================
% En-tête / pied
% ============================================================
\pagestyle{fancy}
\fancyhf{}
\renewcommand{\headrulewidth}{0pt}
\renewcommand{\footrulewidth}{0pt}
\lhead{\raisebox{-0.15ex}{\poplogo\fontsize{13}{13}\selectfont\color{popink}OnBuch\color{poporange}+}}
\rhead{\poppill{\DOCMATIERE\ \textbullet\ \DOCNIVEAU\ \textbullet\ Leçon \DOCLECON}{white}{popink}}
\lfoot{\popxbold\fontsize{7.6}{7.6}\selectfont\color{popmuted}\MakeUppercase{Cours signé OnBuch+}\ \textbullet\ {\popsemi\DOCTITRE}}
\rfoot{\tikz[baseline=-0.6ex]\node[rounded corners=8.5pt,fill=popink,minimum height=17pt,minimum width=17pt,inner xsep=5pt,inner sep=0]{\popxbold\fontsize{8}{8}\selectfont\color{popcream}\thepage\,/\,\pageref*{LastPage}};}

% ============================================================
% Titres
% ============================================================
\newcommand{\chaptitre}[2]{%
  \begin{tcolorbox}[enhanced,colback=poporange,colframe=popink,boxrule=2.6pt,arc=11pt,
      drop shadow=popink,left=15pt,right=15pt,top=12pt,bottom=11pt,before skip=3pt,after skip=18pt]
    \poppill{#2}{popink}{popcream}\par\vspace{9pt}
    {\raggedright\hyphenpenalty=10000\exhyphenpenalty=10000\popdisplay\fontsize{22}{24}\selectfont\color{popcream}\MakeUppercase{#1}\par}
  \end{tcolorbox}
}
% Section numérotée : pastille orange avec le numéro + titre Archivo
\newcounter{popsec}
\newcommand{\coursec}[1]{%
  \stepcounter{popsec}\setcounter{popsub}{0}%
  \par\vspace{16pt}\noindent
  \begin{minipage}{\linewidth}
  \tikz[baseline=-0.7ex]\node[draw=popink,line width=1.6pt,fill=poporange,rounded corners=4pt,minimum size=22pt,inner sep=0]{\popdisplay\fontsize{12}{12}\selectfont\color{popcream}\thepopsec};%
  \hspace{8pt}{\popdisplay\fontsize{15}{17}\selectfont\color{popink}\MakeUppercase{#1}}\par
  \vspace{4pt}\noindent{\color{poporange}\rule{36pt}{3.6pt}}
  \end{minipage}\par\nopagebreak\vspace{8pt}
}
\newcounter{popsub}
\newcommand{\courssub}[1]{%
  \stepcounter{popsub}%
  \par\vspace{9pt}\noindent{\popxbold\fontsize{11.5}{13}\selectfont\color{popink}{\color{poporange}\thepopsec.\thepopsub}\hspace{6pt}#1}\par\nopagebreak\vspace{4pt}
}

% ============================================================
% Boîtes pédagogiques — même recette : fond blanc, bordure épaisse colorée,
% ombre dure encre, badge plein. Titre optionnel [#1].
% ============================================================
\tcbset{popbase/.style={breakable,enhanced,colback=white,boxrule=2.3pt,arc=9pt,
  shadow={3pt}{-3pt}{0pt}{popink},left=14pt,right=14pt,top=11pt,bottom=11pt,before skip=11pt,after skip=15pt,
  fontupper=\popsemi\fontsize{10.6}{14.5}\selectfont\color{popink},
  fontlower=\popsemi\fontsize{10.6}{14.5}\selectfont\color{popink}}}
% Coche dessinée (le glyphe ✓ n'existe pas dans les polices du document)
\newcommand{\popcheck}[1]{\tikz[baseline=-0.5ex]\draw[#1,line width=1.6pt,line cap=round,line join=round](-0.13,0.0)--(-0.04,-0.09)--(0.13,0.1);}
\providecommand{\coche}{\popcheck{popgreen}}
\providecommand{\resultat}[1]{\par\smallskip\noindent\fcolorbox{popgreen}{popgreenL}{\parbox{\dimexpr\linewidth-2\fboxsep-2\fboxrule\relax}{\textbf{Résultat.} #1}}\par\smallskip}
\newcommand{\popboxhead}[3]{\poppill{#1}{#2}{white}\ifx\relax#3\relax\else\hspace{6pt}{\popxbold\fontsize{10.6}{12}\selectfont\color{popink}#3}\fi\par\vspace{7pt}}

\newtcolorbox{aretenir}[1][]{popbase,colframe=popgreen,before upper={\popboxhead{À retenir}{popgreen}{#1}}}
\newtcolorbox{exemplebox}[1][]{popbase,colframe=poporange,before upper={\popboxhead{Exemple}{poporange}{#1}}}
\newtcolorbox{definition}[1][]{popbase,colframe=popblue,before upper={\popboxhead{Définition}{popblue}{#1}}}
\newtcolorbox{propriete}[1][]{popbase,colframe=poppurple,before upper={\popboxhead{Propriété}{poppurple}{#1}}}
\newtcolorbox{methode}[1][]{popbase,colframe=popgold,before upper={\popboxhead{Méthode}{popgold}{#1}}}
\newtcolorbox{attention}[1][]{popbase,colframe=poppink,before upper={\popboxhead{Attention}{poppink}{#1}}}
\newtcolorbox{experience}[1][]{popbase,colframe=popblue,colback=popblueL,before upper={\popboxhead{Expérience}{popblue}{#1}}}
% « Le savais-tu ? » : carte violette pleine (contexte, culture, Cameroun)
\newtcolorbox{savaistu}[1][]{popbase,colback=poppurple,colframe=popink,
  fontupper=\popsemi\fontsize{10.4}{14.2}\selectfont\color{white},
  before upper={\poppill{Le savais-tu\,?}{white}{poppurple}\ifx\relax#1\relax\else\hspace{6pt}{\popxbold\color{white}#1}\fi\par\vspace{7pt}}}
% Exercice résolu : énoncé en haut, \tcblower puis la solution sur fond crème
\newcounter{popexo}\newcounter{popexr}
\newtcolorbox{exoresolu}[1][]{popbase,colframe=poporange,colbacklower=popcream,
  segmentation style={popink,line width=1.2pt,dashed},
  before upper={\stepcounter{popexr}\popboxhead{Exercice résolu \thepopexr}{poporange}{#1}},
  before lower={\poppill{Solution}{popink}{popcream}\par\vspace{6pt}}}
% Exercice (non résolu en place) — la correction est dans la partie Corrigés
\newtcolorbox{exercice}[1][]{popbase,colframe=popink,
  before upper={\stepcounter{popexo}\popboxhead{Exercice \thepopexo}{popink}{#1}}}
% Corrigé
\newtcolorbox{corrige}[1][]{popbase,colframe=popgreen,colback=popgreenL,
  before upper={\popboxhead{Corrigé}{popgreen}{#1}}}
% Objectifs / prérequis / bilan
\newtcolorbox{objectifs}{popbase,colframe=popink,colback=poporangeL,
  before upper={\popboxhead{Objectifs de la leçon}{popink}{}}}
\newtcolorbox{prerequis}{popbase,colframe=popink,colback=popblueL,
  before upper={\popboxhead{Prérequis}{popblue}{}}}
\newtcolorbox{bilan}{popbase,colframe=popink,colback=white,boxrule=2.8pt,
  before upper={\popboxhead{Fiche bilan}{poporange}{L'essentiel à connaître pour l'examen}}}
% Figure encadrée avec légende
\newtcolorbox{popfigure}[1][]{enhanced,breakable=false,colback=white,colframe=popline,boxrule=1.4pt,arc=8pt,
  left=10pt,right=10pt,top=10pt,bottom=8pt,before skip=10pt,after skip=12pt,halign=center,
  lower separated=false,halign lower=center,
  fontlower=\popsemi\fontsize{9}{11.5}\selectfont\color{popmuted},
  #1}
\newcommand{\legende}[1]{\par\vspace{6pt}{\popsemi\fontsize{9}{11.5}\selectfont\color{popmuted}#1}}

% ============================================================
% Signature OnBuch+
% ============================================================
\newcommand{\poplogoinline}[1]{{\poplogo\fontsize{#1}{#1}\selectfont\color{popink}OnBuch\color{poporange}+}}
\newcommand{\signatureonbuch}{%
  \begin{tcolorbox}[enhanced,colback=popink,colframe=popink,boxrule=2.2pt,arc=10pt,
      drop shadow=poporange,left=14pt,right=14pt,top=10pt,bottom=10pt,before skip=0pt,after skip=0pt]
    \tikz[baseline=-0.7ex]\node[circle,fill=popgreen,draw=popcream,line width=1.3pt,minimum size=22pt,inner sep=0]{\popcheck{white}};%
    \hspace{9pt}\begin{minipage}[c]{0.62\linewidth}
      {\popxbold\fontsize{12}{13}\selectfont\color{popcream}Signé\ \textbullet\ {\poplogo OnBuch\color{poporange}+}}\par\vspace{2pt}
      {\raggedright\popsemi\fontsize{8}{10.5}\selectfont\color{popline}\MakeUppercase{Équipe pédagogique OnBuch+}\\\MakeUppercase{Conforme au programme officiel MINESEC}\par}
    \end{minipage}\hfill
    {\poplogo\fontsize{22}{22}\selectfont\color{popcream}OnBuch\color{poporange}+}
  \end{tcolorbox}%
}

% ============================================================
% Page de garde
% #1 titre · #2 matière · #3 niveau · #4 module · #5 n° leçon · #6 durée ·
% #7 description / objectifs (texte ou itemize)
% ============================================================
\newcommand{\pagegarde}[7]{%
  \thispagestyle{empty}
  \newgeometry{a4paper,margin=1.7cm,top=1.6cm,bottom=1.4cm}%
  \begin{tikzpicture}[remember picture,overlay]
    \fill[poporange!18] ([xshift=-3.2cm,yshift=-3.6cm]current page.north east) circle (4.6cm);
    \fill[poppurple!14] ([xshift=2.4cm,yshift=7.6cm]current page.south west) circle (3.8cm);
  \end{tikzpicture}%
  {\poplogo\fontsize{30}{30}\selectfont\color{popink}OnBuch\color{poporange}+}\hfill
  \raisebox{0.6ex}{\poppill{Cours signé \textbullet\ Premium}{popink}{popcream}}\par
  \vspace{1pt}\popsquiggle{poppurple}\par
  \vspace{8pt}
  \begin{tcolorbox}[enhanced,colback=poporange,colframe=popink,boxrule=2.8pt,arc=13pt,
      drop shadow=popink,left=18pt,right=18pt,top=12pt,bottom=13pt,after skip=10pt]
    \poppill{#2 \textbullet\ #3}{popink}{popcream}\hspace{5pt}\poppill{#4}{white}{popink}\par\vspace{8pt}
    {\popdisplay\fontsize{14}{14}\selectfont\color{popink}LEÇON #5}\par\vspace{4pt}
    {\raggedright\hyphenpenalty=10000\exhyphenpenalty=10000\popdisplay\fontsize{25}{27}\selectfont\color{popcream}\MakeUppercase{#1}\par}
    \vspace{8pt}
    {\popxbold\fontsize{9.5}{9.5}\selectfont\color{popink}\MakeUppercase{Durée conseillée : #6}}
  \end{tcolorbox}
  \begin{tcolorbox}[enhanced,colback=white,colframe=popink,boxrule=2.2pt,arc=10pt,
      drop shadow=popink,left=16pt,right=16pt,top=10pt,bottom=10pt,after skip=8pt]
    \poppill{Dans cette leçon}{poppurple}{white}\par\vspace{5pt}
    \popsemi\fontsize{9.3}{12.6}\selectfont\color{popink}#7
  \end{tcolorbox}
  \noindent\begin{minipage}[t]{0.487\linewidth}
    \begin{tcolorbox}[enhanced,colback=popgreen,colframe=popink,boxrule=2.2pt,arc=10pt,
        drop shadow=popink,left=11pt,right=11pt,top=10pt,bottom=10pt,height=3.1cm,valign=center]
      \tcbox[colback=white,colframe=popink,boxrule=1.3pt,arc=5pt,boxsep=0pt,nobeforeafter,
        left=3pt,right=3pt,top=3pt,bottom=3pt,valign=center]{\qrcode[height=1.9cm]{https://play.google.com/store/apps/details?id=com.luvvix.onbuch&hl=fr}}%
      \hspace{8pt}\begin{minipage}[c]{0.5\linewidth}
        {\popdisplay\fontsize{11}{12.5}\selectfont\color{white}\MakeUppercase{Télécharge OnBuch}}\par\vspace{4pt}
        {\raggedright\popsemi\fontsize{8.4}{11}\selectfont\color{white}Gratuit \textbullet\ Google Play\\Tuteur IA, annales, concours, résultats\par}
      \end{minipage}
    \end{tcolorbox}
  \end{minipage}\hfill
  \begin{minipage}[t]{0.487\linewidth}
    \begin{tcolorbox}[enhanced,colback=poppurple,colframe=popink,boxrule=2.2pt,arc=10pt,
        drop shadow=popink,left=11pt,right=11pt,top=10pt,bottom=10pt,height=3.1cm,valign=center]
      \tikz[baseline=-0.7ex]\node[circle,fill=white,draw=popink,line width=1.3pt,minimum size=1.6cm,inner sep=0]{\popdisplay\fontsize{24}{24}\selectfont\color{poppurple}+};%
      \hspace{8pt}\begin{minipage}[c]{0.6\linewidth}
        {\popdisplay\fontsize{11}{12.5}\selectfont\color{white}\MakeUppercase{Passe à OnBuch+}}\par\vspace{4pt}
        {\raggedright\popsemi\fontsize{8.4}{11}\selectfont\color{white}Tous les cours signés, Tuteur IA illimité, fiches PDF — onglet OnBuch+ de l'app\par}
      \end{minipage}
    \end{tcolorbox}
  \end{minipage}
  \vspace{8pt}
  \popcontenu
  \vfill
  \signatureonbuch
  \restoregeometry
  \newpage
}

% Rangée « Ce cours contient » (page de garde)
\newcommand{\poptile}[3]{% #1 couleur #2 gros texte #3 légende
  \begin{tcolorbox}[enhanced,colback=#1,colframe=popink,boxrule=2pt,arc=9pt,shadow={2.5pt}{-2.5pt}{0pt}{popink},
      left=6pt,right=6pt,top=9pt,bottom=9pt,halign=center,valign=center,height=1.75cm,width=\linewidth,nobeforeafter]
    {\popdisplay\fontsize{12.5}{13}\selectfont\color{white}\MakeUppercase{#2}}\par\vspace{3pt}
    {\centering\popsemi\fontsize{7.8}{9.5}\selectfont\color{white}#3\par}
  \end{tcolorbox}}
\newcommand{\popcontenu}{%
  \par\noindent{\popxboldls\fontsize{7.5}{7.5}\selectfont\color{popmuted}CE COURS SIGNÉ CONTIENT}\par\vspace{7pt}
  \noindent\begin{minipage}[t]{0.235\linewidth}\poptile{popblue}{Cours}{complet, illustré, pas à pas}\end{minipage}\hfill
  \begin{minipage}[t]{0.235\linewidth}\poptile{poporange}{Méthodes}{et exercices résolus}\end{minipage}\hfill
  \begin{minipage}[t]{0.235\linewidth}\poptile{poppink}{Exercices}{type examen + corrigés}\end{minipage}\hfill
  \begin{minipage}[t]{0.235\linewidth}\poptile{popgreen}{Fiche bilan}{l'essentiel à retenir}\end{minipage}\par}

% Sommaire
\newtcolorbox{sommaire}{enhanced,colback=white,colframe=popink,boxrule=2.2pt,arc=10pt,
  drop shadow=popink,left=18pt,right=18pt,top=13pt,bottom=13pt,before skip=6pt,after skip=20pt,
  before upper={\poppill{Sommaire}{poppurple}{white}\par\vspace{9pt}\popsemi\fontsize{10.6}{14.5}\selectfont\color{popink}}}

% Signature de fin de cours — poussée tout en bas de la dernière page
\newcommand{\findecours}{%
  \par\vfill\nopagebreak
  \begin{center}\popsquiggle{poporange}\end{center}\vspace{4pt}
  \signatureonbuch
}
\AtBeginDocument{\def\llbracket{\mathopen{[\mkern-3.5mu[}}\def\rrbracket{\mathclose{]\mkern-3.5mu]}}} % intervalles d'entiers (glyphe ⟦ absent de la police)
% Glyphes absents de Fira Math : redéfinis à partir de glyphes disponibles
\AtBeginDocument{%
  \def\setminus{\mathbin{\backslash}}\let\smallsetminus\setminus
  \def\vdots{\mathord{\vbox{\baselineskip3.2pt\lineskiplimit0pt\kern4pt\hbox{.}\hbox{.}\hbox{.}}}}%
  \def\ddots{\mathinner{\mkern1mu\raise6.5pt\hbox{.}\mkern2mu\raise3.6pt\hbox{.}\mkern2mu\raise0.7pt\hbox{.}\mkern1mu}}%
  \def\triangle{\mathord{\scalebox{0.9}{$\symup{\Delta}$}}}%
  \def\nabla{\mathord{\raisebox{\depth}{\scalebox{1}[-1]{$\symup{\Delta}$}}}}%
  \def\checkmark{\popcheck{popdark}}%
  \def\star{\mathbin{\ast}}\def\bigstar{\mathord{\ast}}%
  \def\diamond{\mathbin{\scalebox{0.6}{$\Diamond$}}}%
  \def\bigcup{\mathop{\mathchoice{\raisebox{-0.3em}{\scalebox{1.7}{$\cup$}}}{\scalebox{1.25}{$\cup$}}{\cup}{\cup}}\displaylimits}%
  \def\bigcap{\mathop{\mathchoice{\raisebox{-0.3em}{\scalebox{1.7}{$\cap$}}}{\scalebox{1.25}{$\cap$}}{\cap}{\cap}}\displaylimits}%
  \def\lesseqqgtr{\mathrel{\lessgtr}}\def\bigoplus{\mathop{\oplus}\displaylimits}%
}
\providecommand{\ssi}{si et seulement si} % abréviation fréquente
\AtBeginDocument{\providecommand{\wideparen}{\overparen}} % arc de cercle

%% FIGFIX-BEGIN v5 — correctifs globaux des figures (docs/onbuch-generation/tools/figfix)
\usepackage{adjustbox}\usepackage{etoolbox}\tcbuselibrary{raster}
% toute figure TikZ/pgfplots plus large que la ligne est réduite à la largeur disponible
\BeforeBeginEnvironment{tikzpicture}{\begin{adjustbox}{max width=\linewidth}}
\AfterEndEnvironment{tikzpicture}{\end{adjustbox}}
% légendes pgfplots lisibles par-dessus les courbes
\pgfplotsset{every axis/.append style={legend style={fill=white,fill opacity=0.92,text opacity=1,draw=black!15,inner sep=3pt}}}
% étiquettes d'arêtes (syntaxe edge["..."]) avec fond blanc
\tikzset{every edge quotes/.append style={fill=white,inner sep=1.2pt}}
%% FIGFIX-END
\begin{document}

\pagegarde{Circulation sanguine}{SVTEEHB}{3ème}{SVTEEHB – classe de 3ème}{N09}{7 séances (fiche de progression harmonisée)}{Cette leçon étudie l'appareil circulatoire : les vaisseaux sanguins (artères, veines, capillaires) et le cœur, organe moteur de la circulation du sang. Elle aborde ensuite l'hygiène de la circulation : prévention et premiers secours face aux hémorragies, aux maladies cardiovasculaires et aux accidents vasculaires cérébraux (AVC).
\vspace{4pt}
\begin{multicols}{2}\small
\begin{itemize}
\item À la fin de cette leçon, je sais décrire la structure et le rôle des trois types de vaisseaux sanguins (artères, veines, capillaires).
\item À la fin de cette leçon, je sais décrire la structure externe et interne du cœur à partir d'une dissection réelle ou de schémas.
\item À la fin de cette leçon, je sais schématiser et légender le double circuit de la circulation sanguine (petite et grande circulation).
\item À la fin de cette leçon, je sais interpréter des faits expérimentaux (expériences de Harvey, mesures de pression, pouls) pour établir le sens de la circulation et le rôle du cœur.
\item À la fin de cette leçon, je sais identifier une hémorragie externe et appliquer les gestes de premiers secours (compression, point de compression, garrot en dernier recours).
\item À la fin de cette leçon, je sais citer les principales maladies cardiovasculaires et leurs facteurs de risque, et proposer des règles d'hygiène de vie.
\end{itemize}\end{multicols}}

\begin{sommaire}
\begin{multicols}{2}
\begin{enumerate}
\item Les vaisseaux sanguins, siège de la circulation
\item Le cœur, organe moteur de la circulation
\item Les hémorragies et les moyens de lutte
\item Maladies cardiovasculaires : causes et prévention
\item Les accidents vasculaires cérébraux (AVC) et les premiers secours
\item Activité d'intégration
\item Méthodes et exercices résolus
\item Exercices
\item Corrigés des exercices
\item Fiche bilan
\end{enumerate}
\end{multicols}
\end{sommaire}

\begin{objectifs}
\begin{itemize}
\item À la fin de cette leçon, je sais décrire la structure et le rôle des trois types de vaisseaux sanguins (artères, veines, capillaires).
\item À la fin de cette leçon, je sais décrire la structure externe et interne du cœur à partir d'une dissection réelle ou de schémas.
\item À la fin de cette leçon, je sais schématiser et légender le double circuit de la circulation sanguine (petite et grande circulation).
\item À la fin de cette leçon, je sais interpréter des faits expérimentaux (expériences de Harvey, mesures de pression, pouls) pour établir le sens de la circulation et le rôle du cœur.
\item À la fin de cette leçon, je sais identifier une hémorragie externe et appliquer les gestes de premiers secours (compression, point de compression, garrot en dernier recours).
\item À la fin de cette leçon, je sais citer les principales maladies cardiovasculaires et leurs facteurs de risque, et proposer des règles d'hygiène de vie.
\item À la fin de cette leçon, je sais reconnaître les signes d'un AVC et décrire la conduite à tenir.
\item À la fin de cette leçon, je sais rédiger et justifier une démarche scientifique (hypothèse, interprétation de résultats, conclusion).
\end{itemize}
\end{objectifs}

\begin{prerequis}
\begin{itemize}
\item Le sang : composition et rôles (leçon sur le sang, début d'année de 3ème).
\item Les besoins des cellules en dioxygène et en nutriments, le rejet du dioxyde de carbone (notion de 4ème sur la respiration).
\item Notion d'organe et d'appareil (4ème).
\item Lecture et interprétation de tableaux de résultats et de courbes simples.
\item Hygiène et santé : notions de prévention et de vaccination (4ème).
\end{itemize}
\end{prerequis}

\newpage

\coursec{Les vaisseaux sanguins, siège de la circulation}

Au marché Mokolo à Yaoundé, un vendeur se blesse profondément à la jambe avec une machette : le sang jaillit par saccades, rouge vif. Quelques semaines plus tard, un commerçant de 52 ans s'effondre brutalement au même marché, le bras gauche paralysé : on parle d'« attaque ». Ces deux accidents, si différents en apparence, mettent en cause le même système : l'\cle{appareil circulatoire}. Pourquoi le sang circule-t-il dans tout le corps ? Par quels conduits passe-t-il ? Quelle est la « pompe » qui le propulse ? Et comment réagir face à une hémorragie ou à un AVC pour sauver une vie ? Cette leçon répond à ces questions vitales. Nous commençons par les conduits : les \cle{vaisseaux sanguins}.

\courssub{Mise en évidence expérimentale : le sang circule dans des conduits}

\textbf{Observation.} Le sang est partout dans le corps : la moindre coupure, même superficielle, fait apparaître du sang. Mais ce sang est-il répandu librement entre les organes, ou circule-t-il dans des conduits bien définis ? Et s'il circule, dans quel sens ? Pour répondre, il faut observer un tissu vivant et transparent.

\begin{experience}[Le sang circule dans la queue d'un têtard]
\textbf{Matériel :} un têtard vivant (ou un très petit poisson, dont la nageoire caudale est transparente), une lame de verre, un peu d'eau, un microscope ou une loupe binoculaire.\par
\textbf{Protocole :}
\begin{enumerate}
\item Placer délicatement le têtard sur la lame avec une goutte d'eau, de façon à étaler la membrane de sa queue (membrane natatoire), tissu très fin et transparent.
\item Observer au microscope à faible grossissement, puis à fort grossissement.
\item Noter le sens de déplacement des éléments figurés du sang (globules rouges) dans les différents conduits.
\end{enumerate}
\textbf{Observations :} on voit des conduits de \emph{diamètres très variés} : de gros conduits, des conduits moyens, et de très fins capillaires où les globules rouges passent \emph{en file indienne}. Dans chaque conduit, le sang s'écoule \emph{toujours dans le même sens} : il quitte les gros conduits, se divise dans des conduits de plus en plus fins, traverse le réseau des capillaires, puis se rassemble dans des conduits de plus en plus gros qui ramènent le sang vers le corps.\par
\textbf{Interprétation :} le sang n'est pas répandu librement dans les tissus : il circule dans des \emph{conduits fermés}, les vaisseaux sanguins, et son mouvement a un \emph{sens unique et permanent}.
\end{experience}

\begin{aretenir}[Ce que l'observation établit]
Le sang circule en permanence dans des conduits appelés \cle{vaisseaux sanguins}, toujours dans le même sens. Il existe trois catégories de vaisseaux, distinguées par leur diamètre et leur rôle : les \cle{artères}, les \cle{capillaires} et les \cle{veines}.
\end{aretenir}

\courssub{Les trois types de vaisseaux sanguins}

\begin{definition}[Artère, veine, capillaire]
\begin{itemize}
\item Une \cle{artère} est un vaisseau qui transporte le sang \textbf{du cœur vers les organes}. Elle se ramifie en artères de plus en plus fines, les artérioles.
\item Un \cle{capillaire} est un vaisseau microscopique qui relie les plus fines artères aux plus fines veines. Les capillaires forment un immense réseau au contact des cellules : c'est là que se font les \textbf{échanges} entre le sang et les organes (dioxygène, nutriments, déchets).
\item Une \cle{veine} est un vaisseau qui ramène le sang \textbf{des organes vers le cœur}. Les veinules se réunissent en veines de plus en plus grosses.
\end{itemize}
\textbf{Critère de distinction :} c'est le \emph{sens de la circulation} (vers le cœur ou depuis le cœur) qui définit une artère ou une veine, et \emph{jamais} la couleur ou la nature du sang transporté.
\end{definition}

\begin{attention}[Erreur fréquente : « artère = sang rouge, veine = sang bleu »]
Non ! Une artère ne transporte pas forcément du sang riche en dioxygène. Contre-exemple célèbre : l'\emph{artère pulmonaire} transporte du sang \emph{pauvre} en dioxygène, du cœur vers les poumons ; inversement, les \emph{veines pulmonaires} ramènent du sang \emph{riche} en dioxygène des poumons vers le cœur. Retiens : \textbf{artère = qui part du cœur ; veine = qui arrive au cœur}. Point final.
\end{attention}

\courssub{Structure comparée des parois des vaisseaux}

Chaque type de vaisseau possède une paroi \emph{adaptée à sa fonction}. C'est un grand principe en biologie : la structure épouse la fonction.

\begin{itemize}
\item \textbf{Artère :} paroi \emph{épaisse}, musculeuse et \emph{élastique}. Elle résiste à la forte pression du sang chassé par le cœur et se détend puis se rétracte à chaque battement : c'est ce qui explique le sang qui « jaillit par saccades » lors de la blessure du vendeur de Mokolo.
\item \textbf{Veine :} paroi \emph{mince}, peu musculeuse, car le sang y circule à faible pression. Les veines des membres possèdent des \cle{valvules}, sortes de petits clapets \emph{anti-retour} qui empêchent le sang de redescendre sous l'effet de la pesanteur.
\item \textbf{Capillaire :} paroi réduite à \emph{une seule assise de cellules} très fines. Cette extrême minceur laisse passer facilement les échanges entre le sang et les cellules : c'est le siège des échanges.
\end{itemize}

\begin{exemplebox}[Des diamètres très différents]
L'aorte, la plus grosse artère du corps, mesure environ \SI{2,5}{cm} de diamètre. Un capillaire mesure environ \SI{8}{\micro\metre}, soit à peine le diamètre d'un globule rouge : les globules doivent s'y déformer et passer \emph{un par un}, ce qui ralentit le sang et favorise les échanges. Rapport des diamètres : $\dfrac{\num{2,5}\ \text{cm}}{\num{8}\ \mu\text{m}} = \dfrac{\num{25000}\ \mu\text{m}}{\num{8}\ \mu\text{m}} \approx \num{3000}$. L'aorte est environ \num{3000} fois plus large qu'un capillaire !
\end{exemplebox}

\begin{popfigure}
\begin{center}
\begin{tikzpicture}[scale=1]
% Artère
\begin{scope}[shift={(0,0)}]
\fill[poporangeL] (0,0) circle (1.5);
\fill[poporange!70] (0,0) circle (1.15);
\fill[white] (0,0) circle (0.75);
\node at (0,-2.1) {\textbf{Artère}};
\node[align=center,font=\small] at (0,-2.7) {paroi épaisse,\\élastique};
\draw[->,popdark] (1.9,0.9) -- (1.15,0.55);
\node[font=\small,popdark,align=left] at (3.1,1.1) {paroi\\épaisse};
\end{scope}
% Capillaire
\begin{scope}[shift={(5.5,0)}]
\fill[poppinkL] (0,0) circle (0.55);
\fill[white] (0,0) circle (0.42);
\node at (0,-2.1) {\textbf{Capillaire}};
\node[align=center,font=\small] at (0,-2.7) {1 seule assise\\de cellules};
\draw[->,popdark] (0.9,0.5) -- (0.5,0.28);
\node[font=\small,popdark,align=left] at (1.9,0.7) {paroi très\\fine};
\end{scope}
% Veine
\begin{scope}[shift={(10.5,0)}]
\fill[popblueL] (0,0) circle (1.5);
\fill[popblue!40] (0,0) circle (1.32);
\fill[white] (0,0) circle (1.15);
\node at (0,-2.1) {\textbf{Veine}};
\node[align=center,font=\small] at (0,-2.7) {paroi mince,\\avec valvules};
\draw[->,popdark] (1.9,-0.9) -- (1.32,-0.6);
\node[font=\small,popdark,align=left] at (3.1,-1.1) {paroi\\mince};
\end{scope}
\end{tikzpicture}
\end{center}
\legende{Coupes transversales comparées des trois types de vaisseaux (les épaisseurs de paroi sont respectées à l'échelle relative). Artère : paroi épaisse et élastique, lumière étroite. Capillaire : paroi d'une seule couche de cellules. Veine : paroi mince, grande lumière, valvules anti-retour.}
\end{popfigure}

\courssub{L'expérience historique de William Harvey (1628) : le sang circule en circuit fermé}

Pendant des siècles, on a cru que le sang était produit en permanence par le foie puis « consommé » par les organes. En 1628, le médecin anglais William \cle{Harvey} démontre que c'est impossible, par une expérience simple et un calcul lumineux.

\begin{experience}[La ligature du bras de Harvey]
\textbf{Protocole :} on serre un garrot (ligature) autour du bras, assez lâche pour laisser passer le sang artériel profond mais assez serré pour bloquer le sang des veines superficielles.\par
\textbf{Observations :}
\begin{enumerate}
\item En \emph{aval} de la ligature (côté main), les veines \emph{gonflent} : le sang continue d'arriver dans la main par les artères mais ne peut plus repartir par les veines.
\item Le long des veines gonflées apparaissent des \emph{renflements réguliers} : ce sont les \emph{valvules}. Si l'on chasse le sang d'un segment de veine vers le cœur avec le doigt, le segment reste vide : le sang ne revient pas en arrière, bloqué par la valvule.
\end{enumerate}
\textbf{Interprétation :} le sang des veines circule bien \emph{des organes vers le cœur}, et les valvules imposent un \emph{sens unique} de circulation. Le sang ne peut pas être « consommé » sur place : il doit revenir au cœur.
\end{experience}

\begin{popfigure}
\begin{center}
\begin{tikzpicture}[scale=1]
% avant-bras
\draw[thick,popdark,fill=popcream] (0,0) -- (6.5,0) .. controls (7.3,0.1) and (7.4,0.9) .. (6.6,1.3) -- (0,1.3) -- cycle;
\node[font=\small] at (7.0,1.7) {main};
% artère (profonde)
\draw[very thick,poporange] (0,0.35) -- (6.3,0.35);
\draw[->,very thick,poporange] (3.0,0.35) -- (3.8,0.35);
\node[font=\small,poporange] at (1.2,0.05) {artère (profonde)};
% veine (superficielle)
\draw[very thick,popblue] (0,0.95) -- (6.3,0.95);
\draw[->,very thick,popblue] (3.8,0.95) -- (3.0,0.95);
\node[font=\small,popblue] at (1.2,1.55) {veine (superficielle)};
% ligature
\draw[line width=2.5pt,popink] (2.0,-0.15) -- (2.0,1.45);
\node[font=\small,popink,align=center] at (2.0,-0.5) {ligature\\(garrot)};
% valvules gonflées en aval
\foreach \x in {3.2,4.3,5.4}{
  \draw[popblue,fill=popblueL] (\x,0.95) ellipse (0.18 and 0.28);
}
\node[font=\small,popblue,align=center] at (4.3,2.0) {veine gonflée en aval :\\les valvules apparaissent};
\draw[->,popblue] (4.3,1.6) -- (4.3,1.25);
% sens
\node[font=\small,popdark] at (5.6,0.05) {vers la main};
\end{tikzpicture}
\end{center}
\legende{Expérience de Harvey : la ligature du bras bloque le retour veineux superficiel. En aval (côté main), la veine gonfle et les valvules apparaissent en renflements : le sang veineux circule bien des organes vers le cœur et ne peut pas refluer.}
\end{popfigure}

\begin{savaistu}[Le calcul qui a changé l'histoire de la médecine]
Harvey a mesuré que le cœur éjecte environ \SI{60}{mL} de sang à chaque battement. Avec environ 70 battements par minute, cela fait $60 \times 70 = \num{4200}\ \text{mL} \approx 4\ \text{litres}$ par minute, soit plus de \SI{240}{L} par heure : une masse de sang \emph{bien supérieure à la masse totale du corps} (environ 5 litres de sang seulement chez l'adulte) ! Le foie ne pourrait jamais fabriquer autant de sang. Conclusion inévitable : le même sang est \emph{recyclé} sans cesse — il circule en \textbf{circuit fermé}.
\end{savaistu}

\begin{propriete}[La circulation est un circuit fermé à sens unique]
Établie à partir de l'observation du têtard et de l'expérience de Harvey : le sang circule dans un \textbf{circuit fermé} (cœur $\rightarrow$ artères $\rightarrow$ capillaires $\rightarrow$ veines $\rightarrow$ cœur), toujours \textbf{dans le même sens}. Ce sens unique est garanti par la \emph{pompe cardiaque} qui propulse le sang, et par les \emph{valvules} des veines qui empêchent tout retour en arrière.
\end{propriete}

\courssub{Pression artérielle et pouls}

À chaque contraction, le cœur chasse le sang sous pression dans les artères : c'est la \cle{pression artérielle}. Elle est maximale dans l'aorte au moment de la contraction, puis diminue tout au long du circuit : elle devient très faible dans les capillaires (pour ne pas les déchirer et pour laisser le temps aux échanges) et minimale dans les veines.

\begin{popfigure}
\begin{center}
\begin{tikzpicture}[scale=0.95]
\begin{axis}[
  width=0.85\linewidth, height=5.5cm,
  xlabel={Vaisseaux le long du circuit},
  ylabel={Pression (mmHg)},
  symbolic x coords={Aorte,Artères,Artérioles,Capillaires,Veines},
  xtick=data, ymin=0, ymax=130,
  ytick={0,20,40,60,80,100,120},
  grid=both, grid style={popline!40},
  tick label style={font=\small}, label style={font=\small}
]
\addplot[popcurve,mark=*,mark size=2pt] coordinates {(Aorte,100) (Artères,85) (Artérioles,55) (Capillaires,25) (Veines,8)};
\end{axis}
\end{tikzpicture}
\end{center}
\legende{Évolution de la pression sanguine moyenne le long du circuit (valeurs indicatives chez l'adulte au repos). La pression, élevée dans l'aorte, chute fortement dans les artérioles et devient très faible dans les veines.}
\end{popfigure}

L'onde de pression produite par chaque battement se propage dans les artères : on la sent en appuyant légèrement une artère contre un os, c'est le \cle{pouls}. La pression se mesure au bras avec un brassard (tensiomètre) : on note deux nombres, par exemple 12/8 (pression maximale / pression minimale).

\begin{methode}[Prendre le pouls au poignet]
\begin{enumerate}
\item Retourner le poignet, paume vers le haut, bras détendu.
\item Poser deux doigts (index et majeur, \emph{jamais le pouce} qui a son propre pouls) sur la face interne du poignet, du côté du pouce, sur l'artère radiale.
\item Appuyer légèrement jusqu'à sentir les battements réguliers.
\item Compter les battements pendant 30 secondes et multiplier par 2 (ou pendant 1 minute entière si le rythme est irrégulier).
\item Noter le résultat en battements par minute (batt/min). Au repos, chez l'adolescent, il est d'environ 60 à 90 batt/min.
\end{enumerate}
\end{methode}

\begin{exoresolu}[Le pouls reflète-t-il l'activité du cœur ?]
\textbf{Énoncé.} Un élève de 3ème compte 34 pulsations au poignet en 30 secondes au repos, puis 51 pulsations en 30 secondes juste après avoir monté quatre étages en courant. 1) Calculer sa fréquence cardiaque dans les deux cas. 2) Expliquer pourquoi le pouls permet de connaître l'activité du cœur.
\tcblower
\textbf{Solution.}
\begin{enumerate}
\item Au repos : $34 \times 2 = 68$ batt/min. Après l'effort : $51 \times 2 = 102$ batt/min.
\item Chaque pulsation sentie au poignet correspond à l'onde de pression produite par \emph{une} contraction du cœur, qui se propage dans les artères jusqu'à l'artère radiale. Le nombre de pulsations par minute est donc égal au nombre de battements du cœur par minute : le pouls est un reflet fidèle de la fréquence cardiaque. L'augmentation après l'effort montre que le cœur accélère pour apporter plus de sang (donc plus de dioxygène) aux muscles qui travaillent.
\end{enumerate}
\end{exoresolu}

\begin{savaistu}[L'hypertension artérielle, un problème de santé au Cameroun]
On parle d'\emph{hypertension artérielle} lorsque la tension dépasse durablement environ 14/9 au repos. Très fréquente chez l'adulte camerounais (souvent favorisée par l'excès de sel, l'alcool, le stress et le manque d'activité physique), elle use prématurément le cœur et les vaisseaux et expose aux accidents vasculaires cérébraux. Elle ne provoque souvent aucun symptôme : d'où l'importance de faire mesurer sa tension régulièrement, même quand on se sent bien. (Complément utile, au-delà du strict programme.)
\end{savaistu}

\begin{aretenir}[L'essentiel de la section]
\begin{itemize}
\item Le sang circule dans des \textbf{vaisseaux sanguins}, toujours dans le même sens (observation du têtard).
\item Trois types de vaisseaux : \textbf{artères} (cœur $\rightarrow$ organes, paroi épaisse et élastique), \textbf{capillaires} (réseau microscopique, paroi d'une seule assise de cellules, siège des échanges), \textbf{veines} (organes $\rightarrow$ cœur, paroi mince, valvules anti-retour).
\item Harvey (1628) a démontré que le sang circule en \textbf{circuit fermé} : il ne peut pas être « consommé », il est recyclé.
\item La \textbf{pression artérielle}, maximale dans l'aorte, diminue le long du circuit ; le \textbf{pouls} (au poignet, artère radiale) reflète les battements du cœur.
\end{itemize}
\end{aretenir}

\begin{exercice}[Pour vérifier tes acquis]
1) Un vaisseau ramène le sang des poumons vers le cœur : est-ce une artère ou une veine ? Justifier. \quad 2) Pourquoi la paroi des capillaires est-elle si fine ? \quad 3) Pourquoi le sang jaillit-il « par saccades » quand une artère est sectionnée, alors qu'il s'écoule en nappe continue quand une veine est coupée ? \quad 4) Citer l'expérience qui prouve que le sang veineux circule vers le cœur et décrire ses observations.
\end{exercice}

\coursec{Le cœur, organe moteur de la circulation}

Dans la section précédente, nous avons découvert que le sang circule en circuit fermé dans un réseau de vaisseaux sanguins : artères, veines et capillaires. Mais une question demeure : \emph{qu'est-ce qui pousse le sang dans tout ce réseau ?} Un liquide ne circule pas tout seul dans des tuyaux : il faut une pompe. Cette pompe, c'est le \cle{cœur}, un organe musculaire qui bat sans relâche, du début à la fin de la vie. Cette section est consacrée à l'étude de sa structure, de son fonctionnement et de son rôle de moteur de la double circulation.

\courssub{TP n°7 : dissection d'un cœur de mammifère}

Pour comprendre comment le cœur fonctionne, le mieux est de l'observer directement. Au laboratoire, on utilise un cœur de bœuf, de mouton ou de chèvre, organes faciles à se procurer chez le boucher et dont la structure est semblable à celle du cœur humain.

\begin{methode}[Protocole de dissection d'un cœur de mammifère]
\begin{enumerate}
\item \textbf{Matériel} : un cœur frais de mammifère, une planche à disséquer, un scalpel (ou couteau bien affûté), des ciseaux fins, des pinces, des gants, un entonnoir, un récipient d'eau.
\item \textbf{Observation externe} : noter la forme, la taille, la couleur ; repérer la graisse blanchâtre à la surface, les deux oreillettes et les gros vaisseaux sectionnés (aorte, artère pulmonaire, veines caves).
\item \textbf{Section frontale} : poser le cœur sur la planche et le couper en deux dans le sens de la longueur, de haut en bas, de façon à séparer les cavités droites des cavités gauches.
\item \textbf{Observation interne} : identifier les quatre cavités, comparer l'épaisseur des parois, repérer les valvules entre les cavités et à la base des gros vaisseaux.
\item \textbf{Test des valvules} : placer un entonnoir dans l'aorte et verser de l'eau : observer que l'eau ne pénètre pas dans le ventricule (les valvules se ferment).
\item \textbf{Compte rendu} : dessiner et légender la coupe, rédiger les observations et les conclusions.
\end{enumerate}
\end{methode}

\begin{attention}[Règles de sécurité au laboratoire]
Le scalpel est un instrument tranchant : on coupe toujours \textbf{en s'éloignant de soi}, sur la planche à disséquer, jamais l'organe tenu en main. On porte des \textbf{gants} pour manipuler l'organe frais, on ne porte rien à la bouche, et on se \textbf{lave soigneusement les mains} au savon après la séance. Le matériel est nettoyé et rangé, les déchets biologiques sont éliminés selon les consignes du professeur.
\end{attention}

\courssub{Observation externe du cœur}

Le cœur a la forme d'un \textbf{cône} ou d'un poing fermé, pointe dirigée vers le bas et vers la gauche. Sa taille est remarquable : il est à peu près \textbf{gros comme le poing fermé de son propriétaire}. Chez l'Homme adulte, il mesure environ 12 cm de long et pèse environ 300 g.

À sa surface, on observe :
\begin{itemize}
\item une couche de \textbf{graisse} blanchâtre qui le protège ;
\item deux petites cavités supérieures à paroi fine, les \cle{oreillettes} ;
\item les \textbf{gros vaisseaux} sectionnés à la base du cœur : l'\textbf{aorte} (la plus grosse artère du corps, à paroi très épaisse), l'\textbf{artère pulmonaire} (qui part vers les poumons), les \textbf{veines caves} (qui ramènent le sang du corps) et les \textbf{veines pulmonaires} (qui ramènent le sang des poumons).
\end{itemize}

\courssub{Observation interne : quatre cavités et des parois inégales}

Après la section frontale, on découvre que le cœur est un organe \textbf{creux}, divisé en \textbf{quatre cavités} :
\begin{itemize}
\item en haut, deux \textbf{oreillettes} (droite et gauche), à paroi mince ;
\item en bas, deux \textbf{ventricules} (droit et gauche), à paroi épaisse.
\end{itemize}

Le cœur droit et le cœur gauche sont \textbf{complètement séparés} par une cloison : le sang ne passe jamais directement de l'un à l'autre.

L'observation la plus frappante est la différence d'épaisseur des parois : la paroi de l'\textbf{oreillette} est mince, celle du \textbf{ventricule droit} est épaisse, et celle du \textbf{ventricule gauche} est \emph{encore plus épaisse} (environ trois fois plus que celle du ventricule droit). Cette paroi épaisse est constituée d'un muscle, le \cle{myocarde}.

\begin{definition}[Myocarde, systole, diastole, fréquence cardiaque]
\begin{itemize}
\item Le \cle{myocarde} est le muscle qui forme la paroi du cœur. En se contractant, il chasse le sang hors des cavités.
\item La \cle{systole} est la phase de \textbf{contraction} du cœur : le sang est éjecté dans les artères.
\item La \cle{diastole} est la phase de \textbf{relâchement} du cœur : les cavités se remplissent de sang.
\item La \cle{fréquence cardiaque} est le nombre de battements du cœur par minute. Au repos, elle vaut environ \textbf{60 à 80 battements par minute} chez un adolescent ou un adulte.
\end{itemize}
\end{definition}

\courssub{Les valvules cardiaques : des clapets anti-retour}

Entre chaque oreillette et le ventricule correspondant, et à l'entrée des grosses artères, on observe des membranes fines et résistantes : les \cle{valvules cardiaques}.
\begin{itemize}
\item la valvule \textbf{tricuspide} (à trois pointes) entre l'oreillette droite et le ventricule droit ;
\item la valvule \textbf{mitrale} (ou bicuspide, à deux pointes) entre l'oreillette gauche et le ventricule gauche ;
\item les valvules \textbf{sigmoïdes}, en forme de nids d'oiseau, à l'entrée de l'aorte et de l'artère pulmonaire.
\end{itemize}

\begin{experience}[Le test de l'entonnoir : rôle anti-retour des valvules]
\textbf{Protocole} : on introduit un entonnoir dans l'aorte d'un cœur disséqué et on y verse de l'eau.\par
\textbf{Observation} : l'eau ne pénètre pas dans le ventricule gauche ; elle déborde de l'entonnoir. Les valvules sigmoïdes se sont gonflées et fermées comme un clapet.\par
\textbf{Interprétation} : les valvules ne s'ouvrent que \textbf{dans un seul sens}. Elles laissent passer le sang de l'oreillette vers le ventricule, puis du ventricule vers l'artère, mais elles \textbf{empêchent tout retour en arrière}.\par
\textbf{Conclusion} : grâce aux valvules, le sang circule toujours dans le même sens : veines $\rightarrow$ oreillette $\rightarrow$ ventricule $\rightarrow$ artère.
\end{experience}

\begin{popfigure}
\begin{center}
\begin{tikzpicture}[scale=1.05, every node/.style={font=\small}]
% Paroi du coeur (coupe frontale simplifiee)
\draw[thick, poppink, fill=poppinkL] (0,0.2)
  .. controls (-2.6,0.4) and (-2.8,2.2) .. (-2.2,3.6)
  .. controls (-1.9,4.4) and (-1.2,4.8) .. (-0.6,4.8)
  -- (-0.6,5.6) -- (0.6,5.6) -- (0.6,4.8)
  .. controls (1.2,4.8) and (1.9,4.4) .. (2.2,3.6)
  .. controls (2.8,2.2) and (2.6,0.4) .. (0,0.2) -- cycle;
% Cloison centrale
\draw[thick, poppink, fill=poppinkL] (-0.35,0.4) rectangle (0.35,4.6);
% Cavites internes
\draw[thick, poppink, fill=white] (-1.7,3.1) ellipse (0.75 and 0.85);
\draw[thick, poppink, fill=white] (1.7,3.1) ellipse (0.75 and 0.85);
\draw[thick, poppink, fill=white] (-1.55,1.35) ellipse (0.8 and 1.05);
\draw[thick, poppink, fill=white] (1.55,1.35) ellipse (0.8 and 1.05);
% Gros vaisseaux
\draw[thick, popblue, fill=popblueL] (-2.5,4.1) rectangle (-1.9,5.5);
\draw[thick, poppink, fill=poppinkL] (1.9,4.1) rectangle (2.5,5.5);
\draw[thick, popblue, fill=popblueL] (-0.95,4.5) rectangle (-0.4,5.7);
\draw[thick, poppink, fill=poppinkL] (0.4,4.5) rectangle (0.95,5.7);
% Valvules
\draw[thick, popdark] (-1.95,2.35) -- (-1.55,2.05) -- (-1.15,2.35);
\draw[thick, popdark] (1.15,2.35) -- (1.55,2.05) -- (1.95,2.35);
% Fleches du sens du sang (cote droit = sang desoxygene)
\draw[->, very thick, popblue] (-1.7,4.6) -- (-1.7,3.9);
\draw[->, very thick, popblue] (-1.7,2.9) -- (-1.7,2.45);
\draw[->, very thick, popblue] (-1.55,1.9) -- (-1.55,1.0);
\draw[->, very thick, popblue] (-1.1,1.6) .. controls (-0.9,2.6) .. (-0.68,4.4);
% Fleches cote gauche = sang oxygene
\draw[->, very thick, poppink] (1.7,4.6) -- (1.7,3.9);
\draw[->, very thick, poppink] (1.7,2.9) -- (1.7,2.45);
\draw[->, very thick, poppink] (1.55,1.9) -- (1.55,1.0);
\draw[->, very thick, poppink] (1.1,1.6) .. controls (0.9,2.6) .. (0.68,4.4);
% Legendes
\node[align=center] at (-1.7,3.15) {\textbf{Oreillette}\\ \textbf{droite}};
\node[align=center] at (1.7,3.15) {\textbf{Oreillette}\\ \textbf{gauche}};
\node[align=center] at (-1.55,1.35) {\textbf{Ventricule}\\ \textbf{droit}};
\node[align=center] at (1.55,1.35) {\textbf{Ventricule}\\ \textbf{gauche}};
\node[popblue, rotate=90] at (-2.2,4.85) {\footnotesize Veines caves};
\node[poppink, rotate=90] at (2.2,4.85) {\footnotesize Veines pulmonaires};
\node[popblue, rotate=90] at (-0.68,5.15) {\footnotesize Artère pulmonaire};
\node[poppink, rotate=90] at (0.68,5.15) {\footnotesize Aorte};
\node[popdark, font=\footnotesize] at (-3.3,2.1) {Valvule tricuspide};
\draw[->, popdark, thin] (-2.7,2.15) -- (-1.9,2.2);
\node[popdark, font=\footnotesize] at (3.35,2.1) {Valvule mitrale};
\draw[->, popdark, thin] (2.75,2.15) -- (1.9,2.2);
\node[popdark, font=\footnotesize, align=center] at (0,-0.5) {Cloison de séparation};
\draw[->, popdark, thin] (0,-0.35) -- (0,0.35);
\end{tikzpicture}
\end{center}
\legende{Coupe frontale du cœur. Le cœur comporte quatre cavités : deux oreillettes (en haut, paroi mince) et deux ventricules (en bas, paroi épaisse, surtout à gauche). Les valvules (tricuspide à droite, mitrale à gauche, sigmoïdes à la base des artères) imposent un sens unique au sang. En bleu : sang désoxygéné ; en rouge : sang oxygéné. Attention : sur un schéma de coupe frontale, le cœur droit apparaît à gauche de la feuille, comme face à un patient.}
\end{popfigure}

\courssub{Interprétation : une structure adaptée à la fonction}

Pourquoi le ventricule gauche est-il si épais ? Appliquons la démarche d'investigation : nous avons \textbf{observé} que l'épaisseur du myocarde augmente dans l'ordre : oreillette $<$ ventricule droit $<$ ventricule gauche. \textbf{Hypothèse} : plus une cavité doit propulser le sang loin, plus sa paroi musculaire est épaisse.

Vérifions : l'oreillette pousse le sang seulement vers le ventricule voisin (quelques centimètres) ; le ventricule droit envoie le sang jusqu'aux \textbf{poumons}, organes proches ; le ventricule gauche propulse le sang dans \textbf{tout le corps}, de la tête aux pieds, soit le trajet le plus long et le plus difficile. L'hypothèse est confirmée.

\begin{propriete}[Relation entre la structure et la fonction]
L'épaisseur du myocarde d'une cavité cardiaque est \textbf{proportionnelle au travail} qu'elle doit fournir : plus le sang doit être propulsé loin, plus la paroi est épaisse et puissante. C'est un exemple général en biologie : \emph{la structure d'un organe est adaptée à sa fonction}.
\end{propriete}

\courssub{Le cœur, une double pompe : les deux circulations}

Le sang désoxygéné (pauvre en oxygène) revient du corps par les veines caves, entre dans l'oreillette droite, passe dans le ventricule droit qui le chasse vers les \textbf{poumons} par l'artère pulmonaire. Là, le sang se recharge en oxygène. C'est la \cle{petite circulation} (ou circulation pulmonaire) : cœur droit $\rightarrow$ poumons $\rightarrow$ cœur gauche.

Le sang oxygéné revient des poumons par les veines pulmonaires, entre dans l'oreillette gauche, passe dans le ventricule gauche qui le chasse dans l'\textbf{aorte} vers \textbf{tous les organes}. Après avoir distribué l'oxygène, le sang revient au cœur droit. C'est la \cle{grande circulation} (ou circulation générale) : cœur gauche $\rightarrow$ organes $\rightarrow$ cœur droit.

Le cœur fonctionne donc comme une \textbf{double pompe} : la moitié droite alimente les poumons, la moitié gauche alimente tout le reste du corps.

\begin{popfigure}
\begin{center}
\begin{tikzpicture}[scale=1, every node/.style={font=\small}]
% Coeur
\draw[thick, poppink, fill=poppinkL] (-1.5,-1.1) rectangle (-0.1,1.1);
\draw[thick, poppink, fill=poppinkL] (0.1,-1.1) rectangle (1.5,1.1);
\node[align=center] at (-0.8,0) {\textbf{Cœur}\\ \textbf{droit}};
\node[align=center] at (0.8,0) {\textbf{Cœur}\\ \textbf{gauche}};
% Poumons
\draw[thick, popgreen, fill=popgreenL] (-3.4,2.2) ellipse (0.9 and 0.7);
\draw[thick, popgreen, fill=popgreenL] (-1.6,2.2) ellipse (0.9 and 0.7);
\node at (-2.5,2.2) {\textbf{Poumons}};
% Organes
\draw[thick, popgold, fill=popgoldL] (1.2,-2.6) rectangle (3.6,-3.6);
\node at (2.4,-3.1) {\textbf{Organes du corps}};
% Petite circulation : coeur droit -> poumons (sang desoxygene)
\draw[->, very thick, popblue] (-1.5,0.6) .. controls (-3.2,0.8) and (-3.6,1.2) .. (-3.5,1.8);
% poumons -> coeur gauche (sang oxygene)
\draw[->, very thick, poppink] (-1.4,1.9) .. controls (-0.9,1.7) and (0.4,1.7) .. (0.7,1.15);
% Grande circulation : coeur gauche -> organes (sang oxygene)
\draw[->, very thick, poppink] (1.5,-0.5) .. controls (3.0,-0.8) and (3.4,-1.6) .. (3.3,-2.55);
% organes -> coeur droit (sang desoxygene)
\draw[->, very thick, popblue] (1.3,-3.6) .. controls (0.2,-4.2) and (-1.2,-3.4) .. (-1.0,-1.15);
% Etiquettes
\node[popblue, align=center, font=\footnotesize] at (-4.1,0.9) {sang\\ désoxygéné};
\node[poppink, align=center, font=\footnotesize] at (-0.3,2.15) {sang oxygéné};
\node[poppink, align=center, font=\footnotesize] at (3.9,-1.5) {sang oxygéné};
\node[popblue, align=center, font=\footnotesize] at (-0.2,-4.1) {sang désoxygéné};
\node[popdark, font=\footnotesize, align=center] at (-3.4,0.1) {petite circulation};
\node[popdark, font=\footnotesize, align=center] at (3.9,-3.9) {grande circulation};
\end{tikzpicture}
\end{center}
\legende{Schéma fonctionnel du double circuit. Petite circulation : cœur droit $\rightarrow$ poumons $\rightarrow$ cœur gauche (le sang s'y recharge en oxygène). Grande circulation : cœur gauche $\rightarrow$ organes $\rightarrow$ cœur droit (le sang y distribue l'oxygène). En rouge : sang oxygéné ; en bleu : sang désoxygéné.}
\end{popfigure}

\courssub{Le cœur, un muscle qui bat sans relâche : systole, diastole et pouls}

Le cœur est un \textbf{muscle creux} qui se contracte de façon rythmique : à chaque battement, la systole éjecte le sang dans les artères, puis la diastole permet le remplissage. On peut percevoir ce travail en prenant le \cle{pouls} : on pose deux doigts sur l'artère du poignet (face interne, côté pouce) et on compte les battements pendant une minute.

\begin{experience}[Mesure du pouls avant et après un effort]
\textbf{Protocole} : un élève au repos depuis 5 minutes mesure son pouls pendant 1 minute. Il effectue ensuite un effort de montées-descentes sur une marche pendant 3 minutes, puis mesure son pouls immédiatement après l'effort, puis toutes les minutes pendant la récupération.\par
\textbf{Résultats obtenus} (élève de 14 ans) :
\begin{center}
\begin{tabular}{|l|c|c|c|c|c|}
\hline
\rowcolor{popblueL} Situation & Repos & Fin de l'effort & 1 min après & 3 min après & 5 min après \\
\hline
Fréquence cardiaque (battements/min) & 72 & 138 & 105 & 88 & 76 \\
\hline
\end{tabular}
\end{center}
\textbf{Interprétation} : pendant l'effort, les muscles ont besoin de plus d'oxygène ; le cœur accélère pour leur apporter plus de sang. Après l'effort, la fréquence cardiaque diminue progressivement et revient à sa valeur de repos : c'est la \textbf{récupération}.\par
\textbf{Conclusion} : la fréquence cardiaque s'adapte aux besoins de l'organisme ; elle vaut 60 à 80 battements/min au repos et peut dépasser 130 battements/min à l'effort.
\end{experience}

\begin{popfigure}
\begin{center}
\begin{tikzpicture}
\begin{axis}[
  width=0.85\linewidth, height=6cm,
  xlabel={Temps (min)}, ylabel={Fréquence cardiaque (battements/min)},
  xmin=0, xmax=11, ymin=50, ymax=150,
  xtick={0,1,2,3,4,5,6,7,8,9,10}, ytick={60,80,100,120,140},
  grid=both, grid style={popline!40},
  legend style={at={(0.98,0.95)}, anchor=north east, font=\footnotesize}
]
\addplot[popcurve, mark=*, mark size=1.5pt] coordinates {(0,72) (1,72) (2,100) (3,120) (4,138) (5,105) (6,96) (7,88) (8,82) (9,78) (10,76)};
\legend{Élève de 14 ans}
\addplot[popmuted, dashed, domain=0:11] {72};
\node[font=\footnotesize, popmuted, anchor=south west] at (axis cs:0.2,73) {valeur de repos};
\end{axis}
\end{tikzpicture}
\end{center}
\legende{Évolution de la fréquence cardiaque d'un élève au repos (0--1 min), pendant un effort de montées-descentes (1--4 min) puis pendant la récupération (4--10 min). La fréquence augmente pendant l'effort puis revient progressivement à la valeur de repos.}
\end{popfigure}

\begin{exemplebox}[Un cœur infatigable : quelques chiffres]
Au repos, avec une fréquence d'environ 70 battements par minute, le cœur bat environ $70 \times 60 \times 24 \approx 100\,000$ fois par jour. Chaque battement éjecte environ 75 mL de sang : le cœur propulse donc près de \textbf{7 500 litres de sang par jour} au repos, soit l'équivalent de plus de 35 baignoires ! Et cela sans jamais s'arrêter, pendant toute la vie.
\end{exemplebox}

\begin{aretenir}[Savoir admis : l'automatisme cardiaque]
Le cœur bat \textbf{sans ordre conscient} : on ne décide pas de ses battements. Cet \cle{automatisme cardiaque} est une propriété du myocarde lui-même : un cœur de mammifère prélevé et perfusé avec un liquide nutritif oxygéné peut continuer à battre quelque temps \textbf{hors du corps}. Ce phénomène est admis en classe de 3ème ; son explication détaillée sera étudiée dans les classes supérieures.
\end{aretenir}

\begin{attention}[Erreurs fréquentes à éviter]
\begin{itemize}
\item \textbf{Ne pas inverser oreillettes et ventricules} sur un schéma : les oreillettes sont \emph{en haut} (paroi mince), les ventricules \emph{en bas} (paroi épaisse). Sur une coupe frontale, le cœur droit apparaît à gauche de la feuille.
\item Éviter les expressions « sang pur » et « sang impur » : dire correctement sang \textbf{oxygéné} (riche en oxygène, cœur gauche) et sang \textbf{désoxygéné} (pauvre en oxygène, cœur droit). Le sang désoxygéné n'est pas « sale » : il transporte simplement moins d'oxygène et davantage de gaz carbonique.
\item Le sang ne passe \textbf{jamais directement} du cœur droit au cœur gauche : il doit obligatoirement passer par les poumons.
\end{itemize}
\end{attention}

\begin{savaistu}[Le cœur des athlètes]
Comme tout muscle, le myocarde se développe avec l'entraînement. Le cœur d'un sportif de haut niveau est plus \textbf{volumineux} et plus puissant : chaque battement éjecte plus de sang. Résultat : au repos, son cœur bat plus \textbf{lentement}, parfois seulement 40 à 50 battements par minute ! C'est la \emph{bradycardie d'entraînement}, signe d'un cœur en excellente santé. L'activité physique régulière (marche, football, course) est donc un excellent moyen de protéger son cœur.
\end{savaistu}

\begin{exoresolu}[Exploiter des résultats de mesure du pouls]
\textbf{Énoncé.} Deux élèves, Awa (sportive) et Béa (non sportive), mesurent leur pouls au repos puis juste après un même effort de 3 minutes. Résultats : Awa : 58 battements/min au repos, 120 après l'effort ; Béa : 78 battements/min au repos, 145 après l'effort.
\begin{enumerate}
\item Calcule l'augmentation de la fréquence cardiaque pendant l'effort pour chaque élève.
\item Quelle élève a le cœur le plus entraîné ? Justifie ta réponse.
\end{enumerate}
\tcblower
\textbf{Solution détaillée.}
\begin{enumerate}
\item Augmentation pour Awa : $120 - 58 = 62$ battements/min. Augmentation pour Béa : $145 - 78 = 67$ battements/min.
\item C'est \textbf{Awa} qui a le cœur le plus entraîné. Deux indices le montrent : sa fréquence de repos est plus faible (58 contre 78 battements/min), ce qui correspond à la bradycardie d'entraînement, et son cœur accélère moins pour fournir le même effort (120 contre 145 battements/min), car son myocarde, plus puissant, éjecte plus de sang à chaque battement.
\end{enumerate}
\end{exoresolu}

\begin{exercice}[À toi de jouer]
\begin{enumerate}
\item Lors de la dissection d'un cœur de mouton, tu observes qu'une paroi est trois fois plus épaisse que l'autre. De quelle cavité s'agit-il ? Justifie cette particularité.
\item Reconstitue le trajet d'une goutte de sang depuis la veine cave jusqu'à l'aorte, en citant dans l'ordre toutes les cavités, valvules et vaisseaux traversés.
\end{enumerate}
\end{exercice}

\begin{corrige}[Exercice « À toi de jouer »]
\begin{enumerate}
\item Il s'agit du \textbf{ventricule gauche}. Sa paroi est très épaisse car il doit propulser le sang dans \textbf{tout le corps} (grande circulation), ce qui exige un travail bien supérieur à celui du ventricule droit, qui n'envoie le sang qu'aux poumons voisins. L'épaisseur du myocarde est adaptée au travail fourni.
\item Trajet : veine cave $\rightarrow$ oreillette droite $\rightarrow$ valvule tricuspide $\rightarrow$ ventricule droit $\rightarrow$ valvules sigmoïdes $\rightarrow$ artère pulmonaire $\rightarrow$ poumons $\rightarrow$ veines pulmonaires $\rightarrow$ oreillette gauche $\rightarrow$ valvule mitrale $\rightarrow$ ventricule gauche $\rightarrow$ valvules sigmoïdes $\rightarrow$ aorte.
\end{enumerate}
\end{corrige}

\begin{aretenir}[L'essentiel de la section]
Le cœur est un \textbf{muscle creux} (le myocarde) gros comme le poing, divisé en \textbf{quatre cavités} : deux oreillettes et deux ventricules. Les \textbf{valvules} imposent au sang un sens unique de circulation. L'épaisseur des parois est adaptée au travail fourni : le ventricule gauche, le plus épais, propulse le sang dans tout le corps. Le cœur est une \textbf{double pompe} : le cœur droit assure la petite circulation (vers les poumons), le cœur gauche la grande circulation (vers les organes). Il bat de façon automatique, 60 à 80 fois par minute au repos, et accélère à l'effort.
\end{aretenir}

\coursec{Les hémorragies et les moyens de lutte}

Nous venons de voir que le cœur est une pompe qui propulse le sang dans un réseau de vaisseaux fermés : le sang circule en permanence, sous pression, sans jamais quitter ce circuit. Mais que se passe-t-il lorsqu'un vaisseau est rompu, par exemple lors d'une coupure profonde ou d'un accident de moto-taxi ? Le sang s'écoule alors hors des vaisseaux : c'est une \cle{hémorragie}. Comprendre ce phénomène et savoir réagir vite peut sauver une vie : c'est l'objet de cette section, qui constitue aussi la première partie du TP n°8.

\courssub{Qu'est-ce qu'une hémorragie ?}

\begin{definition}[Hémorragie]
Une \cle{hémorragie} est un écoulement de sang hors des vaisseaux sanguins, à la suite de la rupture de la paroi d'un vaisseau (artère, veine ou capillaire).
\begin{itemize}
\item On parle d'\cle{hémorragie externe} lorsque le sang s'écoule à l'extérieur du corps par une plaie : elle est \textbf{visible}.
\item On parle d'\cle{hémorragie interne} lorsque le sang s'écoule à l'intérieur du corps (dans un organe, une cavité) : elle est \textbf{non visible} et donc plus difficile à détecter.
\end{itemize}
\end{definition}

\begin{attention}[Le danger caché de l'hémorragie interne]
Une hémorragie interne ne se voit pas, mais elle est tout aussi grave qu'une hémorragie externe : le sang perdu ne circule plus dans les vaisseaux et n'irrigue plus les organes. Après un choc violent (accident, chute), une victime apparemment indemne peut être en train de saigner à l'intérieur. Il faut donc toujours surveiller une personne qui a reçu un coup violent, même sans plaie visible.
\end{attention}

\courssub{Reconnaître une hémorragie externe selon le vaisseau touché}

\begin{experience}[Observation : trois plaies, trois aspects du sang]
Trois blessés présentent chacun une plaie qui saigne. On observe attentivement l'aspect du sang et son mode d'écoulement :
\begin{itemize}
\item \textbf{Blessé A} : le sang est \textbf{rouge vif} et jaillit \textbf{par saccades}, en jets réguliers qui suivent le rythme des battements du cœur.
\item \textbf{Blessé B} : le sang est \textbf{rouge sombre} et s'écoule de façon \textbf{continue}, sans jaillissement.
\item \textbf{Blessé C} : le sang \textbf{suinte} lentement à la surface de la plaie, en petites gouttes.
\end{itemize}
\textbf{Interprétation} : rappelons-nous la structure des vaisseaux étudiée précédemment. Les \textbf{artères} reçoivent le sang directement chassé par le cœur : le sang y est sous \textbf{forte pression}, d'où le jaillissement par saccades, synchronisé avec les contractions cardiaques ; il est rouge vif car il vient d'être oxygéné (sauf artère pulmonaire). Les \textbf{veines} transportent le sang sous \textbf{faible pression} : l'écoulement est continu ; le sang est rouge sombre car il est pauvre en dioxygène (sauf veines pulmonaires). Les \textbf{capillaires} sont des vaisseaux microscopiques : leur rupture ne produit qu'un suintement.
\textbf{Conclusion} : l'aspect de l'écoulement permet d'identifier le type de vaisseau touché, donc d'évaluer la gravité.
\end{experience}

\begin{aretenir}[Les trois types d'hémorragies externes]
\begin{center}
\begin{tabularx}{\linewidth}{|l|X|X|X|}
\hline
\rowcolor{popblueL} \textbf{Critère} & \textbf{Artérielle} & \textbf{Veineuse} & \textbf{Capillaire} \\
\hline
Couleur du sang & Rouge vif & Rouge sombre & Rouge \\
\hline
Écoulement & Jaillit par saccades, au rythme du cœur & Continu, régulier & Suintement lent \\
\hline
Gravité & Très grave (perte rapide) & Grave si grosse veine & Bénigne en général \\
\hline
\end{tabularx}
\end{center}
L'hémorragie \textbf{artérielle} est la plus dangereuse : le sang étant sous forte pression, la perte est très rapide.
\end{aretenir}

\begin{popfigure}
\begin{tikzpicture}[scale=0.95]
% --- Artérielle ---
\begin{scope}
\draw[thick, poppink, fill=poppinkL] (0,0) rectangle (1.2,3);
\draw[thick, poppink] (1.2,1.5) -- (1.7,1.5);
\foreach \y in {2.1,1.5,0.9}{
  \draw[-{Stealth[length=2.5mm]}, very thick, poppink] (1.7,\y) .. controls (2.4,\y+0.35) and (3.0,\y+0.5) .. (3.6,\y+0.1);
}
\node[align=center, font=\small\bfseries, popdark] at (1.8,-0.5) {Hémorragie\\artérielle};
\node[align=center, font=\scriptsize, popmuted] at (1.8,-1.25) {sang rouge vif,\\jets par saccades};
\end{scope}
% --- Veineuse ---
\begin{scope}[xshift=5.2cm]
\draw[thick, poppurple, fill=poppurpleL] (0,0) rectangle (1.2,3);
\draw[thick, poppurple] (1.2,1.5) -- (1.7,1.5);
\draw[-{Stealth[length=2.5mm]}, very thick, poppurple] (1.7,1.5) .. controls (2.5,1.4) and (3.1,1.2) .. (3.6,0.6);
\draw[-{Stealth[length=2.5mm]}, very thick, poppurple] (1.7,1.35) .. controls (2.5,1.25) and (3.1,1.0) .. (3.6,0.45);
\node[align=center, font=\small\bfseries, popdark] at (1.8,-0.5) {Hémorragie\\veineuse};
\node[align=center, font=\scriptsize, popmuted] at (1.8,-1.25) {sang rouge sombre,\\écoulement continu};
\end{scope}
% --- Capillaire ---
\begin{scope}[xshift=10.4cm]
\draw[thick, poporange, fill=poporangeL] (0,0) rectangle (1.2,3);
\foreach \x/\y in {1.35/1.7, 1.5/1.3, 1.4/0.9, 1.6/0.5}{
  \fill[poporange] (\x,\y) circle (0.09);
}
\node[align=center, font=\small\bfseries, popdark] at (1.8,-0.5) {Hémorragie\\capillaire};
\node[align=center, font=\scriptsize, popmuted] at (1.8,-1.25) {suintement\\lent};
\end{scope}
\end{tikzpicture}
\legende{Les trois types d'hémorragies externes. Le rectangle représente la paroi du vaisseau sectionné. En haut, l'hémorragie artérielle (sang sous forte pression : jets par saccades) ; au milieu, l'hémorragie veineuse (écoulement continu) ; en bas, l'hémorragie capillaire (simple suintement).}
\end{popfigure}

\courssub{Pourquoi une hémorragie est-elle grave ?}

Le sang est irremplaçable : il transporte le dioxygène et les nutriments vers toutes les cellules. Si le volume de sang diminue trop, le cœur n'a plus assez de sang à propulser, la pression s'effondre et les organes (cerveau, cœur, reins) ne sont plus irrigués : c'est le \cle{choc}.

\begin{exemplebox}[Un exemple chiffré]
Un adulte possède environ \SI{5}{L} de sang. La perte de plus de \SI{1,5}{L} en peu de temps (soit environ 30 \% du volume total) engage le pronostic vital : la vie de la victime est en danger. Chez un enfant, dont le volume sanguin est plus faible, une perte bien moindre suffit à provoquer un état grave.
\end{exemplebox}

Les signes qui doivent alerter devant une perte de sang importante sont :
\begin{itemize}
\item une \textbf{pâleur} du visage et des lèvres ;
\item un \textbf{pouls rapide et faible} (le cœur bat plus vite pour compenser la baisse de volume) ;
\item une \textbf{soif intense}, des sueurs froides, une agitation ou une somnolence.
\end{itemize}

\courssub{TP n°8 (1\iere{} partie) : les gestes de premiers secours face à une hémorragie externe}

Face à une hémorragie externe, chaque minute compte. Le témoin n'a pas besoin d'être médecin : quatre gestes simples, dans le bon ordre, suffisent souvent à sauver la victime.

\begin{definition}[Compression, point de compression, garrot]
\begin{itemize}
\item La \cle{compression} est l'action d'appuyer fortement sur une plaie pour arrêter l'écoulement du sang.
\item Un \cle{point de compression} artérielle est un endroit du corps où l'on peut écraser une artère contre un os pour stopper l'arrivée du sang vers la plaie (exemples : artère brachiale (humérale) sur le bras, artère fémorale au pli de l'aine).
\item Le \cle{garrot} est un lien serré autour d'un membre, en amont de la plaie, pour interrompre la circulation. C'est un \textbf{dernier recours}, réservé aux cas extrêmes.
\end{itemize}
\end{definition}

\begin{methode}[Les 4 étapes du secours en cas d'hémorragie externe]
\begin{enumerate}
\item \textbf{Protéger} : se protéger d'abord soi-même (gants, ou à défaut un tissu, un sac plastique) pour éviter tout contact direct avec le sang de la victime, qui peut transmettre des maladies (hépatites, VIH).
\item \textbf{Comprimer} : appuyer fortement sur la plaie avec un linge propre (ou à main nue protégée) et maintenir la pression sans relâcher. Allonger la victime et, si possible, surélever le membre blessé. Si la compression directe ne suffit pas, comprimer le \textbf{point de compression} artérielle correspondant (artère brachiale pour le bras, artère fémorale pour la cuisse).
\item \textbf{Alerter} : appeler ou faire appeler les secours le plus vite possible (au Cameroun : le \textbf{119} ou le \textbf{112}).
\item \textbf{Surveiller} : rester auprès de la victime, surveiller sa respiration et sa conscience, ne jamais desserrer le pansement compressif, et attendre les secours.
\end{enumerate}
\end{methode}

\begin{popfigure}
\begin{tikzpicture}[scale=1.0]
% Corps stylisé
\draw[thick, popdark, fill=popcream] (0,4.6) circle (0.55); % tête
\draw[thick, popdark, fill=popcream, rounded corners=6pt] (-0.75,1.6) rectangle (0.75,4.0); % tronc
\draw[thick, popdark, fill=popcream, rounded corners=4pt] (-1.7,2.6) rectangle (-0.8,3.9); % bras G
\draw[thick, popdark, fill=popcream, rounded corners=4pt] (0.8,2.6) rectangle (1.7,3.9); % bras D
\draw[thick, popdark, fill=popcream, rounded corners=4pt] (-0.7,0) rectangle (-0.1,1.6); % jambe G
\draw[thick, popdark, fill=popcream, rounded corners=4pt] (0.1,0) rectangle (0.7,1.6); % jambe D
% Points de compression
\fill[poppink] (-1.25,3.6) circle (0.14);
\draw[-{Stealth}, thick, poppink] (-1.25,3.6) -- (-3.1,4.4);
\node[anchor=east, align=right, font=\small, popdark] at (-3.2,4.4) {\textbf{Artère brachiale}\\face interne du bras\\(hémorragie de l'avant-bras)};
\fill[poppink] (0.45,1.55) circle (0.14);
\draw[-{Stealth}, thick, poppink] (0.45,1.55) -- (2.6,2.3);
\node[anchor=west, align=left, font=\small, popdark] at (2.7,2.3) {\textbf{Artère fémorale}\\pli de l'aine\\(hémorragie de la jambe)};
\end{tikzpicture}
\legende{Points de compression artérielle principaux. On écrase l'artère contre l'os, entre le cœur et la plaie : l'artère brachiale (face interne du bras) pour une hémorragie de l'avant-bras ou de la main ; l'artère fémorale (au pli de l'aine) pour une hémorragie de la jambe.}
\end{popfigure}

\begin{attention}[Le garrot : un dernier recours dangereux]
Le garrot ne doit être posé \textbf{qu'en dernier recours}, lorsque la compression directe et les points de compression ont échoué et que la victime risque de se vider de son sang. Un garrot mal posé (trop serré, maintenu trop longtemps, ou posé sans formation) peut détruire les tissus du membre. Si l'on doit y recourir, il faut \textbf{noter l'heure de pose} et le signaler aux secours. En pratique, un pansement compressif bien réalisé suffit dans l'immense majorité des cas.
\end{attention}

\begin{attention}[Ce qu'il ne faut JAMAIS faire]
\begin{itemize}
\item \textbf{Ne jamais retirer un objet planté dans la plaie} (couteau, éclat de verre, tige de fer) : il fait « bouchon » et limite le saignement. Le retirer provoquerait une hémorragie massive. Il faut le laisser en place et comprimer \emph{autour} de l'objet.
\item \textbf{Ne jamais donner à boire} à une victime grave, même si elle réclame : elle pourrait avoir besoin d'une opération chirurgicale, qui exige l'estomac vide.
\item \textbf{Ne jamais desserrer le pansement compressif} « pour voir si ça saigne encore » : cela relancerait l'hémorragie. Si le sang traverse, on ajoute une épaisseur par-dessus et on comprime plus fort.
\end{itemize}
\end{attention}

\begin{popfigure}
\begin{tikzpicture}[
  font=\small,
  decision/.style={diamond, draw=popdark, fill=popgoldL, aspect=2.2, align=center, inner sep=1pt},
  action/.style={rectangle, rounded corners=4pt, draw=popdark, fill=popblueL, align=center, inner sep=4pt},
  danger/.style={rectangle, rounded corners=4pt, draw=popdark, fill=poppinkL, align=center, inner sep=4pt},
  fleche/.style={-{Stealth[length=2.5mm]}, thick, popdark}
]
\node[action, align=center] (start) at (0,0){Plaie qui saigne\\abondamment};
\node[action, align=center] (prot) at (0,-1.6){1. Se protéger\\(gants, tissu)};
\node[action, align=center] (comp) at (0,-3.2){2. Comprimer fort\\avec un linge propre};
\node[decision, align=center] (test) at (0,-5.3){Le saignement\\s'arrête-t-il ?};
\node[action, align=center] (point) at (4.6,-5.3){Comprimer le point\\de compression\\(brachiale ou fémorale)};
\node[action, align=center] (alerte) at (0,-7.3){3. Alerter les secours\\(119 ou 112)};
\node[action, align=center] (surv) at (0,-9.0){4. Allonger, surélever le membre,\\surveiller jusqu'aux secours};
\node[danger, align=center] (garrot) at (4.6,-7.3){Garrot : dernier recours\\noter l'heure de pose};
\draw[fleche] (start) -- (prot);
\draw[fleche] (prot) -- (comp);
\draw[fleche] (comp) -- (test);
\draw[fleche] (test) -- node[above, font=\scriptsize]{non} (point);
\draw[fleche] (test) -- node[right, font=\scriptsize]{oui} (alerte);
\draw[fleche] (point) -- (alerte);
\draw[fleche] (point) -- (garrot);
\draw[fleche] (alerte) -- (surv);
\end{tikzpicture}
\legende{Organigramme de la conduite à tenir face à une hémorragie externe : protéger, comprimer, alerter, surveiller. Le garrot n'intervient qu'en dernier recours, si toutes les compressions ont échoué.}
\end{popfigure}

\courssub{Les hémorragies internes : signes d'alerte et conduite à tenir}

Après un choc violent (accident de la route, chute, coup à l'abdomen), une hémorragie interne peut se produire sans aucune plaie visible. Il faut alors repérer les \textbf{signes d'alerte} :
\begin{itemize}
\item une \textbf{douleur} ou une contracture au niveau de l'abdomen ou du thorax ;
\item une \textbf{pâleur} importante, des sueurs froides ;
\item une \textbf{soif intense} ;
\item un \textbf{pouls rapide et faible} ;
\item parfois des vomissements de sang ou du sang dans les urines.
\end{itemize}

\begin{aretenir}[Conduite à tenir devant une hémorragie interne suspectée]
\begin{enumerate}
\item \textbf{Allonger} la victime au calme, à l'ombre, et la couvrir.
\item \textbf{Ne rien donner à boire ni à manger}, même si la victime réclame.
\item \textbf{Alerter immédiatement} les secours (119 ou 112) et \textbf{transporter d'urgence} la victime à l'hôpital : seule une équipe médicale peut arrêter une hémorragie interne.
\end{enumerate}
\end{aretenir}

\begin{savaistu}[Les accidents de la route au Cameroun]
Au Cameroun, les accidents de la route — en particulier ceux qui impliquent les \textbf{moto-taxis} — sont une cause majeure d'hémorragies graves chez les jeunes. Or, dans beaucoup de cas, la victime décède avant l'arrivée des secours alors qu'une simple compression aurait pu la sauver. Les premiers gestes des témoins sauvent des vies : savoir comprimer une plaie, c'est être capable de garder quelqu'un en vie le temps que l'ambulance arrive. C'est pourquoi les gestes de premiers secours font partie de ton programme de SVTEEHB.
\end{savaistu}

\begin{exoresolu}[Une blessure au marché]
Au marché, un jeune homme se coupe profondément l'avant-bras avec une machette. Le sang est rouge vif et jaillit par saccades. Un témoin veut retirer la lame restée plantée dans la plaie et donner de l'eau au blessé qui dit avoir soif.
\begin{enumerate}
\item Quel type d'hémorragie observe-t-on ? Justifier.
\item Relever les deux erreurs que le témoin s'apprête à commettre et les justifier.
\item Décrire la conduite à tenir correcte.
\end{enumerate}
\tcblower
\textbf{Solution détaillée.}
\begin{enumerate}
\item Il s'agit d'une \textbf{hémorragie artérielle} : le sang est rouge vif et jaillit par saccades, au rythme des battements du cœur, car le sang circule sous forte pression dans les artères.
\item Première erreur : vouloir \textbf{retirer la lame} plantée dans la plaie. L'objet enfoncé fait « bouchon » et limite le saignement ; le retirer provoquerait une hémorragie massive. Deuxième erreur : vouloir \textbf{donner de l'eau} au blessé. Une victime grave ne doit rien boire, car elle pourrait nécessiter une intervention chirurgicale à l'hôpital.
\item Conduite à tenir : (1) se \textbf{protéger} avec des gants ou un tissu ; (2) \textbf{comprimer} fortement la plaie avec un linge propre, \emph{autour} de la lame sans la toucher ; (3) \textbf{alerter} les secours (119 ou 112) ; (4) allonger la victime, surélever le membre et la \textbf{surveiller} jusqu'à l'arrivée des secours, sans jamais desserrer le pansement.
\end{enumerate}
\end{exoresolu}

\begin{exercice}[À toi de jouer]
Lors d'un match de football, un joueur heurte violemment le poteau de but avec le flanc. Il ne présente aucune plaie, mais il se plaint d'une douleur au ventre, il est très pâle, il a soif et son pouls est rapide et faible.
\begin{enumerate}
\item Que soupçonnes-tu ? Justifie ta réponse.
\item Ses coéquipiers veulent lui donner de l'eau et le faire marcher jusqu'au dispensaire. Qu'en penses-tu ?
\item Rédige la conduite à tenir complète.
\end{enumerate}
\end{exercice}

\begin{corrige}[Exercice : À toi de jouer]
\begin{enumerate}
\item On soupçonne une \textbf{hémorragie interne} : le choc violent sur le flanc a pu rompre un vaisseau ou un organe à l'intérieur de l'abdomen. Les signes d'alerte sont tous présents : douleur abdominale, pâleur, soif intense, pouls rapide et faible, sans plaie visible.
\item C'est une double erreur : il ne faut \textbf{jamais donner à boire} à une victime grave (risque d'intervention chirurgicale à venir) et il ne faut \textbf{pas la faire marcher}, car l'effort aggraverait le saignement interne.
\item Conduite à tenir : allonger la victime au calme, à l'ombre, la couvrir ; ne rien lui donner à boire ni à manger ; alerter immédiatement les secours (119 ou 112) et organiser son transport d'urgence à l'hôpital, allongée, en surveillant sa conscience et sa respiration.
\end{enumerate}
\end{corrige}

En résumé, une hémorragie est un écoulement de sang hors des vaisseaux, externe ou interne. Face à une hémorragie externe, la règle d'or tient en quatre mots : \textbf{protéger, comprimer, alerter, surveiller}. Ces gestes simples, à la portée de tous, sauvent des vies chaque jour. Nous verrons dans la section suivante que le sang peut aussi être victime d'accidents \emph{à l'intérieur} des vaisseaux eux-mêmes : ce sont les maladies cardiovasculaires et les accidents vasculaires cérébraux.

\coursec{Maladies cardiovasculaires : causes et prévention}

Après avoir étudié les hémorragies, accidents brutaux où le sang s'échappe du réseau vasculaire, nous nous penchons maintenant sur des pathologies qui s'installent dans la durée : les \cle{maladies cardiovasculaires}. Elles ne surgissent pas par hasard ; elles sont la conséquence d'une altération progressive des vaisseaux et du cœur, souvent silencieuse pendant des années. Comprendre leurs mécanismes, c'est acquérir les clés pour les prévenir efficacement.

\courssub{Les maladies cardiovasculaires : un enjeu majeur de santé publique}

\begin{definition}[Maladies cardiovasculaires (MCV)]
Les \textbf{maladies cardiovasculaires} regroupent l'ensemble des affections qui touchent le \textbf{cœur} (cardiopathies) et les \textbf{vaisseaux sanguins} (vasculopathies). Elles constituent la \textbf{première cause de mortalité dans le monde} (environ 17,9 millions de décès par an selon l'OMS, soit 32 \% de l'ensemble des décès mondiaux). Au Cameroun, elles sont en nette augmentation du fait de l'urbanisation et des changements de mode de vie, touchant des adultes de plus en plus jeunes.
\end{definition}

\begin{savaistu}[L'OMS et la charge des MCV]
L'Organisation Mondiale de la Santé estime que \textbf{la plupart des maladies cardiovasculaires pourraient être évitées} en agissant sur les facteurs de risque comportementaux : tabagisme, alimentation déséquilibrée, inactivité physique, consommation nocive d'alcool. La détection précoce et la prise en charge de l'hypertension artérielle, du diabète et de l'hypercholestérolémie permettent de réduire drastiquement le risque d'accidents graves (infarctus, AVC).
\end{savaistu}

\courssub{L'athérosclérose : le processus silencieux qui obstrue les artères}

La plupart des maladies cardiovasculaires (infarctus du myocarde, AVC ischémique, artérite des membres inférieurs) partagent une cause commune : l'\textbf{athérosclérose}. C'est une maladie chronique de la paroi des \textbf{artères} de gros et moyen calibre.

\begin{experience}[Observation de coupes d'artères (documentaire)]
\textbf{Matériel :} Photographies au microscope optique de coupes transversales d'artères humaines (artère coronaire, aorte) : une artère saine et une artère athéromateuse. Coloration Movat pentachrome (élastine noire, collagène jaune/vert, lipides/athérome rouge vif).
\textbf{Observations :}
\begin{itemize}
    \item \textbf{Artère saine :} La lumière (lumière vasculaire) est large et régulière. L'intima (couche interne) est fine, l'endothélium est intact. La média (couche musculaire) est bien organisée. Le sang circule librement (flux laminaire).
    \item \textbf{Artère athéromateuse :} La lumière est \textbf{rétrécie} (sténose) de façon excentrique. L'intima est \textbf{épaissie} par une \textbf{plaque d'athérome} (ou plaque de stéatose) : cœur lipidique (graisses, cholestérol) recouvert d'une chape fibreuse (collagène, cellules musculaires lisses). La média peut être atrophiée. Des signes d'inflammation (infiltrat lymphocytaire) et de calcification sont visibles.
\end{itemize}
\textbf{Interprétation :} Le dépôt de \textbf{cholestérol} (transporté par les LDL, « mauvais cholestérol ») sous l'endothélium déclenche une réaction inflammatoire chronique. Les macrophages engloutissent les lipides et deviennent des \textbf{cellules spumeuses}. Elles meurent, libérant leur contenu et formant le cœur nécrotique de la plaque. Une chape fibreuse tente de stabiliser l'ensemble. La plaque grandit vers la lumière (obstruction) et vers l'extérieur (remodelage).
\textbf{Conclusion :} L'athérosclérose est une \textbf{maladie inflammatoire chronique} de l'artère, initiée par la rétention de lipides dans l'intima, aboutissant à la formation de plaques qui rétrécissent la lumière artérielle et rigidifient la paroi.
\end{experience}

\begin{definition}[Athérosclérose]
L'\textbf{athérosclérose} est une affection artérielle caractérisée par la formation de \textbf{plaques d'athérome} (dépôts de lipides, principalement du cholestérol, de cellules inflammatoires, de tissu fibreux et de calcium) au niveau de l'intima, entraînant un \textbf{épaississement de la paroi}, une \textbf{perte d'élasticité} et un \textbf{rétrécissement de la lumière} (sténose) qui gêne la circulation sanguine.
\end{definition}

\begin{popfigure}
\begin{tikzpicture}[scale=0.9, every node/.style={font=\small}]
% Couleurs
\colorlet{arterywall}{popblue!20}
\colorlet{endothelium}{popblue}
\colorlet{blood}{poppink}
\colorlet{plaque}{poporange}
\colorlet{fibrous}{popgold}
\colorlet{lipidcore}{poporange!80!black}
\colorlet{muscle}{poppurple!30}
% --- ARTÈRE SAINE ---
\begin{scope}[xshift=0cm]
% Paroi externe (adventice)
\fill[arterywall] (-2.5,-2) rectangle (2.5,2);
% Média (musculaire)
\fill[muscle] (-2,-1.5) rectangle (2,1.5);
% Intima / Endothélium
\fill[endothelium] (-2,-0.2) rectangle (2,0.2);
% Lumière
\fill[blood!30] (-1.8,-0.18) rectangle (1.8,0.18);
% Flux sanguin (flèches)
\foreach \y in {-0.1, 0, 0.1} {
    \draw[->, thick, blood] (-2.2,\y) -- (-1.0,\y);
    \draw[->, thick, blood] (1.0,\y) -- (2.2,\y);
}
% Légendes
\node[align=center, text width=4.5cm] at (0, -2.7) {\textbf{Artère saine}\\\textit{Coupe transversale}};
\draw[<-, popblue] (-1.9,0.5) -- (-2.8,1.2) node[left, align=left] {Endothélium\\(intact)};
\draw[<-, poppurple] (-1.9,1.2) -- (-2.8,1.8) node[left] {Média (muscle lisse)};
\draw[<-, popmuted] (-1.9,-1.2) -- (-2.8,-1.8) node[left] {Adventice};
\draw[<-, blood] (0,0.5) -- (0,1.2) node[above, align=center]{Lumière large\\Flux laminaire};
\end{scope}
% --- ARTÈRE ATHÉROMATEUSE ---
\begin{scope}[xshift=7cm]
% Paroi externe
\fill[arterywall] (-2.5,-2) rectangle (2.5,2);
% Média (amincie par la plaque)
\fill[muscle] (-2,-1.5) rectangle (-0.5,1.5); % Gauche
\fill[muscle] (1.5,-1.5) rectangle (2,1.5);   % Droite (fine)
% Plaque d'athérome
% Chape fibreuse
\fill[fibrous] (-0.5,-0.3) .. controls (-0.2,0.5) and (1.2,0.5) .. (1.5,0.3) -- (1.5,-0.3) -- cycle;
% Cœur lipidique
\fill[lipidcore] (-0.5,-0.25) .. controls (-0.2,0.3) and (1.0,0.3) .. (1.3,0.0) -- (1.3,-0.25) -- cycle;
% Calcifications (points)
\foreach \x/\y in {-0.3/0.0, 0.2/0.1, 0.8/-0.05, 1.1/0.15} {
    \fill[popdark] (\x,\y) circle (0.05);
}
% Lumière rétrécie
\fill[blood!30] (-0.4,-0.28) .. controls (-0.1,0.25) and (1.1,0.25) .. (1.4,-0.02) -- (1.4,-0.28) -- cycle;
% Flux turbulent
\draw[->, thick, blood, decorate, decoration={snake, amplitude=0.5pt, segment length=4pt}] (-2.2,0.1) -- (-0.6,0.1);
\draw[->, thick, blood, decorate, decoration={snake, amplitude=0.5pt, segment length=4pt}] (1.6,0.0) -- (2.2,0.0);
\draw[->, thick, blood, decorate, decoration={snake, amplitude=0.5pt, segment length=4pt}] (-2.2,-0.1) -- (-0.6,-0.1);
\draw[->, thick, blood, decorate, decoration={snake, amplitude=0.5pt, segment length=4pt}] (1.6,-0.1) -- (2.2,-0.1);
% Légendes
\node[align=center, text width=4.5cm] at (0, -2.7) {\textbf{Artère athéromateuse}\\\textit{Coupe transversale}};
\draw[<-, poporange] (0.5,0.4) -- (1.5,1.2) node[right, align=left] {Plaque d'athérome\\(Chape fibreuse + Cœur lipidique)};
\draw[<-, popdark] (0.8,0.0) -- (1.8,-0.5) node[right] {Calcifications};
\draw[<-, poppurple] (-1.9,1.2) -- (-2.8,1.8) node[left] {Média amincie};
\draw[<-, blood] (0,-0.7) -- (0,-1.3) node[below, align=center]{Lumière rétrécie\\(Sténose) - Flux turbulent};
\end{scope}
\end{tikzpicture}
\legende{Comparaison schématique en coupe transversale : à gauche, une artère saine à lumière large et paroi fine ; à droite, une artère atteinte d'athérosclérose avec une plaque d'athérome (chape fibreuse jaune, cœur lipidique orange) rétrécissant la lumière (sténose) et perturbant le flux sanguin (turbulences).}
\end{popfigure}

\begin{aretenir}[Conséquences hémodynamiques de la plaque d'athérome]
\begin{enumerate}
    \item \textbf{Sténose :} Rétrécissement de la lumière $\rightarrow$ augmentation de la résistance vasculaire $\rightarrow$ ischémie (manque d'oxygène) en aval à l'effort (angine de poitrine, claudication).
    \item \textbf{Rigidité artérielle :} Perte de l'élasticité (effet Windkessel diminué) $\rightarrow$ augmentation de la pression artérielle systolique (hypertension systolique isolée du sujet âgé).
    \item \textbf{Rupture de plaque :} La chape fibreuse peut se rompre (plaque instable) $\rightarrow$ exposition du cœur lipidique thrombogène au sang $\rightarrow$ formation d'un \textbf{thrombus} (caillot) $\rightarrow$ \textbf{obstruction brutale totale} (infarctus, AVC ischémique).
\end{enumerate}
\end{aretenir}

\courssub{Les principales maladies cardiovasculaires}

L'athérosclérose, selon sa localisation et sa complication, entraîne des tableaux cliniques distincts.

\begin{definition}[Hypertension artérielle (HTA)]
L'\textbf{hypertension artérielle} est une élévation \textbf{constante et durable} de la pression artérielle au repos. On parle d'HTA lorsque la \textbf{pression artérielle systolique (PAS, le « maximum »)} est $\ge 140$ mmHg ($14$ kPa) \textbf{ET/OU} la \textbf{pression artérielle diastolique (PAD, le « minimum »)} est $\ge 90$ mmHg ($9$ kPa), confirmée sur plusieurs mesures lors de consultations distinctes. C'est le principal facteur de risque cardiovasculaire modifiable.
\end{definition}

\begin{exemplebox}[Repères tensionnels (adulte) - Valeurs moyennes au brassard]
\begin{center}
\begin{tabularx}{\linewidth}{|l|X|X|}\hline
\textbf{Catégorie} & \textbf{Systolique (mmHg)} & \textbf{Diastolique (mmHg)} \\ \hline
Optimale & $< 120$ & $< 80$ \\
Normale & $120 - 129$ & $80 - 84$ \\
Normale haute & $130 - 139$ & $85 - 89$ \\
\rowcolor{poporangeL} \textbf{HTA Grade 1 (modérée)} & \textbf{$140 - 159$} & \textbf{$90 - 99$} \\
\rowcolor{poporangeL} \textbf{HTA Grade 2 (sévère)} & \textbf{$160 - 179$} & \textbf{$100 - 109$} \\
\rowcolor{poppinkL} \textbf{HTA Grade 3 (très sévère)} & \textbf{$\ge 180$} & \textbf{$\ge 110$} \\ \hline
\end{tabularx}
\end{center}
\textbf{Exemple :} Une tension de $12/8$ (soit $120/80$ mmHg) est normale. Une tension mesurée à $14,5/9,2$ ($145/92$ mmHg) à deux reprises signe une HTA Grade 1.
\end{exemplebox}

\begin{attention}[Le piège du « tueur silencieux »]
\textbf{Erreur fréquente :} « Je n'ai pas mal à la tête, donc je n'ai pas d'hypertension. »
\textbf{Réalité :} Dans $> 90\%$ des cas, l'HTA est \textbf{asymptomatique} pendant des années. Les maux de tête, vertiges, bourdonnements d'oreilles ne sont \textbf{PAS} des signes fiables ni constants. Seule la \textbf{mesure répétée de la tension artérielle} permet le diagnostic. Ne pas attendre les symptômes pour se faire dépister !
\end{attention}

\begin{definition}[Infarctus du myocarde (IDM ou crise cardiaque)]
L'\textbf{infarctus du myocarde} est la \textbf{nécrose (mort) irréversible} d'une partie du muscle cardiaque (myocarde) due à l'\textbf{obstruction brutale et complète} d'une \textbf{artère coronaire} (le plus souvent par un thrombus sur une plaque d'athérome rompue). La zone en aval de l'obstruction n'est plus perfusée $\rightarrow$ anoxie $\rightarrow$ mort cellulaire en quelques heures si la circulation n'est pas rétablie en urgence.
\end{definition}

\begin{popfigure}
\begin{tikzpicture}[scale=1.1, every node/.style={font=\small}]
% Cœur simplifié (vue antérieure)
% Oreillettes
\fill[poppink!30] (-2.5,1.5) .. controls (-2.5,2.5) and (-0.5,2.5) .. (-0.5,1.5) -- cycle; % GA
\fill[poppink!30] (0.5,1.5) .. controls (0.5,2.5) and (2.5,2.5) .. (2.5,1.5) -- cycle; % DA
% Ventricules
\fill[poppink!50] (-2.5,1.5) .. controls (-3.0,0) and (-3.0,-2.5) .. (0,-3.0) .. controls (0.5,-2.5) and (0.5,0) .. (-0.5,1.5) -- cycle; % VG
\fill[poppink!50] (0.5,1.5) .. controls (0.5,0) and (0.5,-2.5) .. (0,-3.0) .. controls (-0.5,-2.5) and (-0.5,0) .. (2.5,1.5) -- cycle; % VD (simplifié)
% Artères coronaires
% Coronaire gauche (tronc commun -> interventriculaire antérieure + circonflexe)
\draw[very thick, popblue] (0,1.5) -- (-0.2,1.8) .. controls (-1.0,2.2) and (-1.5,1.5) .. (-1.8,0.5); % Circonflexe
\draw[very thick, popblue] (0,1.5) -- (0.2,1.2) .. controls (0.5,0.5) and (0.5,-1.0) .. (0.3,-2.0); % Interventriculaire antérieure (IVA)
% Coronaire droite
\draw[very thick, popblue] (0,1.5) -- (0.3,1.8) .. controls (1.2,2.2) and (2.0,1.5) .. (2.2,0.2) .. controls (2.0,-0.5) and (1.5,-1.5) .. (0.5,-2.5); % Droit + postérieure
% Zone d'infarctus (antéro-septale, territoire IVA)
\fill[poporange!70!black, opacity=0.7] (0.0,-0.5) .. controls (0.8,-0.2) and (1.0,-1.5) .. (0.2,-2.2) .. controls (-0.5,-2.0) and (-0.5,-0.8) .. cycle;
% Thrombus sur IVA
\fill[popdark] (0.3,0.2) circle (0.15);
\node[popdark, font=\tiny\bfseries] at (0.3,0.2) {×};
% Flèches légende
\draw[<-, thick, popblue] (-2.0,2.0) -- (-3.0,2.5) node[left, align=left] {Artère coronaire\\gauche (Circonflexe)};
\draw[<-, thick, popblue] (0.3,-2.2) -- (0.5,-3.0) node[right, align=right] {Artère interventriculaire\\antérieure (IVA)};
\draw[<-, thick, popblue] (2.5,1.0) -- (3.5,1.5) node[right, align=left] {Artère coronaire\\droite};
\draw[<-, thick, poporange!70!black] (0.5,-1.2) -- (1.5,-1.5) node[right, align=left] {Zone nécrosée\\(Infarctus antéro-septal)};
\draw[<-, thick, popdark] (0.3,0.5) -- (0.5,1.0) node[right, align=left] {Thrombus\\obstructif};
\node at (0, -3.8) {\textbf{Cœur vu de face : artères coronaires et siège d'un infarctus du myocarde}};
\end{tikzpicture}
\legende{Schéma du cœur (vue antérieure) montrant les deux artères coronaires (bleu) qui vascularisent le myocarde. Un thrombus (croix noire) obstrue l'artère interventriculaire antérieure (IVA). La zone en aval (orange foncé) est nécrosée : c'est l'infarctus du myocarde.}
\end{popfigure}

\begin{definition}[Angine de poitrine (Angor) et Insuffisance cardiaque]
\begin{itemize}
    \item \textbf{Angine de poitrine (Angor) :} Douleur thoracique (serrement, brûlure) survenant à l'effort, calmée par le repos. Elle traduit une \textbf{ischémie myocardique réversible} (sténose coronaire critique sans occlusion totale). C'est un signal d'alarme.
    \item \textbf{Insuffisance cardiaque :} Incapacité du cœur à assurer un débit sanguin suffisant pour les besoins de l'organisme. Elle peut être la séquelle d'infarctus répétés, d'une HTA mal contrôlée ou d'une valvulopathie. Signes : dyspnée (essoufflement) à l'effort puis au repos, œdèmes des membres inférieurs, fatigue majeure.
\end{itemize}
\end{definition}

\courssub{Facteurs de risque : ce qui nous expose aux maladies}

On distingue les facteurs sur lesquels on ne peut pas agir et ceux sur lesquels on \textbf{peut} agir. C'est là toute la marge de manœuvre de la prévention.

\begin{propriete}[Facteurs de risque cardiovasculaire (admis)]
\begin{itemize}
    \item \textbf{Non modifiables :} \textbf{Âge} (risque $\uparrow$ après 50 ans homme, 60 ans femme), \textbf{Sexe} (homme $>$ femme avant ménopause), \textbf{Hérédité} (antécédent familial précoce : IDM père $< 55$ ans, mère $< 65$ ans).
    \item \textbf{Modifiables (comportementaux et métaboliques) :}
    \begin{enumerate}
        \item \textbf{Tabagisme} (même passif) : Endommage l'endothélium, favorise la thrombose, baisse le HDL.
        \item \textbf{Alimentation déséquilibrée} : Excès de sel (HTA), graisses saturées/trans (LDL $\uparrow$), sucres ajoutés (diabète, triglycérides), déficit en fruits/légumes (antioxydants, potassium, fibres).
        \item \textbf{Sédentarité / Inactivité physique} : $< 150$ min/semaine d'activité modérée.
        \item \textbf{Surpoids / Obésité} (IMC $\ge 25 / 30$ kg/m$^2$), surtout obésité abdominale (tour de taille $> 94$ cm H, $> 80$ cm F).
        \item \textbf{Consommation d'alcool} excessive.
        \item \textbf{Stress chronique} et manque de sommeil.
    \end{enumerate}
    \item \textbf{Métaboliques (souvent secondaires aux comportementaux) :} Hypertension artérielle, Diabète de type 2, Dyslipidémie (LDL-cholestérol $\uparrow$, HDL-cholestérol $\downarrow$, Triglycérides $\uparrow$).
\end{itemize}
\end{propriete}

\begin{propriete}[Effet cumulatif des facteurs de risque (admis)]
\textbf{Plus le nombre de facteurs de risque est élevé, plus la probabilité de survenue d'une maladie cardiovasculaire augmente de façon exponentielle (synergie).} Un sujet hypertendu, fumeur et diabétique a un risque bien supérieur à la simple addition des risques individuels.
\end{propriete}

\begin{popfigure}
\begin{tikzpicture}
\begin{axis}[
    ybar,
    bar width=12pt,
    width=\linewidth,
    height=8cm,
    enlarge x limits=0.25,
    ylabel={Incidence relative des MCV (\% sur 10 ans)},
    xlabel={Nombre de facteurs de risque modifiables présents},
    symbolic x coords={0, 1, 2, 3, 4},
    xtick=data,
    nodes near coords,
    nodes near coords align={vertical},
    ymin=0, ymax=45,
    ymajorgrids=true,
    grid style=dashed,
    tick label style={font=\small},
    label style={font=\small},
    legend style={at={(0.5,-0.2)}, anchor=north, legend columns=-1, font=\small},
    bar shift=0pt,
]
\addplot[fill=popblue!60, draw=popblue] coordinates {(0,3) (1,7) (2,15) (3,28) (4,42)};
\legend{Population générale (données fictives réalistes type Framingham/INTERHEART)}
\end{axis}
\end{tikzpicture}
\legende{Diagramme en barres montrant l'augmentation drastique de l'incidence des maladies cardiovasculaires (MCV) à 10 ans en fonction du nombre de facteurs de risque modifiables cumulés (tabac, HTA, diabète, dyslipidémie, obésité, sédentarité). Avec 0 facteur, risque faible ($\sim 3\%$) ; avec $\ge 4$ facteurs, risque multiplié par 14 ($\sim 42\%$).}
\end{popfigure}

\begin{exoresolu}[Interprétation de données épidémiologiques : facteurs de risque et fréquence des MCV]
\textbf{Énoncé :} Une étude de cohorte (type Framingham) a suivi 5 000 sujets sains pendant 10 ans. Le tableau ci-dessous rapporte l'incidence cumulative (pourcentage de sujets ayant fait un événement cardiovasculaire majeur : IDM, AVC, mort subite) selon le nombre de facteurs de risque modifiables présents à l'inclusion.

\begin{center}
\begin{tabularx}{\linewidth}{|c|X|X|}\hline
\textbf{Nombre de facteurs de risque} & \textbf{Effectif (n)} & \textbf{Incidence cumulative à 10 ans (\%)} \\ \hline
0 & 850 & 2,8 \\
1 & 1 200 & 6,5 \\
2 & 1 500 & 14,2 \\
3 & 1 000 & 27,8 \\
$\ge$ 4 & 450 & 41,5 \\ \hline
\end{tabularx}
\end{center}

\textbf{Questions :}
\begin{enumerate}
    \item Représentez graphiquement l'évolution de l'incidence en fonction du nombre de facteurs de risque (choisir le graphique adapté).
    \item Décrivez la relation entre le nombre de facteurs de risque et l'incidence des MCV.
    \item Calculez le risque relatif (RR) de faire un événement cardiovasculaire pour un sujet présentant $\ge 4$ facteurs de risque par rapport à un sujet n'en présentant aucun.
    \item En déduire l'importance de la prévention primaire.
\end{enumerate}

\tcblower

\textbf{Solution détaillée :}
\begin{enumerate}
    \item \textbf{Graphique :} Un \textbf{diagramme en barres} (ou histogramme) est le plus adapté car la variable X (nombre de facteurs) est discrète/qualitative ordonnée et la variable Y (incidence \%) est quantitative. (Voir figure ci-dessus réalisée avec pgfplots).
    \item \textbf{Description :} L'incidence des MCV augmente \textbf{régulièrement et de façon non-linéaire (exponentielle)} avec le nombre de facteurs de risque. Elle passe de $2,8\%$ (0 facteur) à $41,5\%$ ($\ge 4$ facteurs), soit une multiplication par environ 15. L'ajout de chaque facteur supplémentaire aggrave le pronostic de manière plus marquée que le précédent (synergie).
    \item \textbf{Calcul du Risque Relatif (RR) :}
    \[ RR = \frac{\text{Incidence (exposés : $\ge 4$ facteurs)}}{\text{Incidence (non exposés : 0 facteur)}} = \frac{41,5}{2,8} \approx 14,8 \]
    Un sujet cumulant $\ge 4$ facteurs de risque a un risque \textbf{presque 15 fois plus élevé} de faire un événement cardiovasculaire majeur à 10 ans qu'un sujet sans facteur de risque.
    \item \textbf{Conclusion préventive :} Comme la majorité de ces facteurs sont \textbf{modifiables} (tabac, alimentation, activité, poids, alcool), agir sur \textbf{un seul} facteur fait déjà baisser le risque de façon significative. La prévention primaire (hygiène de vie) est l'arme la plus puissante et la moins coûteuse contre l'épidémie de MCV.
\end{enumerate}
\end{exoresolu}

\courssub{Prévention et hygiène de vie : agir sur les facteurs modifiables}

La prévention des maladies cardiovasculaires repose sur la \textbf{prise en charge globale du risque vasculaire}. Il ne suffit pas de traiter un chiffre (tension, cholestérol) ; il faut modifier le terrain.

\begin{methode}[Mesurer sa tension artérielle avec un tensiomètre électronique brassard (automesure)]
\textbf{Matériel :} Tensiomètre brassard validé (norme CE), brassard adapté à la circonférence du bras.
\textbf{Conditions (règle des 3 « non » : non fumeur, non café, non effort depuis 30 min ; vessie vide ; repos 5 min assis).}
\begin{enumerate}
    \item \textbf{Installation :} Assis, dos appuyé, pieds au sol, jambes non croisées. Bras nu, soutenu sur une table, paume vers le haut, coude au niveau du cœur.
    \item \textbf{Position du brassard :} Bord inférieur à 2-3 cm au-dessus du pli du coude. Tuyau du brassard dans l'axe de l'artère brachiale (face interne du bras). Serrage : passer un doigt difficilement.
    \item \textbf{Mesure :} Appuyer sur « Start ». Ne pas parler, ne pas bouger pendant le gonflage/dégonflage.
    \item \textbf{Lecture :} L'appareil affiche deux valeurs : \textbf{SYS (Systolique = Maximum)} et \textbf{DIA (Diastolique = Minimum)}. Exemple : \textbf{128 / 82 mmHg} (se lit « 12,8 sur 8,2 »).
    \item \textbf{Répétition :} Effectuer \textbf{3 mesures consécutives} à 1-2 min d'intervalle. Noter la moyenne des 2 dernières (la 1ère est souvent plus haute).
    \item \textbf{Interprétation (automesure domicile) :} Moyenne $< 135/85$ mmHg = Normale. Moyenne $\ge 135/85$ mmHg = HTA confirmée (consulter). \textit{Note : Seuils plus bas qu'au cabinet (140/90) car effet « blouse blanche » absent.}
\end{enumerate}
\end{methode}

\begin{aretenir}[Les piliers de la prévention cardiovasculaire (Hygiène de la circulation)]
\begin{enumerate}
    \item \textbf{Alimentation de type méditerranéen / DASH :} Fruits/légumes $\ge 5$ portions/jour ; céréales complètes ; poisson (oméga-3) 2x/sem ; huiles végétales (olive, colza) ; \textbf{sel $< 5$ g/jour} (1 c. à café) ; limiter charcuteries, plats préparés, boissons sucrées, graisses saturées (beurre, fromage, viandes grasses).
    \item \textbf{Activité physique régulière :} $\ge 150$ min/sem d'endurance modérée (marche rapide, vélo, natation) OU $75$ min/sem intense + renforcement musculaire 2x/sem. Lutte contre la sédentarité (se lever toutes les heures).
    \item \textbf{Arrêt total du tabac :} Bénéfice immédiat (CO $\downarrow$ en 24h, risque IDM divisé par 2 en 1 an). Aide : substituts nicotiniques, varénicline, soutien psychologique.
    \item \textbf{Alcool :} Modération stricte ($\le 2$ verres/jour, pas tous les jours, $\le 10$ verres/sem). Zéro grossesse/conduite/médicaments interactifs.
    \item \textbf{Contrôle pondéral :} Viser IMC $18,5 - 25$ kg/m$^2$ et tour de taille $< 94$ cm (H) / $< 80$ cm (F).
    \item \textbf{Dépistage et suivi médical :} \textbf{Mesure de la TA au moins 1 fois/an} (dès 18 ans, ou plus tôt si facteurs de risque). Bilan lipidique et glycémie à jeun selon facteurs de risque. Observance rigoureuse des traitements prescrits par le médecin (anti-hypertenseurs, hypolipidémiants, antiagrégants).
\end{enumerate}
\end{aretenir}

\begin{savaistu}[Objectif « 25 by 25 » de l'OMS]
L'OMS a fixé pour objectif une réduction relative de \textbf{25 \% de la mortalité prématurée (30-70 ans) due aux maladies non transmissibles (dont MCV) d'ici 2025} (par rapport à 2010). Les cibles clés incluent : réduction de 30 \% de la prévalence du tabagisme, réduction de 30 \% de la consommation moyenne de sel, arrêt de la progression du diabète et de l'obésité, et traitement d'au moins 50 \% des personnes éligibles pour prévenir les IDM/AVC. Au Cameroun, la « Semaine du Cœur » et les campagnes de dépistage gratuit de l'HTA dans les districts de santé s'inscrivent dans cette dynamique.
\end{savaistu}

\begin{exercice}[Évaluation des connaissances : Prévention cardiovasculaire]
\textbf{M. Kamga, 52 ans, chef d'entreprise à Douala, consulte pour un bilan de santé. Il fume 15 cigarettes/jour depuis 30 ans. Il ne pratique aucune activité sportive. Son alimentation est riche en viande rouge, plats frits, sel et alcool (2 bières/jour). Ses mensurations : Taille 1,75 m, Masse 92 kg, Tour de taille 102 cm. Sa tension artérielle (moyenne de 3 mesures au cabinet) : 152/96 mmHg. Glycémie à jeun : 1,28 g/L. LDL-cholestérol : 1,85 g/L. HDL-cholestérol : 0,35 g/L.}
\begin{enumerate}
    \item Calculez l'IMC de M. Kamga. Interprétez-le.
    \item Identifiez \textbf{tous} les facteurs de risque cardiovasculaires modifiables et non modifiables présents chez ce patient.
    \item Classifiez son hypertension artérielle (Grade).
    \item Quel est son nombre total de facteurs de risque majeurs ? En vous référant au graphique de la leçon (incidence relative), estimez son risque cardiovasculaire à 10 ans.
    \item Proposez un plan de prise en charge prioritaire (hygiéno-diététique et suivi médical) en justifiant chaque mesure.
\end{enumerate}
\end{exercice}

\begin{corrige}[Exercice n° : Évaluation des connaissances]
\begin{enumerate}
    \item \textbf{IMC} = Masse / Taille$^2$ = $92 / (1{,}75)^2 = 92 / 3{,}0625 \approx 30{,}0$ kg/m$^2$. \textbf{Interprétation : Obésité (Grade 1, IMC 30-34,9)}. Tour de taille $102$ cm $> 94$ cm $\rightarrow$ \textbf{Obésité abdominale} (facteur de risque métabolique majeur).
    \item \textbf{Non modifiables :} Âge (52 ans), Sexe (Homme).
    \textbf{Modifiables :} Tabagisme actif, Sédentarité, Alimentation déséquilibrée (grasse, salée, alcool), Surpoids/Obésité abdominale, Hypertension artérielle, Dyslipidémie (LDL $\uparrow$, HDL $\downarrow$), Diabète de type 2 (Glycémie 1,28 g/L $\ge 1,26$ g/L $\rightarrow$ diabète confirmé si recontrôlé).
    \item \textbf{HTA Grade 1 (modérée)} : PAS 152 mmHg (dans $140-159$ $\rightarrow$ Grade 1) ET PAD 96 mmHg (dans $90-99$ $\rightarrow$ Grade 1). On classe au grade le plus élevé atteint par l'une ou l'autre : ici Grade 1.
    \item \textbf{Facteurs de risque majeurs comptabilisés :} Âge, Sexe, Tabac, HTA, Diabète, HDL bas, LDL haut, Obésité abdominale. Soit \textbf{au moins 7-8 facteurs majeurs}. Sur le graphique ($\ge 4$ facteurs), l'incidence relative est $\approx 42\%$ à 10 ans. Son risque réel est \textbf{très élevé} ($> 20\%$ à 10 ans).
    \item \textbf{Plan de prise en charge :}
    \begin{itemize}
        \item \textbf{Urgence :} Arrêt du tabac (proposition substituts nicotiniques + soutien psychologique).
        \item \textbf{Consultation médicale spécialisée :} Mise sous traitement antihypertenseur et hypolipidémiant (statine) adaptés, surveillance de la glycémie (régimes + traitement si nécessaire). Observance stricte des traitements prescrits.
        \item \textbf{Hygiène de vie (indispensable) :} Régime DASH/Méditerranéen (sel $< 5$g/j, stop fritures/viande rouge, augmenter légumes/poisson/huiles végétales). Marche rapide 30 min/jour minimum. Perte de poids viser $-5$ à $-10\%$ de la masse initiale. Alcool : réduction drastique ($\le 1$ verre/jour, jours sans).
        \item \textbf{Suivi :} Contrôle TA mensuel jusqu'à normalisation puis trimestriel, LDL 3 mois, HbA1c 3 mois, Fonction rénale, Poids mensuel.
    \end{itemize}
\end{enumerate}
\end{corrige}

\coursec{Les accidents vasculaires cérébraux (AVC) et les premiers secours}

Après avoir étudié les maladies cardiovasculaires qui touchent le cœur et les vaisseaux de manière chronique, nous abordons maintenant une urgence vitale soudaine : l'\cle{accident vasculaire cérébral (AVC)}. C'est la complication la plus redoutée de l'hypertension artérielle et de l'athérosclérose. Au Cameroun, comme dans de nombreux pays d'Afrique subsaharienne, l'incidence des AVC augmente et touche des sujets de plus en plus jeunes, souvent en âge de travailler. Savoir reconnaître les signes et réagir immédiatement peut sauver une vie et limiter les séquelles.

\courssub{Définition et mécanismes de l'AVC}

\begin{definition}[Accident vasculaire cérébral (AVC)]
Un \textbf{accident vasculaire cérébral (AVC)} est un accident brutal de la circulation du sang dans le cerveau, entraînant une interruption brutale de l'apport en oxygène et en nutriments (glucose) vers une zone cérébrale. Cette privation provoque la mort rapide des cellules nerveuses (neurones) de la zone touchée.
\end{definition}

On distingue deux mécanismes principaux, aux conséquences similaires mais aux prises en charge initiales différentes :

\begin{propriete}[Les deux types d'AVC]
\begin{itemize}
    \item \textbf{AVC ischémique (infarctus cérébral) :} représente environ \textbf{80 \% des cas}. Il est dû à l'\textbf{obstruction} d'une artère cérébrale par un caillot (thrombus formé sur une plaque d'athérome ou embolie venue du cœur). Le sang ne passe plus en aval de l'obstacle.
    \item \textbf{AVC hémorragique (hémorragie intra-cérébrale) :} représente environ \textbf{20 \% des cas}. Il est dû à la \textbf{rupture} d'une artère cérébrale, le plus souvent fragilisée par une \textbf{hypertension artérielle} mal contrôlée. Le sang s'épanche dans le tissu cérébral (hématome), comprimant les neurones et interrompant la circulation locale.
\end{itemize}
\end{propriete}

\begin{popfigure}
\begin{tikzpicture}[scale=0.9, every node/.style={font=\small}]
    % Cerveau gauche : AVC ischémique
    \begin{scope}[xshift=-4cm]
        % Contour cerveau
        \draw[popdark, thick] (0,0) .. controls (1,2) and (3,2.5) .. (4,1.5) .. controls (4.5,0) and (3.5,-1.5) .. (2,-2) .. controls (0.5,-2) and (-0.5,-1) .. (0,0);
        % Hémisphères
        \draw[popdark] (2,0) .. controls (2,1) and (2.5,1.5) .. (3,1);
        % Artères
        \draw[popblue, very thick] (1.5,0.5) -- (2.5,1.2) node[midway, above, sloped] {Artère cérébrale};
        % Caillot (obstruction)
        \fill[poporange] (2.2, 0.95) circle (0.15);
        \draw[poporange, thick] (2.2, 0.95) circle (0.15);
        \node[poporange, anchor=south] at (2.2, 1.2) {\textbf{Thrombe}};
        % Zone ischémique (hachurée)
        \pattern[pattern=north east lines, pattern color=poporange!50] (2.5,0.5) .. controls (3,0) and (3.5,-0.5) .. (3.5,-1) .. controls (3,-1.5) and (2.5,-1) .. cycle;
        \node[poporange, align=center] at (3.5, -0.5) {Zone ischémique\\(privée de sang)};
        % Flèche bloquée
        \draw[->, popblue, thick] (1.5,0.5) -- (2.05, 0.85);
        \draw[poporange, very thick, decorate, decoration={zigzag, segment length=2mm, amplitude=1mm}] (2.05, 0.85) -- (2.35, 1.05);
        \node[popblue, anchor=north] at (1.5, 0.2) {Flux sanguin};
        \node[popdark, align=center, font=\bfseries] at (2, -2.5) {AVC Ischémique\\(Obstruction ≈ 80\%)};
    \end{scope}

    % Cerveau droit : AVC hémorragique
    \begin{scope}[xshift=4cm]
        % Contour cerveau
        \draw[popdark, thick] (0,0) .. controls (1,2) and (3,2.5) .. (4,1.5) .. controls (4.5,0) and (3.5,-1.5) .. (2,-2) .. controls (0.5,-2) and (-0.5,-1) .. (0,0);
        % Hémisphères
        \draw[popdark] (2,0) .. controls (2,1) and (2.5,1.5) .. (3,1);
        % Artère rompue
        \draw[popblue, very thick] (1.5,0.5) -- (2.2, 1.0);
        % Rupture
        \draw[popink, very thick, decorate, decoration={zigzag, segment length=2mm, amplitude=1mm}] (2.2, 1.0) -- (2.8, 1.3);
        \node[popink, anchor=west] at (2.8, 1.3) {\textbf{Rupture}};
        % Hématome (sang extravasé)
        \fill[poppink!60] (2.5, 0.8) circle (0.4);
        \draw[popink, thick] (2.5, 0.8) circle (0.4);
        \node[popink, align=center] at (2.5, 1.4) {\textbf{Hématome}\\(compression)};
        % Zone nécrosée autour
        \pattern[pattern=north west lines, pattern color=poppink!50] (2.0, 0.2) .. controls (2.8, 0.5) and (3.0, 0.0) .. (2.8, -0.5) .. controls (2.2, -0.8) and (1.8, -0.3) .. cycle;
        \node[popink, align=center] at (3.5, 0.0) {Zone nécrosée\\(compression + toxique)};
        % Flèche sang qui sort
        \draw[->, popblue, thick] (1.5,0.5) -- (2.0, 0.8);
        \draw[->, poppink, thick] (2.4, 1.0) -- (2.8, 1.1);
        \node[popdark, align=center, font=\bfseries] at (2, -2.5) {AVC Hémorragique\\(Rupture ≈ 20\%)};
    \end{scope}
\end{tikzpicture}
\legende{Schéma comparatif des deux mécanismes d'AVC. À gauche : AVC ischémique par obstruction artérielle (thrombe) entraînant une zone ischémique en aval. À droite : AVC hémorragique par rupture artérielle formant un hématome compressif et toxique pour les neurones environnants.}
\end{popfigure}

\begin{aretenir}[Conséquence cellulaire commune]
Quelle que soit la cause (obstruction ou rupture), les \textbf{neurones} de la zone atteinte sont privés d'oxygène et de glucose. La respiration cellulaire s'arrête, l'ATP n'est plus produit, les pompes ioniques (Na+/K+) cessent de fonctionner → entrée massive d'eau dans la cellule (œdème cytotoxique) et mort cellulaire (nécrose) en \textbf{quelques minutes}. Les fonctions commandées par cette zone (mobilité, langage, vision, sensibilité) sont perdues brutalement.
\end{aretenir}

\begin{savaistu}[L'AVC en Afrique : une urgence de santé publique]
Au Cameroun et en Afrique subsaharienne, l'AVC est devenu la \textbf{première cause de handicap acquis de l'adulte} et une cause majeure de décès. Contrairement aux idées reçues, il ne touche pas seulement les personnes âgées : l'hypertension artérielle sévère, souvent non dépistée ni traitée, provoque des AVC chez des sujets de \textbf{30 à 50 ans}, en plein âge actif. La prise en charge est rendue difficile par le coût des examens (scanner cérébral indispensable pour différencier ischémique et hémorragique) et l'accès limité aux unités neuro-vasculaires spécialisées. La prévention (dépistage et traitement de l'hypertension) et la reconnaissance précoce par la population sont donc les leviers les plus efficaces.
\end{savaistu}

\courssub{Reconnaître un AVC : la méthode « Visage – Bras – Parole » (VBP)}

Le temps, c'est du cerveau (\emph{Time is brain}). Chaque minute de retard entraîne la mort d'environ \textbf{2 millions de neurones}.

\begin{exemplebox}[Donnée chiffrée clé]
\begin{center}
\textbf{Environ 2 millions de neurones meurent chaque minute} pendant un AVC non traité. \\
Le cerveau perd l'équivalent de 3,6 années de vieillissement normal par heure sans traitement.
\end{center}
C'est pourquoi la reconnaissance immédiate par l'entourage (famille, collègues, passants) est le maillon essentiel de la chaîne de survie.
\end{exemplebox}

\begin{methode}[Reconnaître un AVC en 3 gestes (Méthode VBP)]
Face à une personne qui ne va pas bien brutalement (plainte, chute, comportement bizarre), effectuez ces trois tests simples :
\begin{enumerate}
    \item \textbf{V – Visage (Face) :} Demandez à la personne de \textbf{sourire} ou de montrer ses dents. \textbf{Asymétrie ?} Un côté de la bouche ne bouge pas, la lèvre tombe (bave possible).
    \item \textbf{B – Bras (Arms) :} Demandez à la personne de \textbf{lever les deux bras} en même temps, paumes vers le haut, les yeux fermés, pendant 10 secondes. \textbf{Un bras retombe-t-il ?} Ne peut-il pas monter ? Dérive-t-il vers l'extérieur ?
    \item \textbf{P – Parole (Speech) :} Demandez à la personne de \textbf{répéter une phrase simple} (ex. : « Le ciel est bleu » ou « Je vais bien »). \textbf{Parole trouble, incompréhensible ?} Incapacité à trouver les mots ? Ne comprend pas ce qu'on lui dit ?
\end{enumerate}
\textbf{Si UN SEUL de ces signes est présent : c'est un AVC jusqu'à preuve du contraire. Appelez les secours immédiatement (SAMU, pompiers, hôpital le plus proche).}
\end{methode}

\begin{popfigure}
\begin{tikzpicture}[scale=0.85, every node/.style={font=\small}]
    % Geste 1 : Visage
    \begin{scope}[xshift=-5.5cm]
        \draw[popdark, thick] (0,0) circle (1.2);
        % Yeux
        \draw[popdark] (-0.4, 0.2) circle (0.1);
        \draw[popdark] (0.4, 0.2) circle (0.1);
        % Bouche normale (pointillés)
        \draw[popmuted, dashed] (-0.5, -0.4) .. controls (0, -0.7) .. (0.5, -0.4);
        % Bouche déviée (trait plein poporange)
        \draw[poporange, very thick] (-0.5, -0.4) .. controls (0, -0.5) .. (0.5, -0.2);
        \node[poporange, anchor=north] at (0, -0.8) {Bouche déviée};
        \node[popdark, font=\bfseries, anchor=south] at (0, 1.4) {\textbf{1. SOURIRE}};
        \node[popblue, anchor=north] at (0, -1.3) {Asymétrie faciale ?};
    \end{scope}

    % Geste 2 : Bras
    \begin{scope}[xshift=0cm]
        % Personnage simplifié
        \draw[popdark, thick] (0,0) circle (0.6); % Tête
        \draw[popdark, thick] (0,-0.6) -- (0,-1.8); % Tronc
        % Bras gauche (normal)
        \draw[popblue, very thick] (0,-0.8) -- (-0.8, -1.2);
        \draw[popblue, very thick] (-0.8, -1.2) -- (-1.2, -1.6);
        \node[popblue, anchor=east] at (-1.3, -1.6) {Bras tient};
        % Bras droit (retombe)
        \draw[poporange, very thick] (0,-0.8) -- (0.8, -1.0);
        \draw[poporange, very thick, decorate, decoration={zigzag, segment length=1.5mm, amplitude=0.5mm}] (0.8, -1.0) -- (1.0, -1.6);
        \node[poporange, anchor=west] at (1.1, -1.6) {Bras retombe};
        \node[popdark, font=\bfseries, anchor=south] at (0, 0.8) {\textbf{2. LEVER LES BRAS}};
        \node[popblue, anchor=north] at (0, -2.2) {Un bras ne monte pas ?};
    \end{scope}

    % Geste 3 : Parole
    \begin{scope}[xshift=5.5cm]
        \draw[popdark, thick] (0,0) circle (1.2);
        \draw[popdark] (-0.4, 0.2) circle (0.1);
        \draw[popdark] (0.4, 0.2) circle (0.1);
        \draw[popdark] (-0.3, -0.3) .. controls (0, -0.5) .. (0.3, -0.3);
        % Bulle de parole
        \draw[popblue, thick] (1.3, 0.5) .. controls (2, 1) and (2.5, 1) .. (2.8, 0.5) .. controls (3, 0) and (2.5, -0.2) .. (1.8, -0.2) .. controls (1.3, -0.2) and (1.2, 0.2) .. (1.3, 0.5);
        \node[popblue, align=center] at (2.1, 0.3) {« Le ciel... \\ euh... bleu... »};
        \node[popdark, font=\bfseries, anchor=south] at (0, 1.4) {\textbf{3. RÉPÉTER}};
        \node[popblue, anchor=north] at (0, -1.5) {Parole trouble ?};
    \end{scope}
\end{tikzpicture}
\legende{Les trois gestes du test de reconnaissance rapide d'un AVC (méthode VBP). 1. Demander de sourire : rechercher une asymétrie de la bouche. 2. Demander de lever les deux bras : rechercher une faiblesse ou une chute d'un bras. 3. Demander de répéter une phrase : rechercher une parole trouble ou une incompréhension.}
\end{popfigure}

\begin{attention}[Piège fréquent : Ne rien donner à boire ni à manger !]
\textbf{Ne donnez JAMAIS d'eau, de médicament (aspirine, antihypertenseur), ni de nourriture à une victime suspecte d'AVC.}
\begin{itemize}
    \item \textbf{Risque de fausse route :} Les troubles de la déglutition (dysphagie) sont très fréquents et souvent invisibles. Le liquide ou le comprimé peut passer dans les poumons → pneumonie d'inhalation grave, étouffement.
    \item \textbf{Aggravation possible :} Si c'est un AVC hémorragique, l'aspirine (antiagrégant) aggrave le saignement. Si c'est un AVC ischémique, l'aspirine n'est pas le traitement de première intention en urgence et le diagnostic doit être confirmé par scanner à l'hôpital.
    \item \textbf{Anesthésie :} La victime peut devoir être opérée en urgence → estomac vide obligatoire.
\end{itemize}
\end{attention}

\courssub{Conduite à tenir : Premiers secours et Position Latérale de Sécurité (PLS)}

\begin{experience}[TP N°8 – 2e partie : Conduite à tenir face à un AVC (Simulation)]
\textbf{Objectif :} Savoir alerter, installer la victime en PLS si inconsciente, et transmettre les informations clés aux secours.
\textbf{Matériel :} Mannequin ou élève volontaire, téléphone (simulation), chronomètre.
\textbf{Protocole :}
\begin{enumerate}
    \item \textbf{Sécuriser :} Protéger la victime (éloigner du danger, circulation).
    \item \textbf{Évaluer :} Tester VBP (Visage, Bras, Parole). Noter l'\textbf{heure exacte d'apparition des premiers signes} (ou heure du dernier état connu normal).
    \item \textbf{Alerter :} Appeler le \textbf{112} (numéro d'urgence européen, valable au Cameroun) ou le SAMU local / Pompiers (119). Dire : « Suspicion d'AVC, test VBP positif, heure de début : HH:MM, adresse précise, état de conscience ».
    \item \textbf{Installer :}
        \begin{itemize}
            \item \textbf{Victime consciente :} Allongée, tête et buste légèrement surélevés (30°), rassurée, au calme, surveillée en continu.
            \item \textbf{Victime inconsciente qui respire :} \textbf{Position Latérale de Sécurité (PLS)} (voir schéma ci-dessous).
        \end{itemize}
    \item \textbf{Ne rien donner :} Ni boire, ni manger, ni médicament.
    \item \textbf{Surveiller :} Respiration, conscience, jusqu'à l'arrivée des secours. Préparer les papiers d'identité, ordonnances, carte de groupe sanguin.
\end{enumerate}
\textbf{Interprétation :} La rapidité de l'alerte et la qualité de l'information transmise (heure de début, signes VBP, antécédents HTA/diabète, traitements) permettent à l'équipe médicale de l'hôpital de réaliser le scanner cérébral et de décider du traitement adapté (médicamenteux ou chirurgical) pour sauver le tissu cérébral menacé.
\end{experience}

\begin{definition}[Position Latérale de Sécurité (PLS)]
Position stable, sur le côté, permettant de maintenir les voies aériennes supérieures libres (langue en avant, écoulement des liquides par la bouche) chez une victime \textbf{inconsciente qui respire}, en attendant les secours. Elle prévient l'étouffement par obstruction ou inhalation.
\end{definition}

\begin{popfigure}
\begin{tikzpicture}[scale=0.75, every node/.style={font=\small}]
    % Étape 1 : Position initiale
    \begin{scope}[xshift=-7cm]
        \draw[popdark, thick] (0,0) ellipse (1 and 0.5); % Tête
        \draw[popdark, thick] (0,-0.5) -- (0,-2.5); % Corps
        \draw[popdark, thick] (0,-0.8) .. controls (-1,-1) .. (-1.5,-1.5); % Bras G
        \draw[popdark, thick] (0,-0.8) .. controls (1,-1) .. (1.5,-1.5); % Bras D
        \draw[popdark, thick] (0,-2.5) .. controls (-0.5,-3.5) .. (-0.5,-4); % Jambe G
        \draw[popdark, thick] (0,-2.5) .. controls (0.5,-3.5) .. (0.5,-4); % Jambe D
        \node[popdark, font=\bfseries, anchor=south] at (0, 0.8) {1. Allongé sur le dos};
        \node[popblue, anchor=north] at (0, -4.5) {Vérifier respiration};
    \end{scope}

    % Étape 2 : Préparation bras/jambe
    \begin{scope}[xshift=-2cm]
        \draw[popdark, thick] (0,0) ellipse (1 and 0.5);
        \draw[popdark, thick] (0,-0.5) -- (0,-2.5);
        % Bras proche (bras du côté où on va tourner) -> plié, main joue
        \draw[popblue, very thick] (0,-0.8) .. controls (0.5,-0.5) .. (1, 0) node[anchor=west] {Main sur joue};
        % Bras loin -> perpendiculaire
        \draw[popblue, very thick] (0,-0.8) -- (-1.5, -0.8) node[anchor=east] {Bras à 90°};
        % Jambe loin pliée
        \draw[popblue, very thick] (0,-2.5) .. controls (0.5,-3) .. (1, -2.5) node[anchor=west] {Genou plié};
        \node[popdark, font=\bfseries, anchor=south] at (0, 0.8) {2. Préparer le virage};
    \end{scope}

    % Étape 3 : Basculement
    \begin{scope}[xshift=3cm]
        \draw[popdark, thick] (0,0) ellipse (1 and 0.5); % Tête sur le côté
        \draw[popdark, thick] (-0.5,-0.5) -- (-0.5,-2.5); % Corps sur le côté
        % Bras proche (sous la tête)
        \draw[popblue, very thick] (-0.5,-0.8) -- (-1.5, -0.5) node[anchor=east] {Main joue};
        % Bras loin (devant, stabilisateur)
        \draw[popblue, very thick] (-0.5,-0.8) -- (0.5, -1.2) node[anchor=west] {Bras avant};
        % Jambe loin pliée (stabilisatrice)
        \draw[popblue, very thick] (-0.5,-2.5) .. controls (0,-3) .. (0.5, -2.5) node[anchor=west] {Genou à 90°};
        % Flèche rotation
        \draw[->, poporange, thick] (1.5, 0) arc (0:180:1.5 and 0.5);
        \node[popdark, font=\bfseries, anchor=south] at (0, 0.8) {3. Basculer (tirer genou)};
    \end{scope}

    % Étape 4 : PLS finale
    \begin{scope}[xshift=8cm]
        \draw[popdark, thick] (0,0) ellipse (1 and 0.5); % Tête
        \draw[popdark, thick] (-0.5,-0.5) -- (-0.5,-2.5); % Corps
        \draw[popblue, very thick] (-0.5,-0.8) -- (-1.5, -0.5); % Bras sous tête
        \draw[popblue, very thick] (-0.5,-0.8) -- (0.5, -1.2); % Bras avant
        \draw[popblue, very thick] (-0.5,-2.5) .. controls (0,-3) .. (0.5, -2.5); % Jambe pliée
        % Bouche ouverte, écoulement
        \draw[popink, thick] (-0.8, -0.1) -- (-0.4, -0.1);
        \draw[popink, ->, thick] (-0.6, -0.1) -- (-0.6, -0.5) node[below] {Écoulement};
        \node[popgreen, font=\bfseries, anchor=south] at (0, 0.8) {4. PLS FINALE};
        \node[popgreen, align=center, anchor=north] at (0, -3.2) {Voies aériennes libres\\Tête en extension\\Surveillance continue};
    \end{scope}
\end{tikzpicture}
\legende{Mise en place de la Position Latérale de Sécurité (PLS) pour une victime d'AVC inconsciente qui respire. 1. Allongée sur le dos. 2. Bras du côté du secouriste plié main sur joue opposée, bras loin à 90°, jambe loin pliée. 3. Basculer en tirant sur le genou plié. 4. Position finale stable : tête en extension (menton en avant) pour ouvrir les voies aériennes, bouche vers le bas pour l'écoulement, corps stabilisé par le genou à 90° et le bras antérieur.}
\end{popfigure}

\begin{aretenir}[Informations vitales à transmettre aux secours (SAMU / 112 / Pompiers)]
\begin{enumerate}
    \item \textbf{Heure précise de début des signes} (ou heure du dernier contact normal). \textit{C'est le critère majeur pour la décision médicale à l'hôpital.}
    \item \textbf{Signes VBP positifs} (Visage, Bras, Parole).
    \item \textbf{État de conscience} (répond ? ouvre les yeux ? bouge ?).
    \item \textbf{Antécédents connus} : Hypertension artérielle (traitée ?), Diabète, Cardiopathie, AVC antérieur.
    \item \textbf{Traitements en cours} : Anticoagulants, Antiagrégants (Aspirine), Antihypertenseurs, Antidiabétiques.
    \item \textbf{Adresse exacte / Point de repère} et numéro de téléphone du témoin.
\end{enumerate}
\end{aretenir}

\courssub{Prévention des AVC : agir sur les facteurs de risque modifiables}

La meilleure lutte contre l'AVC reste la prévention primaire (avant le premier AVC) et secondaire (après un premier AVC ou un AIT - Accident Ischémique Transitoire).

\begin{propriete}[Principaux facteurs de risque modifiables de l'AVC]
\begin{center}
\begin{tabularx}{\linewidth}{|l|X|}\hline
\textbf{Facteur de risque} & \textbf{Mesure de prévention (Hygiène de la circulation)} \\ \hline
\textbf{Hypertension artérielle (HTA)} & \textbf{Cause n°1.} Dépistage régulier (au moins 1 fois/an). Traitement médicamenteux \textbf{au long cours} (observance stricte). Objectif PA < 140/90 mmHg. \\ \hline
\textbf{Diabète de type 2} & Équilibre glycémique. Alimentation à IG bas, activité physique, traitement prescrit. \\ \hline
\textbf{Dyslipidémie (excès LDL-c)} & Traitement si risque cardiovasculaire élevé (statines). Alimentation pauvre en graisses saturées/trans. \\ \hline
\textbf{Tabagisme} & \textbf{Arrêt total} (sevrage tabagique). Risque divisé par 2 à 5 ans après arrêt. \\ \hline
\textbf{Alcool} & Consommation modérée (max 2 verres/jour, pas tous les jours). \\ \hline
\textbf{Sédentarité / Obésité} & Activité physique régulière (150 min/semaine d'endurance modérée). Perte de poids si IMC > 25. \\ \hline
\textbf{Alimentation riche en sel} & \textbf{< 5 g de sel/jour (OMS).} Limiter bouillons cubes, conserves, plats préparés, sel à table. Privilégier légumes, fruits, poisson, céréales complètes. \\ \hline
\textbf{Fibrillation auriculaire} & Traitement anticoagulant prescrit par le médecin si risque élevé. \\ \hline
\end{tabularx}
\end{center}
\end{propriete}

\begin{savaistu}[L'AIT : Signal d'alarme à ne pas ignorer]
Un \textbf{Accident Ischémique Transitoire (AIT)} se manifeste par des signes neurologiques identiques à l'AVC (paralysie, trouble de la parole, cécité monolatérale) mais \textbf{réversibles en moins d'une heure (souvent quelques minutes)}, sans séquelle visible à l'examen clinique. C'est un \textbf{signal d'alerte majeur} : le risque d'AVC complet dans les 48h à 3 mois est très élevé (10 à 20\%). \textbf{Conduite à tenir :} Appeler le 112 / se rendre aux urgences immédiatement, même si « ça va mieux ». C'est une urgence pour rechercher la cause (sténose carotide, fibrillation auriculaire, HTA) et prévenir l'AVC définitif.
\end{savaistu}

\begin{aretenir}[Neurones et récupération]
Les \textbf{neurones du système nerveux central (cerveau, moelle épinière) ne se divisent pas et ne se régénèrent pas} après une nécrose. La récupération fonctionnelle observée après un AVC repose sur la \textbf{plasticité cérébrale} (réorganisation des réseaux par des zones saines) et la rééducation (kinésithérapie, orthophonie, ergothérapie), mais la zone nécrosée laisse une cicatrice définitive. D'où l'impératif absolu : \textbf{agir vite pour limiter l'étendue de la lésion}.
\end{aretenir}

\courssub{Bilan récapitulatif : Les accidents de la circulation sanguine}

Ce tableau synthétise les trois grands types d'accidents de la circulation étudiés dans cette leçon.

\begin{center}
\begin{tabularx}{\linewidth}{|l|X|X|X|}\hline
\rowcolor{popblueL}
\textbf{Accident} & \textbf{Mécanisme principal / Cause} & \textbf{Signes d'alerte / Conséquences} & \textbf{Moyens de lutte / Prise en charge} \\ \hline
\textbf{Hémorragie externe} & Rupture d'un vaisseau (trauma, varice) → Perte de sang hors du corps. & Flux sanguin visible, pâleur, soif, pouls rapide, hypotension, perte de conscience. & \textbf{Compression manuelle} (point de compression ou garrot si membre), \textbf{Allongement}, \textbf{Appel secours}, \textbf{Transfusion} si choc hémorragique. \\ \hline
\textbf{Maladies cardiovasculaires (Athérosclérose, HTA, Insuffisance cardiaque)} & Dépôt de plaques d'athérome (cholestérol, inflammation) → Rétrécissement artères (coronaires, carotides, membres). HTA → Surcharge cœur. & Angor de poitrine (effort), dyspnée, œdèmes, claudication, signes AVC/AIT. & \textbf{Prévention :} Hygiène de vie (sel, tabac, activité, poids). \textbf{Traitement :} Statines, antihypertenseurs, antiagrégants, bêta-bloquants, chirurgie (pontage, stent, endartériectomie). \\ \hline
\textbf{Accident Vasculaire Cérébral (AVC)} & \textbf{Ischémique (80\%) :} Thrombose/Embolie artère cérébrale. \textbf{Hémorragique (20\%) :} Rupture artère (HTA). & \textbf{Soudain :} Visage asymétrique, Bras faible, Parole trouble (VBP). Hémiparésie, aphasie, troubles visuels, céphalée violente. & \textbf{URGENCE VITALE (112/SAMU).} \textbf{Heure de début.} \textbf{PLS si inconscient.} \textbf{Rien par la bouche.} \textbf{Scanner cérébral urgent.} Traitement médical/chirurgical spécialisé. Rééducation précoce. \\ \hline
\end{tabularx}
\end{center}

\begin{exoresolu}[Cas clinique : M. K., 52 ans, cadre à Douala]
\textbf{Énoncé :} M. K., 52 ans, connu hypertendu (traitement irrégulier), fumeur, diabétique non équilibré, s'effondre brutalement au bureau à 10h15. Son collègue le trouve : visage dévié à droite, bras gauche inerte, ne parvient pas à articuler « Bonjour ». Il est conscient mais ne comprend pas les questions. Le collègue appelle le 112 à 10h18.
\begin{enumerate}
    \item Quel type d'AVC suspectez-vous ? Justifiez.
    \item Citez 3 facteurs de risque présents chez ce patient.
    \item Quelle conduite le collègue a-t-il bien tenue ? Quelle erreur doit-il éviter absolument ?
    \item Pourquoi l'heure 10h15 est-elle cruciale pour l'équipe médicale ?
\end{enumerate}
\tcblower
\textbf{Solution détaillée :}
\begin{enumerate}
    \item \textbf{AVC ischémique probable (infarctus de l'hémisphère droit).} Justification : Installation brutale, déficit neurologique focalisé. Le visage dévié vers la droite (bouche tire à droite) indique une paralysie faciale gauche (atteinte controlatérale de l'hémisphère droit). Le bras gauche inerte confirme une hémiparésie gauche, siège de l'hémisphère droit. Le trouble de la compréhension (ne comprend pas les questions) peut s'expliquer par une atteinte de zones spécifiques de l'hémisphère droit (non dominant pour le langage mais impliqué dans la compréhension prosodique/complexe) ou une extension. Le terrain (HTA, diabète, tabac) favorise l'ischémique (80\% des cas). Seul le scanner à l'hôpital confirmera le diagnostic.
    \item 1. Hypertension artérielle (mal contrôlée/irrégulière). 2. Diabète de type 2 (déséquilibré). 3. Tabagisme actif. (L'âge > 50 ans et le sexe masculin sont des facteurs non modifiables).
    \item \textbf{Bien :} A reconnu les signes (VBP), a appelé le 112 immédiatement (10h18), a noté l'heure de début (10h15).
    \item \textbf{Erreur à éviter :} Ne pas donner d'eau, d'aspirine, ni de médicament pour la tension. Ne pas faire asseoir ou lever le patient. L'allonger, tête légèrement surélevée, surveiller.
    \item L'heure 10h15 (début des signes) détermine la \textbf{fenêtre thérapeutique} pour l'hôpital. Plus le délai est court, plus les chances de bénéficier d'un traitement de réouverture de l'artère (thrombolyse ou thrombectomie) sont grandes, sauvant ainsi du tissu cérébral encore viable. Chaque minute compte.
\end{enumerate}
\end{exoresolu}

\begin{exercice}[Entraînement : Réagir face à l'AVC]
Mme T., 63 ans, prépare le repas. Son mari la trouve assise par terre, la tasse de thé cassée à côté d'elle. Elle ne répond pas quand il lui parle, son œil gauche ne suit pas, sa bouche tire à droite. Elle respire bruyamment.
\begin{enumerate}
    \item En utilisant la méthode VBP, que doit vérifier le mari ? (Il ne peut pas lui demander de lever les bras car elle est par terre).
    \item Quelle est la première action à faire après avoir sécurisé les lieux ?
    \item La victime est inconsciente mais respire. Quelle position installer ? Dessinez le schéma simplifié de cette position en indiquant l'orientation de la tête et le rôle du genou plié.
    \item Listez 3 informations que le mari doit donner au 112.
\end{enumerate}
\end{exercice}

\begin{corrige}[Exercice : Réagir face à l'AVC]
\begin{enumerate}
    \item \textbf{V (Visage) :} Demander de sourire / montrer les dents $\rightarrow$ Asymétrie (bouche tire à droite = paralysie faciale gauche). \textbf{P (Parole) :} Lui parler, lui demander de répéter / nommer un objet $\rightarrow$ Pas de réponse, trouble de la compréhension ou de l'expression. \textbf{B (Bras) :} Impossible de tester le lever des bras (position assise/au sol). On peut tester la préhension des mains (serrer les doigts du secouriste) $\rightarrow$ Faiblesse/main gauche molle (hémiparésie gauche).
    \item \textbf{Appeler le 112 (ou SAMU/Pompiers 119)} immédiatement en précisant : « Suspicion AVC, femme 63 ans, inconsciente, respire, signes VBP positifs (visage dévié droite, hémiparésie gauche, trouble conscience/langage), heure découverte : [heure actuelle], adresse exacte ».
    \item \textbf{Position Latérale de Sécurité (PLS) côté droit (côté sain / côté non paralysé).} Schéma : Corps sur le côté droit, tête en extension (menton en avant) pour ouvrir voies aériennes, bouche ouverte vers le bas (écoulement salive/thé), bras droit (côté sol) plié main sous la joue, bras gauche (côté haut) fléchi coude/main devant le thorax (stabilisateur), jambe gauche (côté haut) pliée genou à 90° au sol (butée anti-retournement), jambe droite tendue.
    \item 1. Heure exacte de survenue (ou heure découverte si inconnue). 2. Signes neurologiques observés (visage dévié droite, hémiparésie gauche, trouble conscience/langage). 3. Antécédents : HTA ? Diabète ? Traitements (anticoagulants ?). Identité et téléphone du mari.
\end{enumerate}
\end{corrige}

\begin{aretenir}[Synthèse : La chaîne de survie AVC]
\begin{center}
\textbf{RECONNAÎTRE (VBP)} $\xrightarrow{\text{Immédiat}}$ \textbf{APPELER (112/SAMU + Heure)} $\xrightarrow{\text{Minutes}}$ \textbf{INSTALLER (PLS si inconscient / Allongé tête 30° si conscient)} $\xrightarrow{\text{Attente}}$ \textbf{SURVEILLER (Respiration, Conscience)} $\xrightarrow{\text{Urgences}}$ \textbf{SCANNER CÉRÉBRAL} $\xrightarrow{\text{Traitement}}$ \textbf{PRISE EN CHARGE MÉDICALE SPÉCIALISÉE} $\xrightarrow{\text{Réanimation}}$ \textbf{RÉÉDUCATION PRÉCOCE}
\end{center}
\textbf{Time is Brain. Chaque minute compte.}
\end{aretenir}

\newpage

\coursec{Activité d'intégration}

\courssub{Objectifs de l'activité}
Mobiliser l'ensemble des savoirs et savoir-faire de la leçon « Circulation sanguine » pour analyser une situation d'urgence vitale réelle : un accident vasculaire cérébral (AVC) survenu en milieu scolaire. L'élève doit être capable d'identifier les facteurs de risque cardiovasculaire, d'expliquer le mécanisme physiopathologique reliant l'athérosclérose à l'ischémie cérébrale, de reconnaître les signes d'alerte (méthode \cle{VITE}), de rédiger une conduite à tenir de premiers secours justifiée scientifiquement, et de proposer des mesures de prévention primaire adaptées au contexte camerounais.

\courssub{Situation : Alerte au lycée — Sauver Monsieur Ngo Bassa}
\label{sit:ngobassa}
\textbf{Contexte.} Nous sommes un mardi 14 mars 2023, 10 h 15, dans la cour du Lycée Bilingue de Nkolbisson (Yaoundé). La température est de \SI{31}{\celsius}, hygrométrie \SI{78}{\percent}. Monsieur \textbf{Ngo Bassa}, 55 ans, surveillant général, connu pour son autorité bienveillante, parcourt la cour pour rappeler des élèves en retard. Soudain, il s'arrête net, porte la main gauche à l'épaule droite, son visage s'affaisse du côté droit, il laisse tomber sa canne et s'effondre au sol, conscient mais incapable de parler distinctement. Trois élèves de 3ème (Aïcha, Kevin, Stéphanie), formés aux gestes qui sauvent lors du TP N°8, interviennent immédiatement pendant qu'un autre court prévenir l'infirmerie et composer le \textbf{112} (numéro d'urgence unique au Cameroun).

\medskip
\noindent \textbf{Document 1 : Cœur et artères coronaires de M. Ngo Bassa (schéma d'angio-IRM simplifié).}
\begin{popfigure}
\begin{tikzpicture}[scale=1.1, every node/.style={font=\small}]
% Heart outline
\draw[popdark, line width=1.5pt, fill=poppinkL!30] 
  (0,0) .. controls (-1.5,1) and (-1.5,2.5) .. (0,3.5) 
  .. controls (1.5,2.5) and (1.5,1) .. (0,0);
% Right ventricle suggestion
\draw[popdark, line width=0.8pt] (0,0) .. controls (1,1) and (1,2) .. (0,3);
% Aorta
\draw[popred, line width=2.5pt] (0,3.5) -- (0,4.5) node[above, popdark] {Aorte};
% Pulmonary trunk
\draw[popblue, line width=2pt] (0.3,3.5) -- (1.2,4.2) node[right, popdark] {Tronc pulmonaire};
% Right Coronary Artery (RCA)
\draw[popred, line width=1.8pt] (0.2,0.2) .. controls (1.5,0.5) and (1.5,2) .. (0.5,3);
% Left Coronary Artery (LCA) - Circumflex
\draw[popred, line width=1.8pt] (-0.2,0.2) .. controls (-1.5,0.5) and (-1.5,2) .. (-0.5,3);
% LAD (Anterior Interventricular)
\draw[popred, line width=1.8pt] (0,0.5) -- (0,3);
% ATHEROMA PLAQUE on LAD
\draw[fill=poporange, poporange, line width=1pt] (0,1.8) ellipse (0.35 and 0.15);
\node[poporange, font=\bfseries\small] at (0,1.3) {Plaque d'athérome};
\draw[<-, poporange, line width=1pt] (0.5,1.3) -- (0.35,1.6);
% Labels
\node[popdark] at (-2.2,1.5) {Oreillette G};
\node[popdark] at (2.2,1.5) {Oreillette D};
\node[popdark] at (-2.2,0) {Ventricule G};
\node[popdark] at (2.2,0) {Ventricule D};
\draw[<-, popdark] (-2,1.5) -- (-0.5,1.8);
\draw[<-, popdark] (2,1.5) -- (0.5,1.8);
\draw[<-, popdark] (-2,0) -- (-0.5,0.5);
\draw[<-, popdark] (2,0) -- (0.5,0.5);
\end{tikzpicture}
\legende{Document 1 : Vue antérieure du cœur. L'artère interventriculaire antérieure (AIV, branche de la coronaire gauche) présente une sténose significative (flèche orange) due à une plaque d'athérome (orange), réduisant le lumen artériel d'environ 70 \%. Ce signe témoigne d'une athérosclérose sévère et systémique.}
\end{popfigure}

\medskip
\noindent \textbf{Document 2 : Habitudes de vie et antécédents médicaux de M. Ngo Bassa (extrait du dossier médical de l'infirmerie scolaire, mis à jour janvier 2023).}
\begin{center}
\begin{tabularx}{\linewidth}{|l|X|}\hline
\rowcolor{popblueL} \textbf{Paramètre} & \textbf{Valeur / Description} \\ \hline
Âge & 55 ans \\ \hline
Sexe & Masculin \\ \hline
Profession & Surveillant général (poste sédentaire, stress fréquent) \\ \hline
Tabagisme & 20 cigarettes/jour depuis 35 ans (\textbf{35 années-paquet}) \\ \hline
Alimentation & Riche en sel (bouillons cubes, poisson fumé, sauces grasses), pauvre en fruits/légumes, 3 repas copieux/jour \\ \hline
Activité physique & Aucune activité sportive régulière (marche < 15 min/jour) \\ \hline
Tension artérielle (au repos, infirmerie, 01/2023) & \textbf{\SI{170/100}{\mmHg}} (Grade 2 HTA) \\ \hline
Glycémie à jeun & \SI{1,15}{\g\per\L} (prédiabète) \\ \hline
Cholestérol total & \SI{2,65}{\g\per\L} (LDL-C : \SI{1,95}{\g\per\L}) \\ \hline
Antécédents familiaux & Père décédé d'AVC à 62 ans ; Mère diabétique type 2 \\ \hline
Traitement actuel & \emph{Irrégulier} : Amlodipine 5 mg (1/j) + Hydrochlorothiazide 12,5 mg (1/j) \\ \hline
\end{tabularx}
\end{center}

\medskip
\noindent \textbf{Document 3 : Observations des élèves témoins et évolution de la pression artérielle (2018–2023).}
\begin{itemize}
  \item \textbf{Témoignage d'Aïcha (3ème A) :} « Sa bouche était tordue vers la gauche, il ne pouvait pas lever le bras droit, il essayait de dire "aidez-moi" mais ça sortait "a... a... zé". Il transpirait beaucoup, regardait vers nous mais ne répondait pas quand on lui serrait la main. »
  \item \textbf{Témoignage de Kevin (3ème B) :} « On a regardé l'heure : 10 h 17. Il respirait fort, 22 cycles/min. Pouls radial gauche : 98 battements/min, irrégulier. »
  \item \textbf{Courbe de la Pression Artérielle Systolique (PAS) et Diastolique (PAD) au repos (mesures infirmerie annuelle).}
\end{itemize}
\begin{popfigure}
\begin{tikzpicture}
\begin{axis}[
  width=\linewidth, height=7cm,
  xlabel={Année}, ylabel={Pression artérielle (mmHg)},
  xmin=2017.5, xmax=2023.5, ymin=110, ymax=190,
  xtick={2018,2019,2020,2021,2022,2023},
  ytick={120,130,140,150,160,170,180},
  grid=major, grid style={popline},
  legend style={at={(0.02,0.98)}, anchor=north west, fill=popcream, draw=popdark, font=\small},
  popcurve/.style={thick, mark=*, mark size=2.5},
]
\addplot[popred, popcurve] coordinates {
  (2018,138) (2019,142) (2020,152) (2021,160) (2022,168) (2023,170)
};
\addlegendentry{PAS (Systolique)}
\addplot[popblue, popcurve] coordinates {
  (2018,85) (2019,88) (2020,92) (2021,96) (2022,98) (2023,100)
};
\addlegendentry{PAD (Diastolique)}
% Seuil HTA
\draw[popmuted, dashed, thick] (2017.5,140) -- (2023.5,140) node[right, popmuted, font=\tiny] {Seuil HTA (\SI{140}{\mmHg})};
\draw[popmuted, dashed, thick] (2017.5,90) -- (2023.5,90) node[right, popmuted, font=\tiny] {Seuil HTA (\SI{90}{\mmHg})};
\end{axis}
\end{tikzpicture}
\legende{Document 3 : Évolution de la tension artérielle de M. Ngo Bassa sur 5 ans. Aggravation progressive malgré traitement prescrit (observance faible).}
\end{popfigure}

\medskip
\noindent \textbf{Document 4 : Rappel des gestes de premiers secours face à un AVC suspecté (affiche Ministère de la Santé Publique / Croix-Rouge Cameroun).}
\begin{center}
\begin{tabularx}{\linewidth}{|l|X|}\hline
\rowcolor{popblueL} \textbf{Étape} & \textbf{Action} \\ \hline
1. \textbf{V} — \textbf{V}isage & Demander à la personne de sourire : \emph{bouche déviée ?} \\ \hline
2. \textbf{I} — \textbf{I}nmobilité & Demander de lever les deux bras : \emph{un bras tombe ou ne monte pas ?} \\ \hline
3. \textbf{T} — \textbf{T}roubles de la parole & Demander de répéter une phrase simple : \emph{parole incompréhensible ou absente ?} \\ \hline
4. \textbf{E} — \textbf{E}xtrême urgence & \textbf{Appeler le 112 (ou 119) IMMÉDIATEMENT}. Noter l'heure de début des signes. \\ \hline
5. \textbf{Position} & Allongé, tête et buste surélevés de 30° (position demi-assise) si conscient ; PLS si inconscient. \\ \hline
6. \textbf{Interdits} & \textbf{Ne rien donner à boire/manger}, ne pas donner de médicament, ne pas faire lever le patient. \\ \hline
7. \textbf{Surveillance} & Surveiller respiration, pouls, niveau de conscience jusqu'à l'arrivée des secours. \\ \hline
\end{tabularx}
\end{center}

\courssub{Consignes}
À partir de l'analyse rigoureuse des quatre documents, rédigez une \textbf{fiche de secouriste} structurée, argumentée et destinée à être archivée au poste de secours du lycée. Cette fiche doit répondre aux quatre tâches suivantes :

\begin{enumerate}
  \item \textbf{Facteurs de risque :} Dressez la liste complète des facteurs de risque cardiovasculaire (modifiables et non modifiables) présentés par M. Ngo Bassa, en vous appuyant \emph{uniquement} sur le Document 2. Classez-les dans un tableau à deux colonnes.
  \item \textbf{Mécanisme physiopathologique :} En vous servant du Document 1 et de vos connaissances sur l'athérosclérose, expliquez \emph{pas à pas} comment la maladie athéroscléreuse a pu provoquer l'accident observé (AVC ischémique). Utilisez les termes : \cle{athérome}, \cle{sténose}, \cle{ischémie}, \cle{thrombose}, \cle{embolie}, \cle{névrones}, \cle{hémiparésie controlatérale}.
  \item \textbf{Reconnaissance et conduite à tenir :} 
  \begin{enumerate}
    \item Identifiez, dans le Document 3, les trois signes cliniques majeurs correspondant à la méthode \cle{VITE}.
    \item Rédigez la conduite à tenir immédiate (chronologie des 5 premières minutes) en justifiant \emph{scientifiquement} chaque geste (position, alerte, surveillance, interdits) par son effet physiologique.
  \end{enumerate}
  \item \textbf{Prévention :} Proposez \textbf{trois mesures de prévention primaire} concrètes, réalistes et adaptées au contexte camerounais (alimentation, activité physique, suivi médical) pour une affiche d'hygiène de vie à destination des personnels et élèves du lycée. Chaque mesure doit être formulée sous forme d'un message court, percutant et scientifiquement fondé.
\end{enumerate}

\courssub{Ressources à mobiliser}
\begin{itemize}
  \item \textbf{Savoirs :} Structure du cœur (cavités, valves, coronaires) ; circulation coronaire (vascularisation du myocarde) ; circulation cérébrale (artères carotides, cercle de Willis) ; définition de l'AVC ischémique (thrombose/embolie) et hémorragique ; facteurs de risque de l'athérosclérose (HTA, tabagisme, dyslipidémie, diabète, sédentarité, âge, sexe, hérédité) ; mécanisme de la plaque d'athérome (endothélium, LDL, macrophages, cellules spumeuses, cap fibreux) ; signes d'alerte AVC (VITE/FAST) ; gestes de premiers secours (PLS, position demi-assise, alerte 112).
  \item \textbf{Savoir-faire :} Lecture d'images médicales (schéma, courbe) ; extraction d'informations d'un tableau ; rédaction d'une argumentation scientifique structurée (observation $\rightarrow$ hypothèse $\rightarrow$ explication) ; application d'un protocole de secours ; communication claire (affiche prévention).
  \item \textbf{Repères numériques :} TA normale $< \SI{140/90}{\mmHg}$ ; HTA Grade 1 : \SI{140-159}{\mmHg} / \SI{90-99}{\mmHg} ; Grade 2 : \SI{160-179}{\mmHg} / \SI{100-109}{\mmHg} ; Cholestérol-LDL cible $< \SI{1,6}{\g\per\L}$ (sujet à risque très élevé) ; Glycémie à jeun normale $< \SI{1,10}{\g\per\L}$ ; Délai thérapeutique thrombolyse $< 4$ h 30.
\end{itemize}

\courssub{Démarche et corrigé commenté}
\begin{exoresolu}[Fiche de secouriste — M. Ngo Bassa — 14/03/2023]
\textbf{ÉNONCÉ DE LA SITUATION :} Voir documents 1 à 4 ci-dessus.

\tcblower

\textbf{SOLUTION DÉTAILLÉE ET COMMENTÉE}

\medskip
\noindent \textbf{1. Facteurs de risque cardiovasculaire de M. Ngo Bassa (Document 2)}

\begin{center}
\begin{tabularx}{\linewidth}{|>{\bfseries}l|X|}\hline
\rowcolor{popblueL} \textbf{Facteurs NON MODIFIABLES} & \textbf{Facteurs MODIFIABLES} \\ \hline
Âge : 55 ans (risque accru après 50 ans chez l'homme) & Tabagisme actif : 20 cig/jour $\times$ 35 ans = \textbf{35 années-paquet} (toxique vasculaire majeur : endothélium, LDL oxydé, plaque instable) \\ \hline
Sexe : Masculin (risque > femme avant ménopause) & Hypertension artérielle \textbf{non contrôlée} : \SI{170/100}{\mmHg} (Grade 2) — \emph{observance thérapeutique irrégulière} \\ \hline
Antécédents familiaux : Père AVC 62 ans ; Mère DT2 (prédisposition génétique athérosclérose/HTA/DT2) & Dyslipidémie : Cholestérol total \SI{2,65}{\g\per\L} ; \textbf{LDL-C \SI{1,95}{\g\per\L}} (cible $< \SI{1,6}{\g\per\L}$ pour très haut risque) \\ \hline
& Prédiabète : Glycémie \SI{1,15}{\g\per\L} (résistance à l'insuline, glycation protéines, atteinte micro/macrovasculaire) \\ \hline
& Sédentarité : $<$ 15 min marche/jour (pas de protection vasculaire par NO, pas de contrôle pondéral) \\ \hline
& Alimentation déséquilibrée : Excès sel (HTA), excès graisses saturées (LDL), déficit fibres/antioxydants (fruits/légumes) \\ \hline
& Stress professionnel chronique (catécholamines $\rightarrow$ vasoconstriction, HTA, plaque instable) \\ \hline
\end{tabularx}
\end{center}

\medskip
\noindent \textbf{2. Mécanisme physiopathologique : De l'athérosclérose systémique à l'AVC ischémique (hémiparésie droite + aphasie)}

\begin{enumerate}
  \item \textbf{Preuve d'athérosclérose systémique (Doc 1) :} La plaque d'\cle{athérome} siégeant sur l'artère interventriculaire antérieure (AIV) avec une \cle{sténose} à 70 \% témoigne d'une atteinte athéromateuse généralisée. Les mêmes facteurs de risque (HTA, tabac, LDL-C élevé, prédiabète) agressent l'endothélium de l'ensemble du réseau artériel, notamment les \textbf{artères carotides} (vascularisation cérébrale antérieure).
  \item \textbf{Genèse de la plaque carotidienne (années d'évolution) :} Au niveau de la bifurcation carotidienne (zone de turbulence), l'endothélium, fragilisé par le cisaillement HTA, le stress oxydatif tabagique et l'infiltration de LDL-C, devient perméable. Monocytes $\rightarrow$ macrophages $\rightarrow$ phagocytose LDL oxydé = \textbf{cellules spumeuses}. Prolifération musculaire lisse, dépôt collagène = \textbf{plaque fibreuse (athérome)} avec cap fibreux et noyau lipidique nécrotique.
  \item \textbf{Rupture de plaque et thrombose artérielle (événement aigu) :} La plaque carotidienne devient \emph{instable} (cap fibreux fin, inflammation). Une fissuration expose le noyau thrombogène au flux sanguin. Activation plaquettaire (adhésion/agrégation via GPIIb/IIIa, fibrinogène) + cascade coagulation $\rightarrow$ formation d'un \textbf{thrombus} (caillot) \emph{in situ} sur la plaque carotidienne.
  \item \textbf{Embolie cérébrale :} Un fragment du thrombus carotidien se détache $\rightarrow$ \textbf{embole}. Il est entraîné par le flux sanguin : Artère carotide interne $\rightarrow$ Cercle de Willis $\rightarrow$ \textbf{Artère cérébrale moyenne (ACM) GAUCHE} (car les signes moteurs/sensitifs touchent l'hémicorps DROIT et le langage est atteint $\rightarrow$ hémisphère dominant gauche).
  \item \textbf{Ischémie cérébrale et nécrose neuronale :} L'embole obstrue l'ACM gauche $\rightarrow$ arrêt brutal de l'apport en $\ce{O2}$ et glucose $\rightarrow$ \textbf{ischémie} du territoire irrigé (cortex moteur/sensitif controlatéral + aire de Broca). Les \textbf{névrones} (cellules à fort métabolisme oxydatif, sans réserve) entrent en défaillance énergétique (arrêt pompe $\ce{Na+/K+}$, dépolarisation, influx $\ce{Ca^{2+}}$, mort cellulaire par nécrose au cœur, apoptose en pénombre) en quelques minutes.
  \item \textbf{Expression clinique :} Perte de la commande motrice volontaire (voie pyramidale croisée) $\rightarrow$ \textbf{hémiparésie controlatérale droite} + atteinte de l'aire de Broca (hémisphère gauche) $\rightarrow$ \textbf{aphasie d'expression} (parole incompréhensible, « a... a... zé »). La déviation de la bouche vers la gauche confirme la paralysie faciale controlatérale (voie corticobulbaire).
\end{enumerate}

\medskip
\noindent \textbf{3. Reconnaissance des signes (VITE) et Conduite à tenir justifiée (Chronologie 10 h 17 – 10 h 22)}

\begin{center}
\begin{tabularx}{\linewidth}{|c|l|X|}\hline
\rowcolor{popblueL} \textbf{Heure} & \textbf{Action (Geste)} & \textbf{Justification scientifique (Physiologie / Pathophysiologie)} \\ \hline
10 h 17 & \textbf{V} : Vérifier Visage (sourire) $\rightarrow$ \textbf{Bouche déviée (gauche)} & Atteinte noyau facial / voie corticobulbaire controlatérale (hémisphère gauche). Signe VITE positif. \\ \hline
10 h 17 & \textbf{I} : Vérifier Inmobilité (lever bras) $\rightarrow$ \textbf{Bras droit ne monte pas} & Hémiparésie controlatérale (voie pyramidale croisée). Signe VITE positif. \\ \hline
10 h 17 & \textbf{T} : Vérifier Troubles parole (répéter phrase) $\rightarrow$ \textbf{Parole incompréhensible (aphasie)} & Atteinte aire de Broca (hémisphère dominant gauche). Signe VITE positif. \\ \hline
10 h 17 & \textbf{E} : \textbf{APPELER LE 112 (SAMU)} — Dire : « AVC suspecté, homme 55 ans, signes VITE +++, conscient, TA haute, heure début 10 h 15, Lycée Nkolbisson. » & \textbf{Temps = Cerveau.} Thrombolyse (rt-PA) possible si arrivée $< 4$ h 30. Prévenir service neurovasculaire (UNV) de l'Hôpital Central ou HGOPED. \\ \hline
10 h 18 & \textbf{Position : Demi-assise (buste 30°), tête latéralisée côté paralysie (droit).} & \emph{Diminue la pression veineuse cérébrale (retour veineux $\uparrow$), limite l'œdème cérébral, prévient la fausse route (salives) si troubles déglutition. \textbf{Pas d'allongé strict} (PIC $\uparrow$).} \\ \hline
10 h 18 & \textbf{Desserer vêtements} (col, ceinture, cravate). & Faciliter ventilation, retour veineux, confort. \\ \hline
10 h 19 & \textbf{Surveillance continue :} FR (22/min), FC (98/min, irrégulier), TA (mesure unique si brassard dispo), Conscience (Glasgow simplifié : ouvre yeux ? répond ? bouge ?). & Détecter aggravation (œdème, hémorragie méningée, arrêt cardiaque). \textbf{Noter l'heure exacte du début (10 h 15) : CRITIQUE pour décision thrombolyse.} \\ \hline
10 h 20 & \textbf{INTERDITS stricts :} Ne rien donner à boire/manger (risque fausse route $\rightarrow$ pneumopathie d'inhalation). Ne pas donner d'aspirine (risque AVC hémorragique non écarté sans scanner). Ne pas faire lever/asseoir brutalement (chute TA orthostatique, chute). & \emph{Principes de précaution : diagnostic différentiel AVC ischémique / hémorragique impossible cliniquement. Scanner cérébral obligatoire avant tout traitement.} \\ \hline
10 h 22 & \textbf{Accueil secours :} Transmettre fiche (heure, signes, TA, traitements, antécédents). & Continuité de la chaîne de survie. \\ \hline
\end{tabularx}
\end{center}

\medskip
\noindent \textbf{4. Trois mesures de prévention primaire pour l'affiche « Protège ton cœur, protège ton cerveau » (Lycée Nkolbisson)}

\begin{enumerate}
  \item \textbf{« Assiette 3 couleurs = Vaisseaux propres »} — \emph{Justification :} Remplacer le bouillon cube (sel caché $\approx$ \SI{2}{\g} NaCl/unité) par épices locales (penja, gingembre, ail) ; ajouter 1 portion légumes verts (feuilles de manioc, folong) + 1 fruit de saison (papaye, ananas) à chaque repas. \textbf{Baisse LDL-C, TA, stress oxydatif.}
  \item \textbf{« 30 min marche active / jour = Cœur costaud »} — \emph{Justification :} Marcher vite (\SI{6}{\km\per\hour}) trajet domicile-lycée ou cour de récréation. \textbf{Effet :} $\uparrow$ NO (vasodilatation), $\downarrow$ TA au repos ($-$5 à $-$8 mmHg), $\uparrow$ HDL-C, $\downarrow$ glycémie, gestion stress. Accessible à tous, gratuit.
  \item \textbf{« Ton traitement HTA = Ta vie, pas une option »} — \emph{Justification :} Observance 100 \% (pilulier hebdo, alarme téléphone). Contrôle TA mensuel infirmerie. \textbf{Cible : $< \SI{140/90}{\mmHg}$.} Arrêt tabac : se faire aider (CSAPA, substituts nicotiniques). \textbf{1 an sans tabac = risque AVC divisé par 2.}
\end{enumerate}

\medskip
\noindent \textbf{Signature secouristes :} Aïcha M. (3ème A), Kevin T. (3ème B), Stéphanie N. (3ème C) — Validé par M. \textit{Infirmier Chef}, Infirmerie Lycée Nkolbisson.
\end{exoresolu}

\courssub{Bilan de l'activité}
Cette activité d'intégration a permis de vérifier l'acquisition des compétences suivantes, conformément au programme officiel de 3ème :
\begin{itemize}
  \item \textbf{Savoirs mobilisés :} Structure et fonction des vaisseaux (artères coronaires, carotides) et du cœur ; physiopathologie de l'athérosclérose (plaques, sténose, thrombose, embolie) ; définition et classification des AVC (ischémique vs hémorragique) ; facteurs de risque cardiovasculaires (modifiables/non modifiables) ; signes cliniques d'alerte (méthode VITE) ; protocole de premiers secours (alerte 112, position, interdits, surveillance).
  \item \textbf{Savoir-faire évalués :} Analyse de documents hétérogènes (schéma médical, tableau clinique, courbe temporelle, affiche institutionnelle) ; rédaction d'une argumentation scientifique causale (chaîne logique plaque $\rightarrow$ ischémie $\rightarrow$ signes) ; application d'un protocole d'urgence avec justification physiologique de chaque geste ; synthèse sous forme de fiche technique professionnelle (fiche secouriste) ; élaboration de messages de prévention fondés sur l'évidence scientifique et adaptés au contexte local.
  \item \textbf{Savoir-être renforcés :} Réactivité face à l'urgence, sang-froid, esprit d'équipe, responsabilité citoyenne (alerte, secours), empathie, rigueur dans la transmission de l'information (heure, signes, antécédents).
\end{itemize}
L'ancrage dans un contexte camerounais réel (lycée de Yaoundé, habitudes alimentaires locales, numéros d'urgence 112/119, pathologies prévalentes HTA/Tabac/Diabète) a donné du sens aux apprentissages et préparé les élèves à agir concrètement comme premiers maillons de la chaîne de survie.

\newpage

\coursec{Méthodes et exercices résolus}

\coursec{Leçon 09 : Circulation sanguine}

\courssub{Les vaisseaux sanguins, siège de la circulation}

\begin{definition}[Artère]
Un vaisseau sanguin qui transporte le sang \textbf{du cœur vers les organes}. Sa paroi est \textbf{épaisse, musculeuse et élastique} pour résister à la forte pression du sang éjecté par le ventricule.
\end{definition}

\begin{definition}[Veine]
Un vaisseau sanguin qui ramène le sang \textbf{des organes vers le cœur}. Sa paroi est \textbf{mince}, peu musculeuse, et elle possède des \textbf{valvules} (clapets) qui empêchent le reflux du sang sous l'effet de la pesanteur.
\end{definition}

\begin{definition}[Capillaire sanguin]
Vaisseau \textbf{microscopique} (diamètre \SI{5}{\micro\metre} à \SI{10}{\micro\metre}) à paroi \textbf{très fine} (une seule couche de cellules endothéliales). C'est le \textbf{siège des échanges} entre le sang et le milieu interstitiel (apport en $O_2$, nutriments ; retrait de $CO_2$, déchets).
\end{definition}

\begin{methode}[Reconnaître une artère, une veine et un capillaire sur un document]
Pour identifier un vaisseau sanguin sur un schéma, une coupe microscopique ou une description, procéder par étapes :
\begin{enumerate}
\item \textbf{Observer le sens de la circulation} par rapport au cœur : l'artère part du cœur (centrifuge), la veine revient vers le cœur (centripète).
\item \textbf{Examiner la paroi} : l'artère a une paroi \textbf{épaisse et musculeuse} ; la veine a une paroi \textbf{mince} et présente souvent des \textbf{valvules} ; le capillaire a une paroi \textbf{unicellulaire}.
\item \textbf{Repérer le calibre et la visibilité} : artères et veines sont macroscopiques ; les capillaires sont microscopiques et forment un réseau dense (lit capillaire).
\item \textbf{Conclure} en nommant le vaisseau et en justifiant par au moins un critère morphologique ou topographique visible.
\end{enumerate}
\textbf{Réflexe à retenir :} Artère = départ du cœur + paroi épaisse ; Veine = retour au cœur + paroi mince + valvules ; Capillaire = échanges + paroi unicellulaire.
\end{methode}

\begin{exoresolu}[Identifier les vaisseaux d'un schéma]
Sur le schéma d'un membre inférieur, on observe trois vaisseaux : le vaisseau A transporte le sang du cœur vers la jambe et possède une paroi épaisse ; le vaisseau B ramène le sang de la jambe vers le cœur et présente des valvules ; le vaisseau C, microscopique, irrigue les muscles de la jambe. Identifier A, B et C en justifiant.
\tcblower
\textbf{Solution détaillée.}
\begin{enumerate}
\item \textbf{Vaisseau A} : il part du cœur vers un organe (la jambe) et sa paroi est épaisse et musculeuse. Seules les artères transportent le sang du cœur vers les organes sous forte pression, ce qui explique leur paroi épaisse. \textbf{A est donc une artère} (ici, l'artère fémorale).
\item \textbf{Vaisseau B} : il ramène le sang vers le cœur et possède des valvules, structures propres aux veines pour empêcher le reflux sanguin. \textbf{B est donc une veine} (ici, la veine saphène ou fémorale).
\item \textbf{Vaisseau C} : il est microscopique et irrigue directement les fibres musculaires : c'est au niveau des capillaires que s'effectuent les échanges gazeux et nutritifs. \textbf{C est donc un capillaire sanguin}.
\end{enumerate}
\textbf{Conclusion :} A est une artère, B une veine et C un capillaire ; l'ensemble constitue le réseau vasculaire assurant la perfusion de la jambe.
\end{exoresolu}

\begin{aretenir}[Critères de distinction des vaisseaux sanguins]
\begin{center}
\begin{tabularx}{\linewidth}{|l|X|X|X|}
\hline
\rowcolor{popblueL} \textbf{Critère} & \textbf{Artère} & \textbf{Veine} & \textbf{Capillaire} \\
\hline
\textbf{Sens} & Cœur $\rightarrow$ Organes & Organes $\rightarrow$ Cœur & Organes (échange) \\
\hline
\textbf{Paroi} & Épaisse, musculeuse, élastique & Mince, peu musculeuse, valvules & Unicellulaire (endothélium) \\
\hline
\textbf{Pression} & Forte & Faible & Très faible \\
\hline
\textbf{Visibilité} & Macroscopique & Macroscopique & Microscopique \\
\hline
\end{tabularx}
\end{center}
\end{aretenir}

\courssub{Le cœur, organe moteur de la circulation}

\begin{definition}[Cœur]
Organe musculaire creux, situé dans le médiastin, qui fonctionne comme une \textbf{pompe aspirante et refoulante}. Il est constitué de \textbf{myocarde} (muscle cardiaque strié involontaire) et comporte \textbf{quatre cavités} : deux oreillettes (supérieures) et deux ventricules (inférieures).
\end{definition}

\begin{definition}[Double circulation]
Le sang traverse le cœur \textbf{deux fois} pour faire le tour complet du corps :
\begin{itemize}
\item \textbf{Petite circulation (cœur-poumons-cœur)} : sang désoxygéné $\rightarrow$ ventricule droit $\rightarrow$ artère pulmonaire $\rightarrow$ poumons $\rightarrow$ veines pulmonaires $\rightarrow$ oreillette gauche.
\item \textbf{Grande circulation (cœur-organes-cœur)} : sang oxygéné $\rightarrow$ ventricule gauche $\rightarrow$ aorte $\rightarrow$ organes $\rightarrow$ veines caves $\rightarrow$ oreillette droite.
\end{itemize}
\end{definition}

\begin{methode}[Décrire le trajet du sang dans le cœur et l'organisme]
Pour décrire correctement la circulation sanguine :
\begin{enumerate}
\item \textbf{Situer les quatre cavités} : deux oreillettes (parois fines, rôle de réservoir) et deux ventricules (parois épaisses, rôle de pompe). La paroi du \cle{ventricule gauche} est la plus épaisse (\SI{10}{mm} à \SI{15}{mm}) car il propulse le sang dans toute la grande circulation.
\item \textbf{Décrire la petite circulation} : le sang pauvre en $O_2$ part du ventricule droit par l'\textbf{artère pulmonaire}, se charge en $O_2$ aux poumons, puis revient à l'oreillette gauche par les \textbf{veines pulmonaires} (sang riche en $O_2$).
\item \textbf{Décrire la grande circulation} : le sang riche en $O_2$ part du ventricule gauche par l'\textbf{aorte}, irrigue les organes (échange au niveau des capillaires), puis revient à l'oreillette droite par les \textbf{veines caves} (sang pauvre en $O_2$).
\item \textbf{Préciser le rôle des valvules cardiaques} : valvules auriculo-ventriculaires (mitrale à gauche, tricuspide à droite) et valvules sigmoïdes (aortique, pulmonaire) ; elles assurent l'unidirectionnalité du flux.
\item \textbf{Conclure} : le cœur assure un double circuit fermé, synchronisé par son système de conduction (nœud sinusal, nœud auriculo-ventriculaire, faisceau de His).
\end{enumerate}
\textbf{Piège classique (BEPC) :} L'artère pulmonaire transporte du sang \emph{pauvre} en $O_2$ et les veines pulmonaires du sang \emph{riche} en $O_2$. Le nom du vaisseau dépend du \textbf{sens de la circulation} (par rapport au cœur), \emph{pas} de la teneur en dioxygène.
\end{methode}

\begin{exoresolu}[Le trajet d'une molécule de dioxygène]
Décrire le trajet complet suivi par le sang depuis le ventricule droit jusqu'à un muscle de la jambe, en citant les cavités cardiaques et les vaisseaux empruntés.
\tcblower
\textbf{Solution détaillée.}
\begin{enumerate}
\item Le sang pauvre en $O_2$ est chassé du \textbf{ventricule droit} dans l'\textbf{artère pulmonaire} (tronc pulmonaire), qui se divise en deux pour aller aux poumons droit et gauche.
\item Dans les \textbf{capillaires pulmonaires}, le sang se charge en $O_2$ et élimine du $CO_2$ (hématose) : c'est la petite circulation.
\item Le sang oxygéné revient au cœur par les \textbf{quatre veines pulmonaires} qui débouchent dans l'\textbf{oreillette gauche}.
\item Le sang passe de l'oreillette gauche dans le \textbf{ventricule gauche} (ouverture de la valvule mitrale, fermeture de la valvule aortique pendant la diastole).
\item Le ventricule gauche, à la paroi très épaisse (myocarde puissant), se contracte (systole) et propulse le sang dans l'\textbf{aorte} (valvule aortique ouverte) : c'est le départ de la grande circulation.
\item De l'aorte, le sang passe dans les artères iliaques, puis l'\textbf{artère fémorale}, les artères du mollet, et enfin les \textbf{capillaires musculaires}, où le $O_2$ est cédé aux cellules (myocytes) pour la respiration cellulaire.
\end{enumerate}
\textbf{Conclusion :} Le trajet est : Ventricule D $\rightarrow$ Artère pulmonaire $\rightarrow$ Poumons $\rightarrow$ Veines pulmonaires $\rightarrow$ Oreillette G $\rightarrow$ Ventricule G $\rightarrow$ Aorte $\rightarrow$ Artères fémorales $\rightarrow$ Capillaires musculaires. Le sang traverse \textbf{deux fois le cœur} (deux pompes en série) avant d'atteindre la jambe.
\end{exoresolu}

\begin{experience}[TP N°7 : Dissection d'un cœur de mammifère (mouton ou bœuf)]
\textbf{Objectif :} Observer l'anatomie externe et interne du cœur, relier la morphologie des cavités à leur fonction.
\textbf{Matériel :} Cœur frais de mouton ou de bœuf, plateau de dissection, ciseaux, scalpel, pince, gants, blouse, règle graduée.
\textbf{Protocole :}
\begin{enumerate}
\item \textbf{Observation externe :} Repérer la forme conique, l'apex (pointe) à gauche, la base (grands vaisseaux) en haut. Identifier l'aorte, le tronc pulmonaire, les veines caves, les veines pulmonaires, les oreillettes (auricules). Différencier face antérieure (sillon interventriculaire antérieur, artères coronaires) et face postérieure.
\item \textbf{Ouverture des cavités :} Inciser la paroi du ventricule droit le long du sillon interventriculaire, puis ouvrir l'oreillette droite. Faire de même pour le côté gauche (ventricule gauche puis oreillette gauche).
\item \textbf{Mesures et observations internes :} Mesurer l'épaisseur des parois (oreillettes, ventricules droit et gauche). Observer les valvules (tricuspide, mitrale, sigmoïdes), les piliers et cordages, les piliers musculaires (trabéculae carneae), le septum interventriculaire.
\item \textbf{Repérage des orifices coronaires :} Au niveau de l'aorte, juste au-dessus de la valvule sigmoïde, repérer les ostia des artères coronaires (vascularisation du myocarde).
\end{enumerate}
\textbf{Consignes de sécurité :} Port des gants et blouse obligatoires. Ne jamais porter les mains au visage. Désinfection du matériel et lavage des mains impératifs après la séance. Élimination des déchets anatomiques selon le protocole du laboratoire.
\textbf{Exploitation attendue :} Tableau comparatif des épaisseurs, schémas légendés (coupes antérieure et postérieure), interprétation fonctionnelle (lien épaisseur/pression de travail).
\end{experience}

\begin{exoresolu}[Exploiter les résultats d'une dissection]
Au cours de la dissection d'un cœur de mouton, un élève mesure l'épaisseur des parois et obtient : oreillette droite : \SI{2}{mm} ; oreillette gauche : \SI{2}{mm} ; ventricule droit : \SI{5}{mm} ; ventricule gauche : \SI{12}{mm}.
\begin{enumerate}
\item Représenter ces résultats dans un tableau.
\item Expliquer les différences observées.
\end{enumerate}
\tcblower
\textbf{Solution détaillée.}

1) \textbf{Tableau des résultats :}
\begin{center}
\begin{tabularx}{\linewidth}{|l|X|X|X|X|}
\hline
\rowcolor{popblueL} \textbf{Cavité cardiaque} & Oreillette droite & Oreillette gauche & Ventricule droit & Ventricule gauche \\
\hline
\textbf{Épaisseur de la paroi (mm)} & 2 & 2 & 5 & 12 \\
\hline
\end{tabularx}
\end{center}

2) \textbf{Interprétation fonctionnelle :}
\begin{itemize}
\item \textbf{Oreillettes (\SI{2}{mm}) :} Parois très fines. Elles reçoivent le sang veineux (basse pression) et ne font que le pousser passivement dans les ventricules (contraction faible, « coup de pied de l'oreillette »).
\item \textbf{Ventricule droit (\SI{5}{mm}) :} Paroi modérément épaisse. Il propulse le sang dans la \textbf{petite circulation} (poumons) : circuit court, basse résistance vasculaire, pression systolique pulmonaire faible (\SI{\sim 25}{mmHg}).
\item \textbf{Ventricule gauche (\SI{12}{mm}) :} Paroi très épaisse (2 à 3 fois le VD). Il propulse le sang dans la \textbf{grande circulation} (aorte $\rightarrow$ tout l'organisme) : circuit long, forte résistance vasculaire périphérique, pression systolique systémique élevée (\SI{\sim 120}{mmHg}). Le travail mécanique (Pression $\times$ Volume) est maximal.
\end{itemize}
\textbf{Conclusion :} L'épaisseur de la paroi (masse myocardique) de chaque cavité est proportionnelle à la \textbf{pression} qu'elle doit générer pour vaincre la résistance du circuit vasculaire qu'elle alimente. Le cœur est une pompe double à pressions différenciées.
\end{exoresolu}

\begin{aretenir}[Architecture et fonctionnement du cœur]
\begin{itemize}
\item \textbf{Myocarde} : muscle strié involontaire, auto-excitable (automatisme), riche en mitochondries (aérobie strict).
\item \textbf{Valvules} : assurent l'unidirectionnalité (pas de reflux). Bruits du cœur : B1 (fermeture valvules AV), B2 (fermeture valvules sigmoïdes).
\item \textbf{Vascularisation propre :} Artères coronaires (naissance sur aorte) $\rightarrow$ capillaires myocardiques $\rightarrow$ veines cardiaques $\rightarrow$ sinus coronaire $\rightarrow$ oreillette droite. \textbf{Ischémie} = obstruction coronaire $\rightarrow$ infarctus.
\item \textbf{Double circulation en série} : Débit identique dans les deux circuits (loi de la continuité), pressions différentes.
\end{itemize}
\end{aretenir}

\courssub{Les hémorragies et les moyens de lutte}

\begin{definition}[Hémorragie]
Perte de sang importante hors du réseau vasculaire, due à la rupture d'un vaisseau. On distingue :
\begin{itemize}
\item \textbf{Hémorragie externe :} le sang s'écoule à l'air libre (plaie).
\item \textbf{Hémorragie interne :} le sang s'accumule dans une cavité naturelle (péritonée, thorax, crâne) ou un tissu (hématome), invisible de l'extérieur. Signes : pâleur, sueurs froides, pouls rapide et filant, hypotension, soif, anxiété.
\end{itemize}
\end{definition}

\begin{methode}[Conduite à tenir devant une hémorragie externe (Premiers secours)]
Face à une hémorragie externe visible, appliquer la règle \textbf{P.C.A.S.} (Protéger, Comprimer, Alerter, Surveiller) :
\begin{enumerate}
\item \textbf{Protéger} : Sécuriser les lieux (circulation, danger). \textbf{Se protéger soi-même} : gants à usage unique (ou sachet plastique propre) pour éviter tout contact avec le sang (risque VIH, Hépatites B/C).
\item \textbf{Comprimer} : \textbf{Compression manuelle directe} sur la plaie avec un linge propre (compresse, pagne, serviette). Appuyer fermement et \textbf{maintenir la pression sans relâcher}. Si un corps étranger est planté : \textbf{ne pas l'extraire}, comprimer autour.
\item \textbf{Allonger la victime} : Position horizontale, surélever le membre blessé si pas de fracture suspectée (favorise le retour veineux et réduit la pression artérielle locale).
\item \textbf{Alerter les secours} : Appeler le \textbf{119} (SAMU Cameroun) ou le centre de santé/hôpital le plus proche. Préciser : lieu, nombre de victimes, type de blessure, état de conscience.
\item \textbf{Surveiller} : Surveiller la respiration, le pouls, l'état de conscience. Rassurer la victime. Ne rien donner à boire ni à manger (risque anesthésie future).
\end{enumerate}
\textbf{Cas particulier :} Hémorragie interne $\rightarrow$ Pas de compression possible. Allonger, alerter d'urgence (119), surveiller, ne rien donner à boire.
\end{methode}

\begin{exoresolu}[Situation d'accident à Yaoundé]
Lors d'un match de football au quartier, un joueur se blesse à la jambe sur une tôle : le sang coule abondamment et de façon continue.
\begin{enumerate}
\item De quel type d'accident s'agit-il ?
\item Décrire, dans l'ordre, les gestes de premiers secours à effectuer par ses camarades.
\end{enumerate}
\tcblower
\textbf{Solution détaillée.}

1) Il s'agit d'une \textbf{hémorragie externe artérielle} (jet pulsatile synchrone du pouls) ou veineuse (flux continu). La description « sang coule abondamment et de façon continue » évoque une hémorragie veineuse ou artérielle de gros calibre. C'est une perte de sang à l'extérieur du corps par rupture de vaisseaux de la jambe.

2) \textbf{Gestes de premiers secours (Ordre strict) :}
\begin{enumerate}
\item \textbf{Protéger} : Écarter la foule, mettre des gants (ou sacs plastiques propres sur les mains) pour éviter la contamination croisée (VIH, Hépatites).
\item \textbf{Comprimer} : Appuyer \textbf{fermement et continûment} sur la plaie avec un tissu propre (pagne plié, serviette). Ne pas soulever le pansement pour regarder.
\item \textbf{Installer la victime} : Allongée sur le dos, jambe blessée surélevée (sur un sac, un ballon) si pas de douleur suspecte de fracture.
\item \textbf{Alerter} : Appeler le \textbf{119} (numéro d'urgence Cameroun) ou organiser l'évacuation vers l'hôpital le plus proche (Hôpital Central, Hôpital Général, Centre de Santé).
\item \textbf{Surveiller} : Maintenir la compression, parler à la victime, vérifier qu'elle respire normalement, noter l'heure de l'accident.
\end{enumerate}
\textbf{Conclusion :} La compression directe immédiate est le geste qui sauve la vie en limitant le volume sanguin perdu (choc hypovolémique).
\end{exoresolu}

\begin{savaistu}[Urgences au Cameroun : Numéros et structures]
\begin{itemize}
\item \textbf{119} : Numéro unique d'urgence médicale (SAMU) au Cameroun, gratuit, accessible 24h/24.
\item \textbf{112} : Numéro d'urgence européen, fonctionne aussi sur les réseaux mobiles camerounais (redirige vers 119).
\item \textbf{Structures de référence :} Hôpital Central de Yaoundé, Hôpital Laquintinie (Douala), Hôpital Général de Douala, Hôpitaux de district, Centres de Santé Intégrés (CSI).
\item \textbf{Croix-Rouge Camerounaise :} Formations aux premiers secours (PSC1), secours en cas de catastrophe.
\end{itemize}
\end{savaistu}

\courssub{Maladies cardiovasculaires : causes et prévention}

\begin{definition}[Athérosclérose]
Maladie chronique des artères de gros et moyen calibre caractérisée par le dépôt de \textbf{plaques d'athérome} (lipides, cholestérol LDL, cellules inflammatoires, tissu fibreux, calcium) sur la \textbf{intima} (paroi interne). Cela entraîne un \textbf{rétrécissement (sténose)} de la lumière artérielle, une perte d'élasticité (raidissement) et un risque de thrombose (caillot) ou de rupture.
\end{definition}

\begin{definition}[Hypertension artérielle (HTA)]
Élévation persistante de la pression artérielle au repos : \textbf{PAS (systolique) $\ge$ \SI{140}{mmHg}} et/ou \textbf{PAD (diastolique) $\ge$ \SI{90}{mmHg}} (seuil OMS). Facteur de risque majeur d'AVC, d'insuffisance cardiaque, d'insuffisance rénale et de rétinopathie.
\end{definition}

\begin{definition}[Infarctus du myocarde (Crise cardiaque)]
Nécrose (mort) d'une partie du muscle cardiaque due à l'obstruction brutale d'une artère coronaire (thrombus sur plaque d'athérome). Signe majeur : douleur thoracique intense, constrictive, irradiant au bras gauche/mâchoire, non calmée par le repos, sueurs froides, angoisse de mort. \textbf{Urgence vitale : 119.}
\end{definition}

\begin{methode}[Analyser un mode de vie et proposer des mesures de prévention cardiovasculaire]
\begin{enumerate}
\item \textbf{Identifier les facteurs de risque} dans la situation : \textbf{Modifiables} (tabagisme, alimentation déséquilibrée riche en graisses saturées/sel, sédentarité, surpoids/obésité, alcool, stress, HTA non traitée, diabète, dyslipidémie) et \textbf{Non modifiables} (âge, sexe masculin, hérédité).
\item \textbf{Expliquer le mécanisme pathophysiologique} : Excès de LDL-cholestérol $\rightarrow$ pénétration dans l'intima $\rightarrow$ oxydation $\rightarrow$ inflammation $\rightarrow$ plaque d'athérome $\rightarrow$ sténose/rupture $\rightarrow$ ischémie (angor, infarctus, AVC ischémique). HTA $\rightarrow$ surcharge de pression $\rightarrow$ hypertrophie ventriculaire gauche $\rightarrow$ insuffisance cardiaque + fragilisation parois artérielles.
\item \textbf{Proposer des mesures de prévention primaire (avant la maladie) et secondaire (après événement)} adaptées : Régime méditerranéen / DASH (fruits, légumes, poisson, huiles végétales, peu de sel/graisses saturées), activité physique \SI{\ge 30}{min}/jour (marche rapide), sevrage tabagique total, limitation alcool, prise en charge médicale de l'HTA/diabète/dyslipidémie, gestion du stress.
\item \textbf{Conclure} : 80\% des maladies cardiovasculaires prématurées sont évitables par l'hygiène de vie (OMS).
\end{enumerate}
\end{methode}

\begin{exoresolu}[Conseiller un patient hypertendu]
Monsieur N., 52 ans, commerçant à Douala, fume un paquet de cigarettes par jour, consomme régulièrement des plats très gras (ndolé, poulet DG, fritures) et très salés, ne pratique aucune activité physique (voiture pour 200 m) et mesure une tension artérielle de \SI{160/100}{mmHg} (valeurs cibles : \textless \SI{140/90}{mmHg}).
\begin{enumerate}
\item Relever les facteurs de risque de maladie cardiovasculaire présents chez ce patient.
\item Lui proposer quatre conseils d'hygiène de vie en les justifiant physiologiquement.
\end{enumerate}
\tcblower
\textbf{Solution détaillée.}

1) \textbf{Facteurs de risque relevés (majoritairement modifiables) :}
\begin{itemize}
\item \textbf{Tabagisme actif majeur} (20 cig/jour) : Endommage l'endothélium, favorise l'athérothrombose, vasoconstriction nicotinique, augmentation $CO$ sanguin.
\item \textbf{Dyslipidémie alimentaire} : Excès d'acides gras saturés (huile de palme non traitée, viande grasse) $\rightarrow$ augmentation LDL-cholestérol $\rightarrow$ athérosclérose accélérée.
\item \textbf{Excès de sel (NaCl)} : Rétention hydrosodée $\rightarrow$ augmentation volémie $\rightarrow$ élévation de la Pression Artérielle (PA).
\item \textbf{Sédentarité / Inactivité physique} : Déconditionnement cardiovasculaire, prise de poids, résistance à l'insuline.
\item \textbf{Hypertension artérielle confirmée} (\SI{160/100}{mmHg}) : Lésion d'organe cœur (hypertrophie VG), reins, yeux, artères. Facteur de risque \textbf{majeur} d'AVC et d'infarctus.
\item \textbf{Âge (52 ans) et sexe (Homme) :} Facteurs non modifiables.
\end{itemize}

2) \textbf{Conseils d'hygiène de vie justifiés :}
\begin{enumerate}
\item \textbf{Arrêt total du tabac} (sevrage aidé si nécessaire) : \textit{Justification :} Supprime l'agression endothéliale quotidienne, réduit le risque d'infarctus de 50\% à 1 an, améliore le goût et la capacité respiratoire.
\item \textbf{Rééquilibrage alimentaire type DASH/Méditerranéen} : Réduire huile de palme/arachide (remplacer par huile rouge non chauffée, huile de soja, arachide), limiter sel de cuisine/bouillons cubes (maggi), augmenter légumes (feuilles : eru, folong, koki ; fruits : papaye, ananas, banane plantain), poisson/viandes maigres. \textit{Justification :} Baisse LDL-c, apport potassium (effet vasodilatateur), contrôle pondéral, baisse PA.
\item \textbf{Activité physique régulière} : Marche rapide \SI{30}{min} $\times$ 5 jours/semaine (ex: trajet marché, terrain de sport). \textit{Justification :} Entraîne le myocarde (baisse FC repos), améliore sensibilité à l'insuline, favorise perte de poids, baisse PA diastolique.
\item \textbf{Suivi médical mensuel} : Contrôle PA, observance du traitement antihypertenseur si prescrit (ex: IEC, Sartan, Diurétique, Calcium-antagoniste), bilan lipidique, glycémie, créatininémie. \textit{Justification :} L'HTA est asymptomatique ; seul le contrôle régulier prévient les complications.
\end{enumerate}
\textbf{Conclusion :} La modification du mode de vie est le \textbf{traitement de première intention} de l'HTA stade 1 et le socle de la prévention cardiovasculaire. Monsieur N. peut normaliser sa PA et éviter un AVC ou un infarctus.
\end{exoresolu}

\begin{aretenir}[Facteurs de risque cardiovasculaires (Classification OMS)]
\begin{center}
\begin{tabularx}{\linewidth}{|l|X|}
\hline
\rowcolor{popblueL} \textbf{Modifiables (Agir dessus)} & \textbf{Non modifiables (Surveiller)} \\
\hline
Tabagisme & Âge (\textgreater 50 ans H, \textgreater 60 ans F) \\
Hypertension artérielle & Sexe (Homme \textgreater Femme avant ménopause) \\
Dyslipidémie (LDL-c haut, HDL-c bas) & Antécédents familiaux précoces (Infarctus \textless 55 ans père/frère, \textless 65 ans mère/sœur) \\
Diabète de type 2 & Ethnie (prédisposition génétique) \\
Surpoids / Obésité abdominale (Tour de taille H \textgreater \SI{94}{cm}, F \textgreater \SI{80}{cm}) & \\
Sédentarité & \\
Alcool excessif & \\
Stress psychosocial & \\
\hline
\end{tabularx}
\end{center}
\end{aretenir}

\begin{savaistu}[Maladies cardiovasculaires au Cameroun : Enquête STEPS 2013-2014]
\begin{itemize}
\item Prévalence de l'HTA : \textbf{\SI{\sim 30}{\%}} des adultes (15-69 ans), souvent méconnue (\textless 20\% traités, \textless 10\% contrôlés).
\item Facteurs de risque combinés : \SI{12}{\%} de la population cumule 3 à 5 facteurs de risque.
\item Charge économique : Coût direct (hospitalisations, dialyse) et indirect (perte de productivité, décès prématuré homme 40-60 ans) énorme pour les ménages et le système de santé.
\item \textbf{Plan d'action national :} Lutte anti-tabac (loi 2006, paquet neutre 2023), promotion activité physique (Journée nationale du sport), réduction sel/huile dans la restauration collective, renforcement prise en charge au niveau des CSI (protocoles simplifiés).
\end{itemize}
\end{savaistu}

\courssub{Les accidents vasculaires cérébraux (AVC) et les premiers secours}

\begin{definition}[Accident Vasculaire Cérébral (AVC)]
Survenue brutale d'un déficit neurologique focal (ou global) d'origine vasculaire, durant plus de 24h (ou entraînant la mort). Deux types principaux :
\begin{itemize}
\item \textbf{AVC Ischémique (\SI{\sim 80}{\%}) :} Obstruction d'une artère cérébrale par un thrombus (thrombose sur athérome) ou un embole (origine cardiaque : fibrillation atriale, ou artère carotide). $\rightarrow$ Ischémie $\rightarrow$ Nécrose (infarctus cérébral) si non revascularisé rapidement.
\item \textbf{AVC Hémorragique (\SI{\sim 20}{\%}) :} Rupture d'une artère cérébrale (souvent sur HTA mal contrôlée, anévrisme, malformation artério-veineuse) $\rightarrow$ Hématome intra-cérébral $\rightarrow$ Compression + toxicité du sang pour les neurones.
\end{itemize}
\textbf{AIT (Accident Ischémique Transitoire) :} Signes identiques mais \textbf{réversibles en \textless 1h (souvent \textless 15 min)}. \textbf{Signal d'alarme majeur} : risque d'AVC complet à \SI{10}{\%} dans les 48h. \textbf{Urgent : 119.}
\end{definition}

\begin{methode}[Conduite à tenir devant un AVC (Règle V.I.T.E. / F.A.S.T.)]
\begin{enumerate}
\item \textbf{Reconnaître les signes d'alerte (Soudains) :}
    \begin{itemize}
    \item \textbf{V} - \textbf{Visage} : Bouche tordue, impossibilité de sourire (parésie faciale).
    \item \textbf{I} - \textbf{Incapacité motrice} : Bras ou jambe paralysé(e) d'un côté (hémiparésie), objet qui tombe.
    \item \textbf{T} - \textbf{Trouble de la parole} : Incompréhensible (dysarthrie), mots inexistants (aphasie), incapacité à comprendre.
    \item \textbf{E} - \textbf{Extrême urgence} : Maux de tête violents soudains (céphalées « coup de massue »), perte de vision, vertiges, perte de conscience.
    \end{itemize}
\item \textbf{Alerter IMMÉDIATEMENT le 119 (SAMU) ou faire transporter d'urgence à l'hôpital (Neurochirurgie/Neurologie).} \textbf{Noter l'heure exacte du début des signes} (décide de l'éligibilité à la thrombolyse : fenêtre \textless 4h30).
\item \textbf{Installer la victime :}
    \begin{itemize}
    \item Consciente : Demi-assise (tronc \SI{30}{\degree}), tête légèrement surélevée, rassurer, ne pas bouger.
    \item Inconsciente mais respirant : \textbf{Position Latérale de Sécurité (PLS)} côté paralysé \textbf{en bas} (pour éviter l'inhalation de salives/vomissements et protéger la voie aérienne).
    \end{itemize}
\item \textbf{Interdits absolus :} Ne rien donner à boire ni à manger (risque fausse route pneumonique), ne pas faire marcher, ne pas injecter de médicament (sauf ordre médical), ne pas mettre d'eau froide sur le visage.
\item \textbf{Surveiller} : Respiration, pouls, niveau de conscience jusqu'à l'arrivée des secours. Transmettre l'heure de début et les antécédents (HTA, diabète, cardiaque, anticoagulants).
\end{enumerate}
\textbf{Règle d'or :} « \textbf{Le temps, c'est du cerveau} » (\SI{1,9}{million} de neurones perdus par minute d'ischémie). La rapidité de la prise en charge hospitalière (Unité Neuro-Vasculaire - UNV) conditionne la survie et la récupération.
\end{methode}

\begin{exoresolu}[Urgence au marché]
Une vendeuse de 60 ans s'effondre brusquement au marché : sa bouche est déviée vers la gauche, elle ne peut plus articuler correctement (parole confuse) et son bras droit reste immobile le long du corps.
\begin{enumerate}
\item Quel accident a probablement subi cette dame ? Justifier par les signes cliniques.
\item Décrire les gestes de premiers secours à accomplir par les témoins.
\end{enumerate}
\tcblower
\textbf{Solution détaillée.}

1) Cette dame a probablement subi un \textbf{Accident Vasculaire Cérébral (AVC) ischémique ou hémorragique de l'hémisphère gauche}.
\textbf{Justification clinique :}
\begin{itemize}
\item \textbf{Brutalité} de l'installation (s'effondre « brusquement »).
\item \textbf{Syndrome hémisphérique gauche} : Atteinte controlatérale (côté droit du corps) $\rightarrow$ \textbf{Hémiplégie droite} (bras droit immobile) + \textbf{Hémiplégie faciale droite} (bouche déviée vers le côté sain = gauche) + \textbf{Aphasie/Dysarthrie} (trouble de la parole, aire de Broca dans hémisphère gauche chez le droitier).
\end{itemize}
Ces signes traduisent la souffrance/mort neuronale dans le territoire de l'artère cérébrale moyenne gauche (sylvienne).

2) \textbf{Gestes de premiers secours (Ordre vital) :}
\begin{enumerate}
\item \textbf{APPELER LE 119 IMMÉDIATEMENT} (ou faire appeler pendant qu'on s'occupe de la victime). Dire : « AVC probable, femme 60 ans, hémiplégie droite, trouble parole, heure de début : [Heure exacte] ». C'est l'action \textbf{la plus importante}.
\item \textbf{Installer en PLS côté droit (côté paralysé) EN BAS} car elle est « effondrée » (altération de la vigilance possible). Cela libère les voies aériennes.
\item \textbf{Desserrer} tout vêtement serrant (collier, foulard, ceinture).
\item \textbf{Ne RIEN donner à boire ni à manger} (dysphagie fréquente = fausse route).
\item \textbf{Surveiller la respiration} (fréquence, amplitude) et l'état de conscience (répond-elle à son nom ? serre-t-elle la main gauche ?).
\item \textbf{Noter l'heure précise} des premiers signes pour l'équipe du SAMU/UNV (critère thrombolyse).
\item \textbf{Rassembler} les ordonnances/médicaments si accessibles (anticoagulants ? antihypertenseurs ?).
\end{enumerate}
\textbf{Conclusion :} La reconnaissance du syndrome VITE et l'appel précoce au 119 sont les seuls gestes qui permettent d'accéder aux traitements de réouverture artérielle (thrombolyse, thrombectomie) et de limiter le handicap définitif.
\end{exoresolu}

\begin{aretenir}[Prise en charge de l'AVC : La chaîne de survie]
\begin{enumerate}
\item \textbf{Témoin} : Reconnaître (VITE) $\rightarrow$ Appeler 119 $\rightarrow$ PLS $\rightarrow$ Surveiller.
\item \textbf{SAMU / SMUR} : Diagnostic clinique $\rightarrow$ Transport vers UNV (Unité Neuro-Vasculaire) certifiée $\rightarrow$ Pré-alerte hôpital.
\item \textbf{Hôpital (UNV) :} IRM/Scanner cérébral URGENT (différencier ischémique/hémorragique) $\rightarrow$ \textbf{Thrombolyse IV (alteplase) si ischémique \textless 4h30} $\rightarrow$ \textbf{Thrombectomie mécanique si gros vaisseau \textless 6h-24h} $\rightarrow$ Surveillance neuro-vitale, prévention complications (pneumonie, phlébite, escarres) $\rightarrow$ Rééducation précoce.
\end{enumerate}
\end{aretenir}

\courssub{Compétence transversale : Rédiger une réponse argumentée au BEPC}

\begin{methode}[Rédiger et justifier une démarche ou une conclusion (Examen BEPC)]
\begin{enumerate}
\item \textbf{Analyser le sujet :} Lire 2 fois. Souligner les \textbf{données fournies} (observations, mesures, schémas, signes cliniques) et la \textbf{question posée}.
\item \textbf{Structurer la réponse :} \textbf{1. Constat / Analyse des données} (Ce que je vois/lis). \textbf{2. Mobilisation des connaissances} (Ce que je sais : définitions, mécanismes, lois). \textbf{3. Raisonnement / Argumentation} (Lien logique Données $\rightarrow$ Connaissances $\rightarrow$ Explication). \textbf{4. Conclusion claire} (Réponse directe à la question).
\item \textbf{Vocabulaire scientifique EXIGÉ :} Artère, veine, capillaire, oreillette, ventricule, aorte, artère pulmonaire, veines caves, veines pulmonaires, hémorragie, compression, AVC, athérosclérose, hypertension, valvules, myocarde, petite/grande circulation. \textbf{Pas de langage familier} (« le cœur pompe », « le sang va dans les veines »).
\item \textbf{Justifier chaque affirmation} : « Car... », « Puisque... », « Ce qui s'explique par... ». Une affirmation sans justification = 0 point ou demi-point.
\item \textbf{Soigner la présentation :} Phrases complètes, écriture lisible, schémas légendés si demandés, unités SI (\SI{}{mm}, \SI{}{mmHg}, \SI{}{min}).
\end{enumerate}
\end{methode}

\begin{exoresolu}[Expliquer un gonflement des jambes (Œdème de déclivité)]
Chez certaines personnes âgées qui restent debout longtemps, on observe un gonflement des jambes en fin de journée. On sait que le sang remonte difficilement des jambes vers le cœur. À l'aide de tes connaissances sur les veines, explique ce phénomène et propose un moyen de le limiter.
\tcblower
\textbf{Solution détaillée (Modèle BEPC).}

\textbf{1. Analyse de la situation :}
Le sang des membres inférieurs doit remonter vers le cœur \textbf{contre la pesanteur}. Ce retour veineux (flux centripète) est assuré par les \textbf{veines}, vaisseaux à paroi mince et distensible, munis de \textbf{valvules} anti-reflux.

\textbf{2. Explication physiopathologique :}
Chez la personne âgée station debout prolongée (orthostatisme) :
\begin{itemize}
\item La \textbf{pression hydrostatique} dans les veines des jambes est très élevée (colonne de sang de \SI{\sim 1}{m} de haut).
\item Les \textbf{valvules veineuses deviennent incompétentes} (vieillissement, dilatation veineuse) : elles ne ferment plus étanche.
\item Le sang \textbf{reflue et stagne} dans le réseau veineux superficiel et profond des jambes (insuffisance veineuse).
\item L'augmentation de la \textbf{pression veineuse capillaire} force la sortie du plasma (eau + électrolytes) vers l'espace interstitiel $\rightarrow$ \textbf{Œdème de déclivité} (gonflement mou, indolore, en « godet », siégeant aux chevilles/pieds, régressant au décubitus).
\end{itemize}

\textbf{3. Moyens de limitation (Hygiène veineuse) :}
\begin{itemize}
\item \textbf{Activité musculaire régulière (Marche) :} La contraction des muscles du mollet (pompé veineux / « cœur périphérique ») comprime les veines profondes $\rightarrow$ chasse le sang vers le haut (valvules s'ouvrent) $\rightarrow$ empêche la stase.
\item \textbf{Surélévation des jambes au repos} (pied du lit surélevé \SI{10}{cm}, ou jambes sur coussin) : La gravité aide le retour veineux $\rightarrow$ résorption de l'œdème.
\item \textbf{Bas de contention (contention élastique) :} Pression dégressive cheville $\rightarrow$ cuisse, augmente la vitesse veineuse, réduit la distension.
\item \textbf{Hygiène de vie :} Perte de poids si obésité, éviter la chaleur excessive (dilatation), hydratation.
\end{itemize}

\textbf{Conclusion :} Le gonflement des jambes (œdème veineux) résulte de la \textbf{stase sanguine} dans les veines des membres inférieurs due à l'incompétence valvulaire et à l'orthostatisme prolongé. La \textbf{marche} et la \textbf{surélévation} activent le retour veineux et préviennent ce trouble.
\end{exoresolu}

\begin{attention}[Les 5 erreurs qui coûtent des points au BEPC]
\begin{enumerate}
\item \textbf{Confondre artère et veine sur le critère du sang} : Écrire « l'artère transporte du sang riche en $O_2$ ». \textbf{FAUX.} L'artère pulmonaire transporte du sang pauvre. \textbf{Critère unique :} Sens par rapport au cœur (Artère = Part, Veine = Revient).
\item \textbf{Inverser les cavités cardiaques (Gauche/Droite)} : Dire « le sang oxygéné arrive dans l'oreillette droite » ou « le ventricule droit envoie le sang dans l'aorte ». \textbf{Repère :} Poumon = Gauche (Oxygéné) ; Corps = Droit (Désoxygéné) $\rightarrow$ \textbf{Gauche = Riche en $O_2$, Droit = Pauvre en $O_2$}.
\item \textbf{Identifier un vaisseau sans justifier} : Écrire « C'est une veine » sans citer « paroi mince » ou « valvules » ou « revient au cœur ». \textbf{0 point.} Toujours écrire : « C'est une veine \textbf{car} sa paroi est mince et elle possède des valvules. »
\item \textbf{Désordre dans les gestes de secours} :
    \begin{itemize}
    \item Hémorragie : \textbf{Comprimer AVANT d'alerter} (si seul témoin, compresser 1 min puis alerter, ou alerter en compressant). Ne pas faire de garrot (sauf membre arraché/écrasé).
    \item AVC : \textbf{Alerter (119) EN PREMIER}. Ne pas perdre de temps à prendre la tension ou donner du sucre.
    \end{itemize}
\item \textbf{Vocabulaire imprécis / Phrases télégraphiques} : « Cœur pompe sang partout » au lieu de « Le ventricule gauche éjecte le sang dans l'aorte pour la grande circulation ». « Vaisseau bouché » au lieu de « Artère coronaire obstruée par un thrombus sur plaque d'athérome ». \textbf{Le BEPC sanctionne le manque de rigueur terminologique.}
\end{enumerate}
\end{attention}

\begin{exercice}[Entraînement BEPC - Circulation sanguine]
\begin{enumerate}
\item \textbf{Schéma annoté :} Reproduire le schéma simplifié du cœur (4 cavités, 4 gros vaisseaux, valvules). Légender : Oreillette D/G, Ventricule D/G, Veines caves, Artère pulmonaire, Veines pulmonaires, Aorte. Indiquer par des flèches rouges le sang riche en $O_2$, bleues le sang pauvre en $O_2$.
\item \textbf{QCM :} Une artère se distingue d'une veine car :
    \begin{enumerate}
    \item Elle transporte toujours du sang oxygéné.
    \item Sa paroi est plus fine.
    \item Elle part du cœur.
    \item Elle n'a pas de valvules.
    \end{enumerate}
\item \textbf{Mise en situation :} Un cycliste chute et présente une plaie large au cuisse avec un jet de sang pulsatile.
    \begin{enumerate}
    \item Quel type d'hémorragie ? Justifier.
    \item Décrire la conduite à tenir précise (ordres, gestes, numéro d'urgence).
    \end{enumerate}
\item \textbf{Argumentation :} « L'hypertension artérielle est un facteur de risque majeur d'AVC. » Expliquer ce lien de causalité en utilisant les termes : paroi artérielle, athérosclérose, pression, rupture/obstruction, neurones.
\item \textbf{Calcul :} Au cours d'une dissection, on mesure le diamètre interne de l'aorte (\SI{25}{mm}) et de l'artère pulmonaire (\SI{22}{mm}). Sachant que le débit cardiaque est de \SI{5}{L/min} au repos, comparer la vitesse moyenne du sang dans ces deux artères (formule $v = Q/S$, $S = \pi r^2$). Interpréter la différence.
\end{enumerate}
\end{exercice}

\begin{corrige}[Exercice n°1 - Entraînement BEPC]
\begin{enumerate}
\item \textbf{Schéma :} (Description textuelle attendue) Cœur vu de face. Base en haut. Oreillette D (supérieure droite) recevant Veines Caves S/I (bleu). Ventricule D (inférieur droit) envoyant Tronc Pulmonaire (bleu) $\rightarrow$ Artères Pulmonaires D/G. Oreillette G (supérieure gauche) recevant 4 Veines Pulmonaires (rouge). Ventricule G (inférieur gauche, paroi épaisse) envoyant Aorte (rouge) $\rightarrow$ Crosse $\rightarrow$ Artères coronaires. Valvules : Tricuspide (OD-VD), Pulmonaire (VD-Art Pulm), Mitrale (OG-VG), Aortique (VG-Aorte). Flèches : Bleu (VD $\rightarrow$ Poumons $\rightarrow$ OG), Rouge (VG $\rightarrow$ Aorte $\rightarrow$ Corps $\rightarrow$ OD).
\item \textbf{QCM :} Bonne réponse : \textbf{3} (Elle part du cœur). Justification : 1 faux (artère pulmonaire), 2 faux (paroi plus épaisse), 4 faux (valvules sigmoïdes à la base).
\item \textbf{Hémorragie :}
    \begin{enumerate}
    \item \textbf{Hémorragie externe artérielle.} Justification : Sang rouge vif, jet pulsatile (synchrone systole cardiaque), débit important $\rightarrow$ rupture artère fémorale superficielle ou profonde.
    \item \textbf{Conduite :} 1. Protéger (gants). 2. \textbf{Compression manuelle FORTE} point de compression (racine cuisse / pli de l'aine) AU-DESSUS de la plaie + compression directe plaie. 3. Allonger, surélever jambe. 4. \textbf{Appeler 119} (ou faire appeler). 5. Surveiller, couvrir, ne rien donner. \textbf{Pas de garrot} sauf si compression inefficace et vie menacée (noter heure pose).
    \end{enumerate}
\item \textbf{Argumentation :} L'HTA chronique $\rightarrow$ hypertrophie de la média artérielle + lésion endothéliale $\rightarrow$ favorise dépôt LDL $\rightarrow$ \textbf{athérosclérose accélérée} (artères carotides, cérébrales). Parois fragilisées (anévrismes micro-artériels - lipohyalinose). \textbf{Deux mécanismes d'AVC :}
    \begin{itemize}
    \item \textbf{Ischémique :} Plaque d'athérome instable $\rightarrow$ rupture $\rightarrow$ thrombose $\rightarrow$ obstruction artère cérébrale $\rightarrow$ ischémie nécrotique neurones.
    \item \textbf{Hémorragique :} Rupture micro-anévrisme (maladie des petits vaisseaux) ou artère de gros calibre $\rightarrow$ hématome intra-cérébral $\rightarrow$ compression + toxicité hémoglobine $\rightarrow$ mort neuronale.
    \end{itemize}
    Conclusion : Contrôler la PA (\textless 140/90) prévient les deux types d'AVC.
\item \textbf{Calcul :}
    \begin{align*}
    Q &= \SI{5}{L/min} = \SI{5000}{cm^3/min} = \SI{83.3}{cm^3/s} \quad (\text{débit identique dans les deux circuits, loi de la continuité}) \\
    S_{\text{aorte}} &= \pi \times (\SI{1.25}{cm})^2 = \SI{4.91}{cm^2} \\
    v_{\text{aorte}} &= \frac{83.3}{4.91} \approx \SI{17}{cm/s} \\
    S_{\text{art pulm}} &= \pi \times (\SI{1.1}{cm})^2 = \SI{3.80}{cm^2} \\
    v_{\text{art pulm}} &= \frac{83.3}{3.80} \approx \SI{22}{cm/s}
    \end{align*}
    \textbf{Interprétation :} La vitesse est plus élevée dans l'artère pulmonaire car sa section est plus faible (pour une même pression moyenne plus basse, le circuit pulmonaire a une résistance plus faible, d'où un calibre légèrement plus petit pour un débit égal). \textit{Note : En réalité, l'aorte a souvent un débit légèrement supérieur (shunt bronchique), mais au programme 3ème, on considère débit égal.}
\end{enumerate}
\end{corrige}

\newpage

\coursec{Exercices}

\coursec{Circulation sanguine : Entraînement et préparation au BEPC}

\courssub{Je vérifie mes connaissances}

\begin{exercice}[QCM : les vaisseaux et la circulation]
Pour chaque question, choisis la (ou les) bonne(s) réponse(s).
\begin{enumerate}
\item Les artères sont des vaisseaux qui :
\begin{enumerate}
\item ramènent le sang vers le c\oe ur ;
\item transportent le sang du c\oe ur vers les organes ;
\item ont une paroi épaisse et élastique ;
\item contiennent toujours du sang riche en dioxygène.
\end{enumerate}
\item Les capillaires :
\begin{enumerate}
\item ont une paroi très fine ;
\item sont le siège des échanges entre le sang et les organes ;
\item ont un diamètre très grand ;
\item relient les artérioles aux veinules.
\end{enumerate}
\item La petite circulation (pulmonaire) :
\begin{enumerate}
\item va du c\oe ur vers les poumons ;
\item va du c\oe ur vers tout le corps ;
\item part du ventricule droit ;
\item part du ventricule gauche.
\end{enumerate}
\item Face à une personne victime d'un AVC, il faut :
\begin{enumerate}
\item lui donner à boire ;
\item alerter immédiatement les secours (SAMU 112 ou 119) ;
\item la coucher sur le dos, la tête légèrement surélevée (30°) ;
\item la faire marcher pour la réveiller.
\end{enumerate}
\end{enumerate}
\end{exercice}

\begin{exercice}[Vrai ou faux, en justifiant]
Réponds par \textbf{vrai} ou \textbf{faux}, puis justifie ta réponse en une ou deux phrases.
\begin{enumerate}
\item Le c\oe ur est un muscle creux qui se contracte et se relâche automatiquement.
\item Les veines possèdent des valvules qui empêchent le sang de revenir en arrière.
\item Le ventricule gauche a une paroi plus mince que le ventricule droit.
\item Une hémorragie externe grave peut entraîner la mort si l'on n'intervient pas rapidement.
\end{enumerate}
\end{exercice}

\begin{exercice}[Définitions]
Définis en une phrase claire et complète chacun des termes suivants : \cle{artère}, \cle{veine}, \cle{capillaire}, \cle{hémorragie}, \cle{accident vasculaire cérébral (AVC)}.
\end{exercice}

\begin{exercice}[Calculs directs : fréquence cardiaque et débit cardiaque]
Au repos, le c\oe ur d'un élève bat 72 fois par minute et éjecte à chaque battement \SI{70}{mL} de sang (volume d'éjection systolique).
\begin{enumerate}
\item Calcule le volume de sang éjecté par le c\oe ur en une minute (débit cardiaque), en litres par minute.
\item Pendant un effort, la fréquence cardiaque passe à 150 battements par minute, le volume éjecté par battement restant de \SI{70}{mL}. Calcule le nouveau débit cardiaque en litres par minute.
\item Compare les deux résultats et conclus sur le rôle du c\oe ur pendant l'effort.
\end{enumerate}
\end{exercice}

\courssub{Je m'entraîne}

\begin{exercice}[L'expérience de Harvey : sens de la circulation veineuse]
Un élève réalise l'expérience suivante : il pose un garrot (ligature) serré autour de son bras, au niveau du milieu du bras. Au bout de quelques instants, il observe que les veines de l'avant-bras gonflent et deviennent bien visibles, avec des renflements espacés.
\begin{enumerate}
\item Explique pourquoi les veines gonflent en aval du garrot (vers la main).
\item En appuyant avec un doigt sur une veine gonflée et en poussant le sang vers le coude, la veine se vide jusqu'à un renflement, puis se remplit à nouveau par le côté main. Que montre cette observation sur le sens de circulation du sang dans les veines ?
\item Quel rôle jouent les valvules des veines observées sous forme de renflements ?
\end{enumerate}
\end{exercice}

\begin{exercice}[Structure et fonction des cavités cardiaques]
Lors de la dissection d'un c\oe ur de mouton, un groupe d'élèves mesure l'épaisseur des parois des quatre cavités et obtient les résultats suivants :
\begin{center}
\begin{tabularx}{\linewidth}{|l|X|X|X|X|}\hline
\rowcolor{popblueL}
\textbf{Cavité} & Oreillette droite & Oreillette gauche & Ventricule droit & Ventricule gauche \\ \hline
\textbf{Épaisseur (mm)} & 2 & 3 & 4 & 10 \\ \hline
\end{tabularx}
\end{center}
\begin{enumerate}
\item Classe les quatre cavités par épaisseur croissante de leur paroi.
\item Rappelle le rôle de chaque cavité dans la circulation du sang (provenance et destination du sang).
\item Établis une relation entre l'épaisseur de la paroi d'une cavité et le travail mécanique qu'elle doit fournir (pression à générer).
\end{enumerate}
\end{exercice}

\begin{exercice}[Premiers secours : hémorragie externe grave]
Sur la route de Yaoundé, un motard est victime d'un accident. Allongé au sol, il présente une plaie à la cuisse d'où le sang s'écoule en jets rouges vifs, synchrones avec les battements du c\oe ur, de façon abondante.
\begin{enumerate}
\item De quel type d'hémorragie s'agit-il ? Justifie ta réponse par deux arguments.
\item Cite les gestes de premiers secours à effectuer, dans l'ordre chronologique strict, en précisant l'action pour chaque étape.
\item Explique pourquoi il faut agir très vite face à ce type d'hémorragie.
\end{enumerate}
\end{exercice}

\begin{exercice}[Fréquence cardiaque et effort : analyse graphique]
On mesure la fréquence cardiaque d'un élève de 3\textsuperscript{e} avant, pendant et après une course de 5 minutes :
\begin{center}
\begin{tabularx}{\linewidth}{|l|X|X|X|X|X|X|}\hline
\rowcolor{popblueL}
\textbf{Temps (min)} & 0 & 2 & 5 & 7 & 10 & 15 \\ \hline
\textbf{Fréquence (battements/min)} & 70 & 120 & 150 & 110 & 85 & 72 \\ \hline
\end{tabularx}
\end{center}
L'effort commence à $t = 0$ min et s'arrête à $t = 5$ min.
\begin{enumerate}
\item Trace la courbe de la fréquence cardiaque en fonction du temps (échelle suggérée : 1 cm pour 2 min en abscisse ; 1 cm pour 20 battements/min en ordonnée).
\item Décris précisément les variations de la fréquence cardiaque au cours de l'expérience en distinguant la phase d'effort et la phase de récupération.
\item Explique physiologiquement pourquoi la fréquence cardiaque augmente pendant l'effort.
\item Calcule le temps de récupération (durée entre la fin de l'effort et le retour à la valeur de repos initiale).
\end{enumerate}
\end{exercice}

\courssub{Je me prépare au BEPC}

\begin{exercice}[Le c\oe ur et la double circulation]
\textbf{Partie A : Anatomie du c\oe ur.}\\
Le schéma ci-dessous représente une coupe frontale schématique du c\oe ur humain. Les structures sont désignées par les lettres A à F.
\begin{popfigure}
\begin{tikzpicture}[scale=1.1, every node/.style={font=\small}]
% Couleurs : Bleu = Sang désoxygéné (Cœur droit), Orange = Sang oxygéné (Cœur gauche)
% Ventricules
\draw[thick, popblue, fill=popblue!10] (-2.2,0) ellipse (1.0 and 1.8); % VD
\draw[thick, poporange, fill=poporange!10] (2.2,0) ellipse (1.0 and 1.8);  % VG
% Oreillettes
\draw[thick, popblue, fill=popblue!10] (-2.2,2.4) ellipse (0.9 and 0.7); % OD
\draw[thick, poporange, fill=poporange!10] (2.2,2.4) ellipse (0.9 and 0.7); % OG
% Septum interventriculaire
\draw[thick, popdark] (0,-1.8) -- (0,1.8);
% Valvules (schématique)
\draw[thick, popdark] (-2.8,1.2) -- (-1.6,1.2); % Tricuspide
\draw[thick, popdark] (1.6,1.2) -- (2.8,1.2);   % Mitrale
\draw[thick, popdark] (-1.2,0.2) -- (-0.2,0.2); % Sigmoïde pulmonaire
\draw[thick, popdark] (0.2,0.2) -- (1.2,0.2);   % Sigmoïde aortique
% Gros vaisseaux
% Veines caves (vers OD)
\draw[thick, popblue] (-2.2,3.1) -- (-2.2,4.0);
\draw[thick, popblue] (-3.0,3.1) -- (-3.0,4.0);
% Artère pulmonaire (depuis VD)
\draw[thick, popblue] (-1.2,1.8) .. controls (-0.5,2.5) and (0.5,2.5) .. (1.2,1.8); % Trunk commun
\draw[thick, popblue] (-1.2,1.8) -- (-1.2,3.0); % Branche gauche
\draw[thick, popblue] (1.2,1.8) -- (1.2,3.0);   % Branche droite
% Veines pulmonaires (vers OG)
\draw[thick, poporange] (3.0,3.1) -- (3.0,4.0);
\draw[thick, poporange] (3.8,3.1) -- (3.8,4.0);
% Aorte (depuis VG)
\draw[thick, poporange] (1.2,1.8) -- (1.2,3.0) -- (1.2,4.2); % Montante
\draw[thick, poporange] (1.2,3.6) .. controls (2.5,3.6) and (2.5,2.5) .. (3.2,2.5); % Crosse aortique (schématique)
% Lettres
\node[popblue] at (-2.2,-0.2) {\textbf{A}};
\node[poporange] at (2.2,-0.2) {\textbf{B}};
\node[popblue] at (-2.2,2.4) {\textbf{C}};
\node[poporange] at (2.2,2.4) {\textbf{D}};
\node[popblue] at (-1.2,3.3) {\textbf{E}};
\node[poporange] at (1.2,4.5) {\textbf{F}};
% Légende interne
\node[align=left, anchor=west, font=\small] at (3.5,1.5) {
\textbf{A, B} : Ventricules\\
\textbf{C, D} : Oreillettes\\
\textbf{E, F} : Gros vaisseaux artériels
};
\end{tikzpicture}
\legende{Coupe frontale schématique du cœur humain. En bleu : territoire veineux / sang désoxygéné (cœur droit). En orange : territoire artériel / sang oxygéné (cœur gauche).}
\end{popfigure}
\begin{enumerate}
\item Nomme les cavités A, B, C et D en précisant « droite » ou « gauche ».
\item Le vaisseau E part de la cavité A et se dirige vers les poumons : nomme-le. Le vaisseau F part de la cavité B et irrigue tout le corps : nomme-le.
\item Indique, en utilisant les lettres du schéma, le trajet suivi par une molécule de dioxygène depuis les poumons jusqu'à un muscle de la jambe (on admettra qu'elle arrive au c\oe ur par les veines pulmonaires vers la cavité D).
\item Explique pourquoi on parle de \og double circulation \fg{} chez l'Homme.
\end{enumerate}
\textbf{Partie B : Le c\oe ur, organe moteur.}
\begin{enumerate}[resume]
\item Rappelle ce qu'on appelle \cle{systole} et \cle{diastole} (précise pour les oreillettes et les ventricules).
\item Le c\oe ur d'un adulte au repos bat 75 fois par minute. Combien de battements effectue-t-il en une journée de 24 heures ? En une année de 365 jours ? Donne les résultats en notation scientifique.
\end{enumerate}
\end{exercice}

\begin{exercice}[Un cas clinique : l'AVC de M. Ngoa]
M. Ngoa, 52 ans, commerçant à Douala, est hypertendu depuis plusieurs années, fume un paquet de cigarettes par jour et ne pratique aucune activité physique. Un matin, il perd brutalement la parole et ne peut plus lever le bras droit. Sa famille le conduit d'urgence à l'hôpital, où le médecin diagnostique un accident vasculaire cérébral (AVC) ischémique (infarctus cérébral).
\begin{enumerate}
\item Définis un accident vasculaire cérébral (AVC).
\item Explique, à l'aide de tes connaissances, le mécanisme de cet accident ischémique au niveau des vaisseaux du cerveau (athérosclérose, thrombus, ischémie, nécrose).
\item Relève dans le texte trois facteurs de risque qui ont favorisé l'AVC de M. Ngoa, et justifie pour chacun son rôle pathophysiologique.
\item Quels signes observés par la famille (troubles de la parole, paralysie) auraient dû alerter immédiatement ? Cite le mnémotechnique FAST (ou VITE en français).
\item Propose quatre mesures de prévention des maladies cardiovasculaires que tu donnerais à un adulte de ton entourage.
\end{enumerate}
\end{exercice}

\begin{exercice}[Synthèse : circulation sanguine et hygiène de vie]
Lors d'une campagne de santé scolaire dans un collège de Bafoussam, on relève la tension artérielle et la fréquence cardiaque de repos de deux élèves de 3\textsuperscript{e} :
\begin{center}
\begin{tabularx}{\linewidth}{|l|X|X|}\hline
\rowcolor{popblueL}
 & \textbf{Élève 1 (sportif régulier)} & \textbf{Élève 2 (sédentaire)} \\ \hline
Fréquence cardiaque de repos (battements/min) & 60 & 82 \\ \hline
Temps de récupération après un effort standardisé (min) & 3 & 9 \\ \hline
\end{tabularx}
\end{center}
\begin{enumerate}
\item Compare les résultats des deux élèves et interprète les différences observées en lien avec l'adaptation cardiaque à l'entraînement.
\item Explique pourquoi la pratique régulière d'une activité physique est bénéfique pour le c\oe ur (hypertrophie physiologique, volume d'éjection, fréquence de repos) et pour les vaisseaux (vascularisation, élasticité, profil lipidique).
\item Cite trois autres règles d'hygiène de vie (alimentation, tabac, stress, sommeil) qui protègent l'appareil circulatoire.
\item Un élève se blesse profondément au bras pendant la coupe d'un bois en travaux pratiques : le sang coule en continu, de façon abondante, de couleur rouge foncé (veineuse) mais le débit est important. Décris la conduite à tenir, étape par étape, jusqu'à l'arrivée des secours.
\item Rédige un court texte (8 à 10 lignes) de sensibilisation destiné aux élèves du collège, intitulé \og Protégeons notre c\oe ur et nos vaisseaux \fg.
\end{enumerate}
\end{exercice}



\newpage

\coursec{Corrigés des exercices}

\begin{corrige}[Exercice 1 — QCM : les vaisseaux et la circulation]
\begin{enumerate}
\item Bonnes réponses : \textbf{b} et \textbf{c}. Les artères transportent le sang du \cle{cœur} vers les organes et possèdent une paroi épaisse et élastique qui résiste à la pression artérielle. La réponse a décrit les veines (paroi mince, valvules, sang vers le cœur) ; la réponse d est fausse car l'artère pulmonaire contient du sang pauvre en \cle{dioxygène}.
\item Bonnes réponses : \textbf{a}, \textbf{b} et \textbf{d}. Les capillaires ont une paroi très fine (une seule couche de cellules endothéliales), ce qui permet les échanges entre le sang et les organes ; ils relient les artérioles aux veinules. Leur diamètre est au contraire très petit (environ \SI{5}{\micro\meter} à \SI{10}{\micro\meter}), ce qui rend la réponse c fausse.
\item Bonnes réponses : \textbf{a} et \textbf{c}. La \cle{petite circulation} va du ventricule droit vers les poumons (par l'artère pulmonaire) et revient à l'oreillette gauche (par les veines pulmonaires). Les réponses b et d décrivent la \cle{grande circulation}.
\item Bonnes réponses : \textbf{b} et \textbf{c}. Face à un \cle{AVC}, il faut alerter immédiatement les secours (112 ou 119 au Cameroun) et installer la victime sur le dos, la tête légèrement surélevée (30°). Il ne faut \textbf{jamais} donner à boire (risque de fausse route) ni faire marcher la victime.
\end{enumerate}
\textbf{Point de vigilance :} ne pas confondre « artère = sang oxygéné » : c'est le \emph{sens de la circulation} (qui part du cœur) qui définit une artère, pas la richesse en dioxygène. L'artère pulmonaire transporte du sang désoxygéné.
\end{corrige}

\begin{corrige}[Exercice 2 — Vrai ou faux, en justifiant]
\begin{enumerate}
\item \textbf{Vrai.} Le cœur est un muscle creux (le \cle{myocarde}) qui se contracte (\cle{systole}) et se relâche (\cle{diastole}) de façon automatique et rythmée, sans ordre volontaire, environ 60 à 100 fois par minute au repos (fréquence cardiaque de repos).
\item \textbf{Vrai.} Les veines des membres possèdent des \cle{valvules}, sortes de petits clapets, qui empêchent le reflux du sang et l'obligent à circuler vers le cœur (retour veineux).
\item \textbf{Faux.} C'est le contraire : le ventricule gauche a une paroi beaucoup plus épaisse (environ \SI{10}{\milli\meter} à \SI{15}{\milli\meter} chez l'adulte, \SI{10}{\milli\meter} contre \SI{4}{\milli\meter} chez le mouton disséqué en TP), car il doit propulser le sang dans tout le corps (grande circulation), ce qui demande une pression plus forte que la circulation vers les poumons (petite circulation).
\item \textbf{Vrai.} Une hémorragie externe grave fait perdre rapidement une grande quantité de sang ; la pression sanguine s'effondre (\cle{choc hypovolémique}), les organes ne sont plus irrigués, et la victime peut mourir en quelques minutes si l'on ne comprime pas la plaie et si les secours n'arrivent pas vite.
\end{enumerate}
\end{corrige}

\begin{corrige}[Exercice 3 — Définitions]
\begin{itemize}
\item \cle{Artère} : vaisseau sanguin à paroi épaisse, musculaire et élastique qui transporte le sang du cœur vers les organes.
\item \cle{Veine} : vaisseau sanguin à paroi mince, souvent muni de valvules, qui ramène le sang des organes vers le cœur.
\item \cle{Capillaire} : vaisseau sanguin microscopique, à paroi très fine (une seule couche cellulaire), reliant les artérioles aux veinules et siège des échanges (dioxygène, nutriments, déchets, dioxyde de carbone) entre le sang et les organes.
\item \cle{Hémorragie} : écoulement de sang hors des vaisseaux, à la suite d'une rupture d'un vaisseau sanguin ; elle est \emph{externe} si le sang sort du corps, \emph{interne} sinon.
\item \cle{Accident vasculaire cérébral (AVC)} : arrêt brutal de la circulation du sang dans une partie du cerveau, par obstruction (AVC ischémique, \SI{80}{\%} des cas) ou rupture (AVC hémorragique) d'un vaisseau cérébral, entraînant la mort des cellules nerveuses de la zone non irriguée (\cle{nécrose} ou infarctus cérébral).
\end{itemize}
\end{corrige}

\begin{corrige}[Exercice 4 — Fréquence cardiaque et débit cardiaque]
\begin{enumerate}
\item Débit cardiaque au repos (volume d'éjection systolique \cle{VES} = \SI{70}{\milli\liter}) :
\[ Q = f \times \text{VES} = \SI{72}{\per\minute} \times \SI{70}{\milli\liter} = \SI{5040}{\milli\liter\per\minute} = \mathbf{\SI{5,04}{\liter\per\minute}}. \]
\item Débit cardiaque pendant l'effort (VES supposé constant à \SI{70}{\milli\liter}) :
\[ Q' = \SI{150}{\per\minute} \times \SI{70}{\milli\liter} = \SI{10500}{\milli\liter\per\minute} = \mathbf{\SI{10,5}{\liter\per\minute}}. \]
\item Comparaison : $\SI{10,5}{\liter\per\minute} \div \SI{5,04}{\liter\per\minute} \approx 2,1$. Le débit cardiaque a plus que doublé pendant l'effort. \textbf{Conclusion :} pendant l'effort, le cœur accélère sa fréquence (chronotropie positive) pour augmenter son débit, afin d'apporter plus de dioxygène et de nutriments aux muscles qui travaillent et d'éliminer le dioxyde de carbone produit.
\end{enumerate}
\textbf{Point de vigilance :} bien convertir les millilitres en litres (\SI{1000}{\milli\liter} = \SI{1}{\liter}) avant de donner le résultat final ; ne pas oublier l'unité (\SI{}{\liter\per\minute} ou \SI{}{\milli\liter\per\minute}).
\end{corrige}

\begin{corrige}[Exercice 5 — L'expérience de Harvey]
\begin{enumerate}
\item Le garrot comprime les veines superficielles du bras et bloque le retour du sang vers le cœur. Le sang continue d'arriver dans l'avant-bras par les artères (plus profondes, moins comprimées) mais ne peut plus repartir : il s'accumule donc dans les veines \emph{en aval} du garrot (côté main), qui gonflent et deviennent visibles.
\item Quand on pousse le sang vers le coude, la veine se vide jusqu'au premier renflement, puis se remplit à nouveau \emph{par le côté main}. Cela montre que le sang circule dans les veines \textbf{de la périphérie (la main) vers le cœur}, et qu'il ne peut pas revenir en arrière.
\item Les renflements correspondent aux \cle{valvules} : ce sont des clapets qui laissent passer le sang vers le cœur mais se ferment si le sang tend à refluer vers la main. Elles assurent donc la circulation du sang veineux dans un seul sens (un sens : périphérie $\to$ cœur).
\end{enumerate}
\textbf{Point de vigilance :} « en aval du garrot » signifie ici du côté de la main (l'aval du flux veineux) ; bien orienter le raisonnement avec le sens réel de la circulation veineuse (vers le cœur).
\end{corrige}

\begin{corrige}[Exercice 6 — Structure et fonction des cavités cardiaques]
\begin{enumerate}
\item Ordre croissant d'épaisseur (valeurs moyennes observées sur cœur de mouton en TP) :
\cle{Oreillette droite} (\SI{\approx 2}{\milli\meter}) $<$ \cle{Oreillette gauche} (\SI{\approx 3}{\milli\meter}) $<$ \cle{Ventricule droit} (\SI{\approx 4}{\milli\meter}) $<$ \cle{Ventricule gauche} (\SI{\approx 10}{\milli\meter}).
\item Rôles :
\begin{itemize}
\item \textbf{Oreillette droite} : reçoit le sang pauvre en dioxygène venant de tout le corps par les veines caves (supérieure et inférieure) et l'envoie dans le ventricule droit.
\item \textbf{Ventricule droit} : propulse ce sang vers les poumons par l'artère pulmonaire (petite circulation).
\item \textbf{Oreillette gauche} : reçoit le sang riche en dioxygène venant des poumons par les veines pulmonaires et l'envoie dans le ventricule gauche.
\item \textbf{Ventricule gauche} : propulse le sang oxygéné dans l'aorte vers tout le corps (grande circulation).
\end{itemize}
\item Relation structure-fonction : plus une cavité doit fournir un travail mécanique important (propulser le sang loin et sous forte pression), plus sa paroi musculaire (myocarde) est épaisse. Les oreillettes n'ont qu'à pousser le sang dans les ventricules voisins (pression faible, parois minces) ; le ventricule droit envoie le sang aux poumons, proches et à basse résistance (paroi moyenne) ; le ventricule gauche doit irriguer tout l'organisme contre une forte résistance vasculaire, ce qui exige la pression la plus forte : sa paroi est la plus épaisse. \textbf{La structure est adaptée à la fonction.}
\end{enumerate}
\end{corrige}

\begin{corrige}[Exercice 7 — Premiers secours : hémorragie externe grave]
\begin{enumerate}
\item Il s'agit d'une \textbf{hémorragie artérielle}. Deux arguments : le sang est \emph{rouge vif} (sang riche en dioxygène, artériel) et il s'écoule \emph{en jets synchrones des battements du cœur}, car il est sous forte pression dans les artères (pression systolique).
\item Gestes de premiers secours, dans l'ordre strict :
\begin{enumerate}
\item \textbf{Alerter} immédiatement les secours (SAMU : \textbf{112} ou \textbf{119} au Cameroun) ou faire alerter par un témoin, en précisant le lieu, le nombre de victimes et la gravité (jet artériel).
\item \textbf{Comprimer fortement la plaie} avec les mains protégées (gants à usage unique ou tissu propre), directement sur le point qui saigne (compression manuelle digitale ou palmaire).
\item Si possible, maintenir la compression avec un \textbf{pansement compressif} serré (bande + compresse) sans relâcher la pression manuelle avant sa mise en place.
\item \textbf{Allonger la victime} en position horizontale (ou demi-assise si détresse respiratoire), la couvrir pour éviter le refroidissement (risque de choc), et \textbf{ne jamais} lui donner à boire ni à manger (risque de vomissements et d'anesthésie ultérieure).
\item Surveiller la victime (respiration, conscience, pouls radial) jusqu'à l'arrivée des secours, sans jamais relâcher la compression.
\end{enumerate}
\item Il faut agir très vite car le sang artériel jaillit sous pression (pression systolique \SI{\approx 120}{\milli\meter Hg}) : la perte de sang est très rapide (jusqu'à \SI{1}{\liter} en quelques minutes). En quelques minutes, la volémie résiduelle devient insuffisante pour assurer la perfusion des organes vitaux (cerveau, cœur, reins) : la victime perd conscience et risque la mort. La compression immédiate est le seul geste qui stoppe cette perte.
\end{enumerate}
\textbf{Point de vigilance :} ne pas retirer un objet planté dans la plaie (comprimer autour) et ne pas poser de garrot (geste dangereux réservé à des cas exceptionnels : membre sectionné, compression impossible) : la compression directe est le geste de référence enseigné au BEPC.
\end{corrige}

\begin{corrige}[Exercice 8 — Fréquence cardiaque et effort : analyse graphique]
\textbf{1. Tracé de la courbe :}
\begin{popfigure}
\begin{tikzpicture}
\begin{axis}[
    width=\linewidth, height=8cm,
    xlabel={Temps (min)}, ylabel={Fréquence cardiaque (battements/min)},
    xmin=-0.5, xmax=16, ymin=50, ymax=170,
    xtick={0,2,5,7,10,15}, ytick={50,70,90,110,130,150,170},
    grid=major, grid style={popline},
    axis lines=middle,
    legend pos=north east,
    tick label style={font=\small},
    label style={font=\small}
]
\addplot[popcurve, popblue, mark=*, mark size=2pt] coordinates {
    (0,70) (2,120) (5,150) (7,110) (10,85) (15,72)
};
\addlegendentry{Fréquence cardiaque mesurée}
\draw[dashed, poporange, thick] (axis cs:5,50) -- (axis cs:5,170);
\node[poporange, anchor=south, font=\small\bfseries] at (axis cs:2.5,155) {Phase d'effort};
\node[popgreen, anchor=south, font=\small\bfseries] at (axis cs:10,155) {Phase de récupération};
\end{axis}
\end{tikzpicture}
\legende{Variations de la fréquence cardiaque avant, pendant (0 à 5 min) et après (5 à 15 min) un effort physique de 5 minutes.}
\end{popfigure}

\textbf{2. Description :} au repos ($t = \SI{0}{\minute}$), la fréquence cardiaque est de \SI{70}{\per\minute}. Pendant la phase d'effort (\SI{0}{\minute} à \SI{5}{\minute}), elle augmente rapidement : \SI{120}{\per\minute} à \SI{2}{\minute}, puis \SI{150}{\per\minute} à \SI{5}{\minute} (maximum). Pendant la phase de récupération (\SI{5}{\minute} à \SI{15}{\minute}), elle diminue d'abord vite (\SI{110}{\per\minute} à \SI{7}{\minute}, \SI{85}{\per\minute} à \SI{10}{\minute}), puis plus lentement, jusqu'à retrouver environ sa valeur de repos (\SI{72}{\per\minute} à \SI{15}{\minute}).

\textbf{3. Explication physiologique :} pendant l'effort, les muscles squelettiques consomment beaucoup plus de dioxygène et de nutriments (glucose, acides gras) et rejettent plus de dioxyde de carbone. Pour répondre à ces besoins accrus, le système nerveux autonome (orthosympathique) accélère la fréquence cardiaque (et augmente le volume d'éjection) : le \cle{débit cardiaque} augmente, ce qui accélère l'apport de dioxygène aux muscles et l'élimination des déchets métaboliques. C'est une adaptation automatique de l'organisme à l'exercice.

\textbf{4. Temps de récupération :} fin de l'effort à $t = \SI{5}{\minute}$ ; retour à la valeur de repos (\SI{\approx 70}{\per\minute}) à $t = \SI{15}{\minute}$ :
\[ t_{\text{récupération}} = \SI{15}{\minute} - \SI{5}{\minute} = \mathbf{\SI{10}{\minute}}. \]
\textbf{Point de vigilance :} le temps de récupération se compte à partir de la \emph{fin} de l'effort, pas à partir du début. Une récupération rapide (moins de \SI{5}{\minute}) signe un bon entraînement cardiovasculaire.
\end{corrige}

\begin{corrige}[Exercice 9 — Le cœur et la double circulation]
\textbf{Partie A : Anatomie du cœur (coupe frontale).}
\begin{enumerate}
\item \textbf{A :} ventricule droit ; \textbf{B :} ventricule gauche ; \textbf{C :} oreillette droite ; \textbf{D :} oreillette gauche. (Rappel : sur une coupe frontale vue de face, le côté droit du cœur apparaît à \emph{gauche} du schéma.)
\item \textbf{E :} artère pulmonaire (elle part du ventricule droit vers les poumons). \textbf{F :} aorte (elle part du ventricule gauche vers tout le corps).
\item Trajet d'une molécule de dioxygène : poumons $\to$ veines pulmonaires $\to$ \textbf{D} (oreillette gauche) $\to$ \textbf{B} (ventricule gauche) $\to$ \textbf{F} (aorte) $\to$ artères puis artérioles de la jambe $\to$ capillaires du muscle $\to$ cellules musculaires (diffusion dans la mitochondrie pour la respiration cellulaire).
\item On parle de \og \cle{double circulation} \fg{} parce que le sang passe \textbf{deux fois par le cœur} au cours d'un tour complet : la \textbf{petite circulation} (cœur droit $\to$ poumons $\to$ cœur gauche), où le sang se charge en dioxygène, et la \textbf{grande circulation} (cœur gauche $\to$ organes $\to$ cœur droit), où le sang distribue le dioxygène et récupère le dioxyde de carbone. Le cœur droit et le cœur gauche fonctionnent comme deux pompes en série.
\end{enumerate}
\textbf{Partie B : Le cœur, organe moteur.}
\begin{enumerate}[resume]
\item La \cle{systole} est la phase de \textbf{contraction} du muscle cardiaque : la systole des oreillettes (brève) chasse le sang vers les ventricules, puis la systole des ventricules (puissante) éjecte le sang dans les artères (aorte et artère pulmonaire). La \cle{diastole} est la phase de \textbf{relâchement} : les cavités se remplissent de sang (les oreillettes par les veines, puis les ventricules par les oreillettes ouvertes).
\item Nombre de battements en une journée (fréquence \SI{75}{\per\minute}) :
\[ 75 \times 60 \times 24 = 108\,000 = \mathbf{\num{1,08e5}} \ \text{battements}. \]
En une année de 365 jours :
\[ \num{1,08e5} \times 365 = 39\,420\,000 = \mathbf{\num{3,942e7}} \ \text{battements}. \]
Le cœur accomplit donc un travail considérable, sans jamais s'arrêter : cela justifie l'importance de l'hygiène de vie pour le protéger.
\end{enumerate}
\textbf{Barème indicatif (type BEPC) :} Partie A : 1 pt par cavité correctement nommée, 1 pt par vaisseau, 2 pts pour le trajet complet, 2 pts pour la double circulation. Partie B : 2 pts pour systole/diastole, 2 pts pour les calculs avec notation scientifique.
\textbf{Points de vigilance :} ne pas inverser droite et gauche sur le schéma ; vérifier que le trajet demandé commence bien aux poumons (donc par le cœur gauche) ; en notation scientifique, un seul chiffre non nul avant la virgule.
\end{corrige}

\begin{corrige}[Exercice 10 — Un cas clinique : l'AVC de M. Ngoa]
\begin{enumerate}
\item Un \cle{accident vasculaire cérébral (AVC)} est l'arrêt brutal de la circulation du sang dans une région du cerveau, par obstruction d'un vaisseau (AVC ischémique, thrombose ou embolie) ou par rupture d'un vaisseau (AVC hémorragique). Privées de sang, donc de dioxygène et de glucose, les cellules nerveuses (neurones) de cette région meurent en quelques minutes (\cle{nécrose} / infarctus cérébral).
\item Mécanisme de l'AVC ischémique chez M. Ngoa : l'hypertension artérielle et le tabac ont favorisé le dépôt de plaques de graisse (cholestérol, cellules) sur la paroi interne des artères cérébrales : c'est l'\cle{athérosclérose}. La paroi s'épaissit, perd son élasticité, le calibre de l'artère diminue (sténose). Un caillot de sang (\cle{thrombus}) peut alors se former sur la plaque d'athérome et boucher complètement l'artère : la zone du cerveau en aval n'est plus irriguée, c'est l'\cle{ischémie}. Sans dioxygène, les cellules nerveuses de cette zone meurent par nécrose (infarctus cérébral), responsable des troubles observés (paralysie, aphasie).
\item Trois facteurs de risque relevés et justifiés :
\begin{itemize}
\item \textbf{L'hypertension artérielle (HTA)} : la pression élevée use et fragilise les parois des artères (micro-lésions) et accélère le dépôt des plaques d'athérome.
\item \textbf{Le tabagisme} (un paquet par jour) : le monoxyde de carbone et la nicotine favorisent le dépôt des plaques d'athérome, augmentent la coagulabilité du sang (risque de thrombus) et provoquent des spasmes artériels.
\item \textbf{La sédentarité} (aucune activité physique) : elle favorise la prise de poids, l'hypertension, le diabète de type 2 et le mauvais état des vaisseaux (perte d'élasticité).
\end{itemize}
\item Les signes qui auraient dû alerter immédiatement (survenue brutale) : la \textbf{perte de la parole} (aphasie) et la \textbf{paralysie du bras droit} (hémiplégie controlatérale). Le mnémotechnique \textbf{VITE} (recommandé au Cameroun) :
\begin{itemize}
\item \textbf{V} : Visu (trouble brutal de la vision : perte d'un côté, vision double) ;
\item \textbf{I} : Incapacité (faiblesse ou paralysie brutale d'un bras, d'une jambe, d'un côté du visage) ;
\item \textbf{T} : Trouble de la parole (impossible de parler, mots incompréhensibles) ;
\item \textbf{E} : Extrême urgence : appeler immédiatement les secours (\textbf{112} ou \textbf{119}).
\end{itemize}
\item Quatre mesures de prévention primaire :
\begin{enumerate}
\item pratiquer une activité physique régulière (marche rapide, sport, au moins \SI{30}{\minute} par jour, 5 fois par semaine) ;
\item ne pas fumer et éviter la consommation excessive d'alcool ;
\item adopter une alimentation équilibrée, pauvre en sel (limite l'HTA), en graisses saturées et en sucres, riche en fruits, légumes et poissons ;
\item contrôler régulièrement sa tension artérielle (au moins une fois par an) et suivre scrupuleusement son traitement antihypertenseur en cas d'HTA.
\end{enumerate}
\end{enumerate}
\textbf{Barème indicatif (type BEPC) :} définition 2 pts ; mécanisme (athérosclérose $\to$ thrombus $\to$ ischémie $\to$ nécrose) 4 pts ; 3 facteurs justifiés 3 pts ; signes + VITE 2 pts ; prévention 2 pts.
\textbf{Points de vigilance :} bien enchaîner les étapes du mécanisme dans l'ordre logique (cause $\to$ lésion vasculaire $\to$ obstruction $\to$ conséquence cellulaire) ; chaque facteur de risque doit être \emph{justifié physiologiquement}, pas seulement cité ; insister sur l'urgence : chaque minute compte pour limiter l'étendue de la nécrose (thrombolyse possible < 4h30).
\end{corrige}

\begin{corrige}[Exercice 11 — Synthèse : circulation sanguine et hygiène de vie]
\begin{enumerate}
\item Comparaison : l'élève 1 (sportif régulier) a une fréquence cardiaque de repos plus basse (\SI{60}{\per\minute} contre \SI{82}{\per\minute}, \cle{bradycardie} physiologique) et récupère beaucoup plus vite (\SI{3}{\minute} contre \SI{9}{\minute}). Interprétation : l'entraînement régulier rend le cœur plus puissant et plus efficace : à chaque battement, il éjecte un plus grand volume de sang (volume d'éjection systolique augmenté), donc il a besoin de battre moins souvent au repos pour assurer le même débit cardiaque (économie cardiaque). Après l'effort, son cœur revient rapidement à son rythme de repos, signe d'une bonne adaptation cardiaque et d'un système nerveux parasympathique réactif. L'élève sédentaire, au contraire, a un cœur moins entraîné, qui doit battre plus vite et plus longtemps pour fournir l'effort et récupérer.
\item Bénéfices de l'activité physique régulière sur l'appareil circulatoire :
\begin{itemize}
\item \textbf{Pour le cœur :} le muscle cardiaque se renforce (hypertrophie physiologique, augmentation de la masse ventriculaire gauche), le volume d'éjection systolique augmente, et la fréquence cardiaque de repos diminue : le cœur travaille plus efficacement (meilleur rendement) et se fatigue moins pour un même travail.
\item \textbf{Pour les vaisseaux :} la vascularisation (capillarisation) des organes s'améliore, les parois artérielles gardent leur élasticité (lutte contre la rigidité artérielle), et l'activité physique limite le dépôt de graisses (plaques d'athérome), ce qui réduit le risque d'hypertension artérielle, d'infarctus du myocarde et d'AVC.
\end{itemize}
\item Trois autres règles d'hygiène de vie essentielles :
\begin{itemize}
\item avoir une alimentation équilibrée, pauvre en sel ($< \SI{5}{\gram}$/jour), en graisses saturées et en sucres ajoutés, riche en fruits, légumes, légumineuses et poissons (régime méditerranéen type) ;
\item ne pas fumer (le tabac détruit l'endothélium vasculaire, favorise l'athérosclérose, la thrombose et le spasme artériel) ;
\item limiter le stress chronique et dormir suffisamment (\SIrange{7}{9}{\hour} par nuit), car le stress et le manque de sommeil augmentent la tension artérielle et la fréquence cardiaque au repos.
\end{itemize}
\item Conduite à tenir face à l'hémorragie veineuse abondante (sang rouge foncé, écoulement continu, débit important) :
\begin{enumerate}
\item se protéger (gants à usage unique ou interposition d'un tissu propre / sac plastique) et \textbf{alerter ou faire alerter} les secours (\textbf{112} ou \textbf{119}) et un adulte responsable ;
\item \textbf{comprimer fortement la plaie} directement avec un linge propre (compresse si possible), sans relâcher la pression (compression manuelle maintenue) ;
\item allonger l'élève blessé, surélever légèrement le bras blessé (si pas de fracture suspectée) pour favoriser le retour veineux, et maintenir la compression (pansement compressif si disponible, sans garrot) ;
\item ne rien donner à boire ni à manger, couvrir le blessé pour le garder au chaud (prévention du choc hypothermique) ;
\item le rassurer et surveiller sa conscience et sa respiration jusqu'à l'arrivée des secours.
\end{enumerate}
Même si le sang est rouge foncé (hémorragie veineuse) et coule en continu (pas de jets), le débit important impose la même conduite : compression immédiate et appel des secours.
\item Exemple de texte de sensibilisation (affichage classe / journal scolaire) :\\[2mm]
\og \textbf{Protégeons notre cœur et nos vaisseaux.} Notre cœur bat plus de \num{100000} fois par jour pour faire circuler le sang dans tout notre corps : c'est un organe précieux et infatigable. Pour le garder en bonne santé, adoptons de bonnes habitudes dès maintenant : pratiquons un sport régulièrement, mangeons équilibré, avec moins de sel, de graisses et de sucres, et plus de fruits et de légumes. Refusons le tabac et l'alcool, qui bouchent et fragilisent nos vaisseaux. Dormons bien et apprenons à gérer notre stress. Apprenons aussi les gestes qui sauvent : comprimer une plaie qui saigne et appeler vite les secours (\textbf{112} ou \textbf{119}). Un cœur bien entretenu aujourd'hui, c'est une vie longue et en bonne santé demain ! \fg
\end{enumerate}
\textbf{Point de vigilance :} dans la comparaison des deux élèves, toujours relier chaque différence chiffrée (FC repos, temps de récupération) à une explication physiologique sur le fonctionnement du cœur (volume d'éjection, système nerveux autonome) ; une simple lecture du tableau ne suffit pas au BEPC.
\end{corrige}

\newpage

\coursec{Fiche bilan}

\begin{popfigure}
\centering
\begin{center}\tcbox[colback=poporange,colframe=poporange,coltext=white,arc=9pt,boxsep=4pt,left=6pt,right=6pt]{\bfseries\small Circulation sanguine}\end{center}
\vspace{-2pt}
\begin{tcbraster}[raster columns=3,raster equal height=rows,raster column skip=6pt,raster row skip=6pt,enhanced,blankest]
\begin{tcolorbox}[enhanced,colback=white,colframe=poporange,boxrule=0.9pt,arc=3pt,title={Vaisseaux sanguins},fonttitle=\bfseries\scriptsize,coltitle=white,colbacktitle=poporange,left=4pt,right=4pt,top=2pt,bottom=2pt,boxsep=1pt,toptitle=1.5pt,bottomtitle=1.5pt,halign=left]
\begin{itemize}[nosep,leftmargin=*,itemsep=1pt,font=\scriptsize]
\item Artères (haute pression)
\item Veines (valves, réserve)
\item Capillaires (échanges)
\end{itemize}
\end{tcolorbox}
\begin{tcolorbox}[enhanced,colback=white,colframe=popgreen,boxrule=0.9pt,arc=3pt,title={Cœur (moteur)},fonttitle=\bfseries\scriptsize,coltitle=white,colbacktitle=popgreen,left=4pt,right=4pt,top=2pt,bottom=2pt,boxsep=1pt,toptitle=1.5pt,bottomtitle=1.5pt,halign=left]
\begin{itemize}[nosep,leftmargin=*,itemsep=1pt,font=\scriptsize]
\item 4 cavités (2 oreillettes, 2 ventricules)
\item Valves (anti-retour)
\item Double circulation (poumon / systémique)
\item Cycle cardiaque (systole / diastole)
\end{itemize}
\end{tcolorbox}
\begin{tcolorbox}[enhanced,colback=white,colframe=poppink,boxrule=0.9pt,arc=3pt,title={Hémorragies \& secours},fonttitle=\bfseries\scriptsize,coltitle=white,colbacktitle=poppink,left=4pt,right=4pt,top=2pt,bottom=2pt,boxsep=1pt,toptitle=1.5pt,bottomtitle=1.5pt,halign=left]
\begin{itemize}[nosep,leftmargin=*,itemsep=1pt,font=\scriptsize]
\item Types (artérielle, veineuse, capillaire)
\item Geste vital \textbf{Compression}
\item Garrot (dernier recours)
\end{itemize}
\end{tcolorbox}
\begin{tcolorbox}[enhanced,colback=white,colframe=poppurple,boxrule=0.9pt,arc=3pt,title={Maladies cardiovasculaires},fonttitle=\bfseries\scriptsize,coltitle=white,colbacktitle=poppurple,left=4pt,right=4pt,top=2pt,bottom=2pt,boxsep=1pt,toptitle=1.5pt,bottomtitle=1.5pt,halign=left]
\begin{itemize}[nosep,leftmargin=*,itemsep=1pt,font=\scriptsize]
\item Facteurs risque (tabac, sel, sédentarité)
\item Athérosclérose (plaques athérome)
\item HTA (tension > 14/9)
\item Prévention (hygiène de vie)
\end{itemize}
\end{tcolorbox}
\begin{tcolorbox}[enhanced,colback=white,colframe=popgold,boxrule=0.9pt,arc=3pt,title={AVC (urgence vitale)},fonttitle=\bfseries\scriptsize,coltitle=white,colbacktitle=popgold,left=4pt,right=4pt,top=2pt,bottom=2pt,boxsep=1pt,toptitle=1.5pt,bottomtitle=1.5pt,halign=left]
\begin{itemize}[nosep,leftmargin=*,itemsep=1pt,font=\scriptsize]
\item Signes : \textbf{VITE} (Visage, Incapacité, Trouble parole, Extrême urgence)
\item Types (ischémique / hémorragique)
\item Premiers secours (PLS, 15/112, noter heure)
\end{itemize}
\end{tcolorbox}
\end{tcbraster}

\legende{Carte mentale récapitulative de la leçon 9 : Circulation sanguine (5 axes principaux).}
\end{popfigure}

\begin{bilan}

\courssub{Les vaisseaux sanguins, siège de la circulation}

\begin{definition}[Vaisseaux sanguins]
Tuyaux fermés assurant le transport du sang entre le cœur et les organes. Trois types histologiques distincts.
\end{definition}

\begin{center}
\begin{tabularx}{\linewidth}{|l|X|X|X|}\hline
\rowcolor{popblueL}
\textbf{Caractère} & \textbf{Artère} & \textbf{Veine} & \textbf{Capillaire} \\ \hline
\textbf{Paroi} & Épaisse, musculaire, élastique & Mince, peu musculaire, valves & Très fine (1 couche endothélium) \\ \hline
\textbf{Lumière} & Étroite & Large & Très fine ($\approx 8\,\mu m$) \\ \hline
\textbf{Pression} & Élevée (pulsatile) & Basse & Basse \\ \hline
\textbf{Sens} & Cœur $\to$ Organes & Organes $\to$ Cœur & Réseau d'échange \\ \hline
\textbf{Rôle principal} & Conduite rapide & Réserve, retour veineux & Échanges O$_2$/CO$_2$, nutriments/déchets \\ \hline
\end{tabularx}
\end{center}

\courssub{Le cœur, organe moteur de la circulation}

\begin{definition}[Cœur]
Organe musculaire creux (myocarde), divisé en 4 cavités, assurant la propulsion du sang par contractions rythmiques involontaires (automatisme).
\end{definition}

\begin{itemize}
\item \textbf{Cavités :} 2 oreillettes (supérieures, réception) + 2 ventricules (inférieures, expulsion). Séparation stricte gauche/droite (cloison interventriculaire).
\item \textbf{Gros vaisseaux :} Veines caves (sup/inf) $\to$ Oreillette droite ; Artère pulmonaire $\leftarrow$ Ventricule droit ; Veines pulmonaires $\to$ Oreillette gauche ; Aorte $\leftarrow$ Ventricule gauche.
\item \textbf{Valves :} Tricuspide (OD/VD), Sigmoïde pulmonaire (VD/AP), Bicuspide/Mitrale (OG/VG), Sigmoïde aortique (VG/Ao). Rôle : \cle{anti-retour}.
\item \textbf{Double circulation :}
      \begin{enumerate}
      \item \textbf{Petite circulation (poumon) :} Cœur droit $\to$ Artère pulmonaire $\to$ Poumons (hématose) $\to$ Veines pulmonaires $\to$ Cœur gauche. Sang désoxygéné $\to$ oxygéné.
      \item \textbf{Grande circulation (systémique) :} Cœur gauche $\to$ Aorte $\to$ Organes (échanges) $\to$ Veines caves $\to$ Cœur droit. Sang oxygéné $\to$ désoxygéné.
      \end{enumerate}
\item \textbf{Cycle cardiaque :} \textbf{Systole} (contraction, éjection) / \textbf{Diastole} (relâchement, remplissage). Fréquence au repos $\approx 60-80\,\text{bpm}$.
\item \textbf{Coronaires :} Artères vascularisant le myocarde. Obstruction = infarctus.
\end{itemize}

\courssub{Les hémorragies et les moyens de lutte}

\begin{definition}[Hémorragie]
Écoulement de sang hors du réseau vasculaire. Urgence vitale si volume $> 500\,mL$ (adulte) ou signes de choc.
\end{definition}

\begin{center}
\begin{tabularx}{\linewidth}{|l|X|X|}\hline
\rowcolor{popblueL}
\textbf{Type} & \textbf{Aspect du sang} & \textbf{Gravité / Conduite} \\ \hline
\textbf{Capillaire} & Suintement lent, rouge clair & Bénigne, arrêt spontané ou pansement \\ \hline
\textbf{Veineuse} & Flot continu, rouge foncé & Modérée à grave, \textbf{compression} directe \\ \hline
\textbf{Artérielle} & Jet pulsatile, rouge vif & \textbf{Extrême urgence}, \textbf{compression} forte en amont \\ \hline
\end{tabularx}
\end{center}

\begin{methode}[Geste qui sauve : Compression manuelle]
\begin{enumerate}
\item Mettre des gants (ou sac plastique) si possible.
\item Comprimer \textbf{directement} sur la plaie avec un linge propre (ou main nue).
\item Allonger la victime, surélever le membre si possible (sauf fracture).
\item Maintenir la compression \textbf{jusqu'à l'arrivée des secours} (ne pas retirer le pansement imbibé, en ajouter un par-dessus).
\item Alerter : \textbf{112} (ou 15 SAMU, 18 Pompiers).
\item Garrot \textbf{uniquement} si compression impossible (amputation, membre écrasé) : noter l'heure de pose.
\end{enumerate}
\end{methode}

\courssub{Maladies cardiovasculaires : causes et prévention}

\begin{definition}[Athérosclérose]
Maladie dégénérative des artères : dépôt de \cle{plaques d'athérome} (lipides, cholestérol, cellules) sur l'intima $\to$ épaississement, perte d'élasticité, rétrécissement (sténose) $\to$ ischémie.
\end{definition}

\begin{propriete}[Principaux facteurs de risque modifiables]
\begin{itemize}
\item Tabagisme (toxiques vasculaires, vasoconstriction).
\item Hypertension artérielle (HTA) : Pression $> 140/90\,mmHg$ (au cabinet).
\item Dyslipidémie : Excès LDL-cholestérol (« mauvais »), manque HDL.
\item Diabète déséquilibré (glycation protéines vasculaires).
\item Sédentarité, surpoids/obésité (IMC $> 25$), alimentation riche en sel/graisses saturées.
\item Stress chronique, consommation excessive d'alcool.
\end{itemize}
\end{propriete}

\begin{propriete}[Prévention (Hygiène de vie)]
Activité physique régulière ($30\,min/j$ marche rapide) ; Alimentation méditerranéenne (fruits, légumes, poisson, huile d'olive, peu de sel) ; Arrêt tabac ; Contrôle tension/ glycémie/ lipides ; Gestion du stress.
\end{propriete}

\courssub{Les accidents vasculaires cérébraux (AVC) et les premiers secours}

\begin{definition}[AVC (Accident Vasculaire Cérébral)]
Atteinte brutale d'une zone cérébrale par interruption de l'irrigation sanguine (AVC \textbf{ischémique} $\approx 80\%$ : thrombose/embolie) ou rupture d'une artère (AVC \textbf{hémorragique} $\approx 20\%$). \textbf{Urgence vitale absolue : « Le temps, c'est le cerveau ».}
\end{definition}

\begin{methode}[Reconnaissance : Mnémotechnique \textbf{V.I.T.E.}]
\begin{itemize}
\item \textbf{V} : \textbf{Visage} déformé (bouche tordue, impossible de sourire).
\item \textbf{I} : \textbf{Incapacité} motrice (bras qui tombe, hémiparésie).
\item \textbf{T} : \textbf{Trouble} de la parole (incompréhensible, aphasie).
\item \textbf{E} : \textbf{Extrême urgence} $\to$ Appeler \textbf{15 (SAMU)} ou \textbf{112} \textbf{immédiatement}, noter l'heure de début des signes.
\end{itemize}
\end{methode}

\begin{methode}[Conduite à tenir devant un AVC]
\begin{enumerate}
\item Appeler \textbf{15 / 112} (ne pas appeler le médecin traitant d'abord).
\item Installer la victime : \textbf{demi-assise} (tête surélevée $30^\circ$) si consciente, \textbf{PLS (Position Latérale de Sécurité)} si inconsciente mais respirant.
\item \textbf{Ne rien donner à manger/boire/avaler} (risque fausse route).
\item Rassurer, surveiller la respiration, noter l'heure exacte du début des signes (décide du traitement thrombolyse < 4h30).
\end{enumerate}
\end{methode}

\end{bilan}

\begin{aretenir}[Checklist avant le BEPC]
\begin{itemize}
\item $\square$ Je sais différencier artère, veine et capillaire (structure, pression, sens, rôle).
\item $\square$ Je nomme les 4 cavités du cœur et les 4 gros vaisseaux (entrées/sorties).
\item $\square$ Je trace le parcours du sang dans la petite et la grande circulation (sens, oxygénation).
\item $\square$ J'explique le rôle des valves cardiaques (anti-retour) et je les nomme.
\item $\square$ Je décris le cycle cardiaque (systole/diastole) et je donne la FC au repos.
\item $\square$ Je reconnais les 3 types d'hémorragies (aspect, gravité) et j'applique le geste vital : \textbf{Compression}.
\item $\square$ Je cite au moins 4 facteurs de risque cardiovasculaires modifiables.
\item $\square$ Je définis l'athérosclérose et ses conséquences (sténose, ischémie, thrombose).
\item $\square$ J'identifie les signes d'un AVC grâce à la méthode \textbf{VITE}.
\item $\square$ Je sais alerter (15/112), installer la victime (demi-assise ou PLS) et noter l'heure.
\item $\square$ Je sais expliquer pourquoi le tabac, le sel et la sédentarité abîment les vaisseaux.
\end{itemize}
\end{aretenir}

\findecours

\end{document}
